% !TeX program = xelatex
% 数学 LaTeX 排版 Skill v4.0.0；版式保持 v3 金标准. XeLaTeX 编译至少两次. 不分发字体.
% WZ-LOCKED-CLASS-BEGIN
\begin{filecontents*}[overwrite]{elegantbook-original-adapter.cls}
\NeedsTeXFormat{LaTeX2e}
\ProvidesClass{elegantbook-original-adapter}[2026/09/23 v3.0 fixed mathematical book environments]
\LoadClass[10pt,letterpaper,oneside]{book}
\RequirePackage[UTF8,scheme=plain,fontset=fandol]{ctex}
\RequirePackage{fontspec}
\setmainfont{cmunrm.otf}[BoldFont=cmunbx.otf,ItalicFont=cmunti.otf,BoldItalicFont=cmunbi.otf]
\setCJKfamilyfont{sourceheader}[ItalicFont=FandolKai-Regular.otf]{FandolKai-Regular.otf}
\RequirePackage{amsmath,amssymb,mathrsfs}
\RequirePackage{amsthm,etoolbox}
\RequirePackage{xcolor,geometry,setspace,fancyhdr,enumitem,needspace,xurl}
\definecolor{structurecolor}{HTML}{880E4F}
\definecolor{main}{HTML}{9C27B0}
\geometry{paperwidth=612bp,paperheight=792bp,left=20mm,right=20mm,top=95.04bp,bottom=20mm,headheight=12pt,headsep=25pt,footskip=30pt}
\setstretch{1.6}
\setlength{\parindent}{2em}
\setlength{\parskip}{0pt}
\setlist{nosep,leftmargin=2em}
\AtBeginDocument{\setlength{\abovedisplayskip}{4pt plus 1pt minus 1pt}\setlength{\belowdisplayskip}{4pt plus 1pt minus 1pt}\setlength{\abovedisplayshortskip}{2pt}\setlength{\belowdisplayshortskip}{3pt}}
\fancyhf{}
\fancyhead[L]{{\itshape\CJKfamily{sourceheader}\rightmark}}
\fancyhead[R]{{\itshape\CJKfamily{sourceheader}\leftmark}}
\fancyfoot[C]{\thepage}
\renewcommand{\headrulewidth}{0.4pt}
\renewcommand{\headrule}{{\color{structurecolor}\hrule height\headrulewidth width\headwidth\vskip-\headrulewidth}}
\pagestyle{fancy}
\fancypagestyle{cover}{\fancyhf{}\fancyfoot[C]{\thepage}\renewcommand{\headrulewidth}{0pt}\renewcommand{\headrule}{}}
\RequirePackage{hyperref}
\hypersetup{unicode=true,colorlinks=true,linkcolor=structurecolor,urlcolor=structurecolor,citecolor=structurecolor,linktoc=section,bookmarksnumbered=true,bookmarksopen=true,pdfborder={0 0 0},pdfstartview=FitH}
\setcounter{tocdepth}{1}
\renewcommand*\l@section{\@dottedtocline{1}{1.5em}{2.4em}}
\renewcommand{\thesection}{\thechapter.\arabic{section}}
\newcommand{\originalchapter}[1]{\par\addvspace{14pt}\Needspace{5\baselineskip}\refstepcounter{chapter}\setcounter{section}{0}\addcontentsline{toc}{chapter}{\protect\numberline{\thechapter}#1}\markboth{\thechapter\quad #1}{}\begingroup\centering\color{structurecolor}\bfseries\fontsize{14.4}{20}\selectfont\thechapter\quad #1\par\endgroup\nobreak\vspace{8pt}}
\newcommand{\originalsection}[1]{\par\addvspace{9pt}\Needspace{4\baselineskip}\refstepcounter{section}\addcontentsline{toc}{section}{\protect\hspace*{1.5em}\protect\numberline{\protect\textcolor{structurecolor}{\thesection}}#1}\markright{\thesection\quad #1}\noindent{\fontsize{12}{17}\selectfont\color{structurecolor}\thesection\quad}{\bfseries\fontsize{12}{17}\selectfont #1}\par\nobreak\vspace{4pt}}

% Semantic environments are registered once here, not re-designed in manuscripts.
\newtheoremstyle{wzstatement}{6pt}{5pt}
  {\normalfont\kaishu}{0pt}{\normalfont\bfseries\color{structurecolor}}
  {}{.7em}{\thmname{#1}\thmnumber{ #2}\thmnote{\normalfont\ (#3)}}
\newtheoremstyle{wzdefinition}{6pt}{5pt}
  {\normalfont\songti}{0pt}{\normalfont\bfseries\color{structurecolor}}
  {}{.7em}{\thmname{#1}\thmnumber{ #2}\thmnote{\normalfont\ (#3)}}
\newtheoremstyle{wzproblem}{7pt}{5pt}
  {\normalfont\songti}{2em}{\normalfont\bfseries\color{main}}
  {}{.7em}{\thmname{#1}\thmnumber{ #2}}
\newtheoremstyle{wzremark}{4pt}{4pt}
  {\normalfont\songti}{0pt}{\normalfont\bfseries}
  {}{.7em}{\thmname{#1}}
\theoremstyle{wzstatement}
\newtheorem{theorem}{定理}[chapter]
\newtheorem{lemma}[theorem]{引理}
\newtheorem{proposition}[theorem]{命题}
\newtheorem{corollary}[theorem]{推论}
\theoremstyle{wzdefinition}
\newtheorem{definition}[theorem]{定义}
\theoremstyle{wzproblem}
\newtheorem{example}{例题}[chapter]
\newtheorem{exercise}{习题}[chapter]
\theoremstyle{wzremark}
\newtheorem*{remark}{注}
\renewcommand{\proofname}{证明}
% Retain the native amsthm QED stack; only adjust the fixed typography.
\expandafter\patchcmd\csname\string\proof\endcsname{\normalfont \topsep6\p@\@plus6\p@\relax}
  {\normalfont\songti \topsep5\p@\@plus1\p@\relax}
  {}{\ClassError{elegantbook-original-adapter}{Cannot patch proof spacing}{Check the amsthm version.}}
\expandafter\patchcmd\csname\string\proof\endcsname{\itshape}{\normalfont\bfseries}
  {}{\ClassError{elegantbook-original-adapter}{Cannot patch proof heading}{Check the amsthm version.}}
\expandafter\patchcmd\csname\string\proof\endcsname{\ignorespaces}
  {\setlength{\parindent}{2em}\ignorespaces}
  {}{\ClassError{elegantbook-original-adapter}{Cannot patch proof paragraphs}{Check the amsthm version.}}
\AtBeginEnvironment{proof}{\par\Needspace{3\baselineskip}}
\renewcommand{\qedsymbol}{\ensuremath{\square}}
\newenvironment{solution}{\begin{proof}[解答]}{\end{proof}}
\newcommand{\wzcheckchapter}{\ifnum\value{chapter}<1
  \ClassError{elegantbook-original-adapter}{Numbered mathematics before chapter 1}{Move it after the first originalchapter.}\fi}

\newcommand{\wzregisterguard}[1]{\AtBeginEnvironment{#1}{\wzcheckchapter\par\Needspace{3\baselineskip}}}
\forcsvlist{\wzregisterguard}{theorem,lemma,proposition,corollary,definition,example,exercise}
% Front matter and bibliography are not numbered mathematical chapters.
\newcommand{\originalpreface}{\par\addvspace{10pt}\Needspace{5\baselineskip}
  \noindent{\bfseries\fontsize{12}{17}\selectfont 前言}\par\nobreak\vspace{4pt}}
\newcommand{\originalcontents}{\par\addvspace{14pt}
  \begin{center}{\color{structurecolor}\bfseries\fontsize{14.4}{20}\selectfont 目录}\end{center}
  \vspace{-10pt}\@starttoc{toc}}
\renewcommand{\bibname}{参考文献}
\renewenvironment{thebibliography}[1]{
  \par\addvspace{12pt}\Needspace{3\baselineskip}\phantomsection
  \addcontentsline{toc}{chapter}{\bibname}
  \markboth{\bibname}{}
  \noindent{\color{structurecolor}\bfseries\fontsize{14.4}{20}\selectfont\bibname}\par\nobreak\vspace{6pt}
  \list{\@biblabel{\@arabic\c@enumiv}}
    {\settowidth\labelwidth{\@biblabel{#1}}\leftmargin\labelwidth
     \advance\leftmargin\labelsep\usecounter{enumiv}
     \let\p@enumiv\@empty\renewcommand\theenumiv{\@arabic\c@enumiv}
     \setlength{\itemsep}{3pt}\setlength{\parsep}{0pt}}
  \sloppy\clubpenalty4000\widowpenalty4000\sfcode`\.=1000\relax
}{\def\@noitemerr{\@latex@warning{Empty bibliography}}\endlist}

\numberwithin{equation}{chapter}
\renewcommand{\theequation}{\thechapter.\arabic{equation}}
\allowdisplaybreaks[2]
\emergencystretch=2em
\clubpenalty=6000
\widowpenalty=6000
\raggedbottom
\endinput
\end{filecontents*}
% WZ-LOCKED-CLASS-END
\documentclass{elegantbook-original-adapter}
% ===== 元数据内容槽 =====
\newcommand{\BookTitle}{概率论与随机变量}
\newcommand{\BookAuthor}{Scott Sheffield 原讲义; 中文重构与推导}
\newcommand{\BookDate}{MIT 18.600, Fall 2019}
\hypersetup{pdftitle={概率论与随机变量: MIT 18.600 中文重构本},pdfauthor={Scott Sheffield; 中文整理}}
\usepackage{tikz,pgfplots,booktabs,longtable}
\pgfplotsset{compat=1.18}
\usetikzlibrary{arrows.meta,calc,positioning}
\newcommand{\R}{\mathbb R}
\newcommand{\N}{\mathbb N}
\newcommand{\Z}{\mathbb Z}
\newcommand{\dd}{\,\mathrm d}
\newcommand{\ee}{\mathrm e}
\newcommand{\ii}{\mathrm i}
\newcommand{\eps}{\varepsilon}
\newcommand{\E}{\mathbb E}
\newcommand{\Prob}{\mathbb P}
\newcommand{\ind}{\mathbf 1}
\newcommand{\Var}{\operatorname{Var}}
\newcommand{\Cov}{\operatorname{Cov}}
\newcommand{\Bin}{\operatorname{Bin}}
\newcommand{\Pois}{\operatorname{Pois}}
\newcommand{\Exp}{\operatorname{Exp}}
\newcommand{\Geom}{\operatorname{Geom}}
\newcommand{\Normal}{\mathcal N}
\newcommand{\sourcelecture}[1]{\par\noindent 对应原课程第 #1 讲.\par}
\begin{document}
\frontmatter
\begin{center}
{\color{structurecolor}\bfseries\fontsize{18}{25}\selectfont\BookTitle\par}
{\bfseries MIT 18.600 中文重构本\par}
\BookAuthor\par
\BookDate\par
\end{center}
\originalpreface
概率模型把不确定的结果写成样本空间上的事件, 随机变量再把这些结果变成可以计算的量. 本书从有限计数出发, 经条件概率、分布与期望, 进入极限定理、马尔可夫链、熵和鞅. 读者需掌握中学代数、排列组合的基本语言及一元微积分; 多元积分和所需的线性代数在使用处说明.

主要依据为 Scott Sheffield 在 Massachusetts Institute of Technology 的 MIT OpenCourseWare 课程 \href{https://ocw.mit.edu/courses/18-600-probability-and-random-variables-fall-2019/pages/lecture-notes/}{18.600 Probability and Random Variables, Fall 2019} 的37份主讲义及1份补充讲义. 第16、29讲是考试, 没有讲义. 本书按数学依赖关系重组章节, 合并重复复习页, 补充原课件略去的条件、推导和证明; 不是逐页影印或外部教科书题库的翻译. 课程中的复习任务在书末建立逐讲对应. 独立作业集与考试卷不属于本次来源范围.

原件如含第三方版权保留声明, 整份文件仅登记官方来源链接, 不随仓库或源码包复制. 本书不复用其受限图片、外书题面或文字片段, 相关数学概念用自行推导的正文与另设练习说明. 图片采用原生 TikZ. 关于金融模型的部分是概率论中的理想化数学模型, 其无套利条件和连续时间极限的适用范围在正文明确给出.

MIT OCW 默认许可为 \href{https://creativecommons.org/licenses/by-nc-sa/4.0/}{CC BY-NC-SA 4.0 International}; 本中文改编以相同许可作非商业使用及相同方式共享. 原作者为 Scott Sheffield, 原机构为 MIT; 翻译、重排、增补证明和自设例题属于本整理本的改动, 不表示 MIT 或作者认可本项目. 单独权利声明优先于默认许可. 官方条款见 \href{https://ocw.mit.edu/pages/privacy-and-terms-of-use/}{MIT OCW 使用条款}. 源清单记录逐资源页数、校验和、版权核查及持有状态. 版本日期为2026-10-08.
\clearpage
\originalcontents
\clearpage
\mainmatter
\originalchapter{计数、样本空间与概率公理}
\sourcelecture{1--5}

概率论的第一项工作不是代入公式, 而是说明一次试验究竟记录什么. 同样是抽取若干物品, 记录抽取顺序与只记录最后的集合, 得到的是不同的样本空间. 若两种记录之间每个结果恰好有同样多的对应关系, 两种算法可以给出相同的概率; 若对应数不同, 就不能把粗略记录后的结果一律看成等可能. 本章先建立计数工具, 再说明什么样的数字分配可以称为概率.

\originalsection{分阶段选择与排列}
设完成一件事需要依次作出若干选择. 若第一步有 $a_1$ 种选择, 而无论前面如何选择, 第 $j$ 步总有 $a_j$ 种选择, 则总数为 $a_1\cdots a_m$. 这是乘法原理. 它可由一棵选择树解释: 每个第一层结点有同样多的后继, 逐层数到叶子即可. 条件中的``无论前面如何选择''不能省略. 如果某一步的选择数会变化, 应先分类, 对各类求乘积, 再把互不重叠的各类相加.

\begin{example}
用符号 $0,1,2,3$ 写长度为六的序列, 要求相邻符号不同. 若另要求首位不能为零, 各有多少个?
\end{example}
\begin{solution}
第一种情形首位有四种选择. 后面的每一位只须避开紧邻的前一位, 因而总有三种选择, 答案为 $4\cdot3^5$. 第二种情形只有首位的选择数变为三, 答案为 $3^6$. 后一位究竟可用哪些符号会随前一位变化, 但可用符号的\emph{数量}始终为三; 乘法原理要求的是数量恒定.
\end{solution}

从 $n$ 个不同元素中有序选出 $k$ 个, 不允许重复, 总数为
\[
 n(n-1)\cdots(n-k+1)=\frac{n!}{(n-k)!},\qquad 0\leq k\leq n.
\]
允许重复时则为 $n^k$. 这里 $n!=1\cdot2\cdots n$, 并约定 $0!=1$. 空序列只有一个, 从空集合到自身也只有一个函数, 所以这个约定表达了真实的计数事实, 同时使边界 $k=0,n$ 无须另改公式.

\begin{definition}
集合 $\{1,\ldots,n\}$ 的排列是该集合到自身的双射 $\sigma$. 用序列 $\sigma(1),\ldots,\sigma(n)$ 可以表示它, 因而排列总数为 $n!$. 两个排列的复合定义为 $(\sigma\circ\rho)(j)=\sigma(\rho(j))$.
\end{definition}
排列不只是把元素排成一列的记号. 反复施用 $\sigma$ 可以看出它的结构. 从元素 $j$ 出发, 依次写 $j,\sigma(j),\sigma^2(j),\ldots$. 有限性保证某时重复; 双射性保证第一次重复必定回到 $j$, 否则对重复位置施用逆映射会产生更早的重复. 得到的闭合序列叫一个循环. 对尚未出现的元素重复操作, 就把集合分成互不相交的循环. 循环内部换一个起点、交换各循环书写次序, 不改变排列.

例如 $(1\ 4\ 2)(3\ 5)(6)$ 表示 $1\mapsto4\mapsto2\mapsto1$, $3\leftrightarrow5$, 而 $6$ 不变. 满足 $\sigma^2(j)=j$ 的排列叫对合. 它恰好是循环长度只有一或二的排列: 循环长度为 $r$ 时, 施用两次回到原点等价于 $r$ 整除二. 长度为一的循环称固定点. 没有固定点的排列称错排, 我们将在容斥原理之后计算其数量.

\originalsection{消除顺序与组合数}
从 $n$ 个不同元素中选出 $k$ 个组成集合, 与有序选择的区别在于同一集合的 $k!$ 种次序全部合并. 每个集合恰好对应 $k!$ 个有序序列, 因而
\[
 \binom nk=\frac{n!}{k!(n-k)!}.
\]
这种``先数带标签的对象, 再除以固定重复倍数''的方法很常用. 除法的依据必须是每个最终对象都有相同的原像数. 若把两个颜色相同的物品看成不可区分, 就应除以它们内部的排列数; 若不同对象有不同数目的对称性, 则不能直接除以一个猜出的常数.

\begin{proposition}[Pascal 递推与二项式定理]
约定 $\binom nk=0$ 当 $k<0$ 或 $k>n$. 对 $n\geq1$,
\[
 \binom nk=\binom{n-1}{k-1}+\binom{n-1}{k},\qquad
 (x+y)^n=\sum_{k=0}^n\binom nk x^ky^{n-k}.
\]
特别地, $\sum_{k=0}^n\binom nk=2^n$.
\end{proposition}
\begin{proof}
固定一个指定元素. 含有它的 $k$ 元子集, 对应于从其余 $n-1$ 个元素选 $k-1$ 个; 不含它的, 对应于选 $k$ 个. 两类互不重叠且穷尽所有选择, 得递推式.

展开 $n$ 个因子 $(x+y)$ 的乘积时, 每一项来自在每个因子中选 $x$ 或 $y$. 要得到 $x^ky^{n-k}$, 恰好要指定哪 $k$ 个因子选 $x$, 所以系数为 $\binom nk$. 令 $x=y=1$ 得最后的等式. 这也说明 $n$ 元集合共有 $2^n$ 个子集, 因为每个元素独立地有``选入''和``不选入''两种计数选择.
\end{proof}

\begin{example}
十张卡片具有两种类别, 每类各五张且每张有不同编号. 不计顺序地选四张, 要求恰有三张来自第一类. 若改成逐张抽取并记录顺序, 如何验证同一比例?
\end{example}
\begin{solution}
无序选择的总数为 $\binom{10}{4}$, 合格数为 $\binom53\binom51$. 有序总数为 $10\cdot9\cdot8\cdot7$. 合格有序数为
\[
 \binom43(5\cdot4\cdot3)\cdot5.
\]
这里先指定三张第一类卡片出现的位置, 再依次填入卡片. 前后两种总数及合格数分别相差 $4!$, 因而比例相同. 不能把有序合格数除以无序总数.
\end{solution}

\originalsection{多项式系数与整数分配}
把 $n$ 个不同元素分给 $r$ 个有名称的组, 组大小依次是非负整数 $n_1,\ldots,n_r$, 且 $n_1+\cdots+n_r=n$. 把全部元素先排列, 再按规定大小切成各组, 同一分组被各组内部排列重复 $n_1!\cdots n_r!$ 次. 所以分组数为
\[
 \binom{n}{n_1,\ldots,n_r}=\frac{n!}{n_1!\cdots n_r!}.
\]
组的名称很重要: 给甲、乙两个部门分配与把人分成两个无名称的团队是不同问题. 当无名称组的大小相同, 还会有交换这些组产生的重复.

\begin{theorem}[多项式定理]
对非负整数 $n$,
\[
 (x_1+\cdots+x_r)^n
 =\sum_{\substack{n_1,\ldots,n_r\geq0\\n_1+\cdots+n_r=n}}
 \frac{n!}{n_1!\cdots n_r!}x_1^{n_1}\cdots x_r^{n_r}.
\]
\end{theorem}
\begin{proof}
乘积的每个展开项对应一个长度为 $n$ 的符号序列, 每个位置选 $r$ 个符号之一. 若符号 $i$ 出现 $n_i$ 次, 消去位置标签之后这一项变为 $x_1^{n_1}\cdots x_r^{n_r}$. 同一单项式出现的次数, 正是把 $n$ 个位置分成各个符号所占位置组的数量, 即上面的多项式系数. 对所有可能的次数向量求和即可.
\end{proof}
令所有 $x_i=1$, 可知多项式系数之和为 $r^n$. 比如 $x^2yz^3$ 在 $(x+y+z)^6$ 中的系数为 $6!/(2!1!3!)=60$.

\begin{proposition}[隔板计数]
对整数 $n\geq0,k\geq1$, 非负整数解
\[
 a_1+\cdots+a_k=n
\]
共有 $\binom{n+k-1}{k-1}$ 个. 若各 $a_i\geq1$ 且 $n\geq k$, 则共有 $\binom{n-1}{k-1}$ 个.
\end{proposition}
\begin{proof}
写出 $n$ 个星号和 $k-1$ 条隔板. 第一条隔板之前的星号数为 $a_1$, 相邻隔板之间的数依次为 $a_2,\ldots,a_{k-1}$, 最后为 $a_k$. 隔板可以相邻, 也可位于两端, 正好允许零值. 这个对应可以反过来唯一恢复星号隔板串. 从 $n+k-1$ 个位置中选出隔板位置, 得第一式. 正整数情形令 $b_i=a_i-1$, 就变为非负整数和 $n-k$ 的问题, 因而数量为 $\binom{n-1}{k-1}$.
\end{proof}
这里数的是有序的大小向量, 又叫弱组合, 不是通常将不同次序视为同一对象的整数分拆. 也不要把它与多项式系数混淆: 隔板数的是可能的组大小; 多项式系数是在组大小已给定时分配不同元素的数量.

\begin{remark}[阶乘与积分]
微积分还给出一个与以后连续分布相连的表达式. 设 $I_n=\int_0^\infty t^n\ee^{-t}\dd t$. 当 $n=0$ 时 $I_0=1$. 对 $n\geq1$ 分部积分, 边界项 $-t^n\ee^{-t}$ 在零和无穷处都为零, 所以 $I_n=nI_{n-1}=n!$. 对实数 $z>0$, 积分 $\Gamma(z)=\int_0^\infty t^{z-1}\ee^{-t}\dd t$ 收敛, 且同一计算给 $\Gamma(z+1)=z\Gamma(z)$. 因而 $\Gamma(n+1)=n!$. 这是一种把阶乘延伸到非整数的方式, 本章的计数仍只使用整数阶乘.
\end{remark}

\originalsection{样本空间与事件}
一次试验所有可能记录组成的集合称为样本空间, 记为 $\Omega$. 其中一个元素 $\omega$ 是一个完整结果. 事件是我们能够询问``是否发生''的结果集合. 掷两枚有标号的骰子时, 可用 $\Omega=\{1,\ldots,6\}^2$; ``点数和为四''是集合 $\{(1,3),(2,2),(3,1)\}$. 若只记录点数和, 空间为 $\{2,\ldots,12\}$, 但这些和并不等可能.

事件 $A\cup B$ 表示至少一个发生, $A\cap B$ 表示两者都发生, $A^c=\Omega\setminus A$ 表示 $A$ 不发生. 有时简写 $AB=A\cap B$. $A\setminus B=A\cap B^c$. 两事件互斥是说 $A\cap B=\varnothing$, 即不可能同时发生. 对任意事件族, De Morgan 法则为
\[
 \left(\bigcup_i A_i\right)^c=\bigcap_i A_i^c,\qquad
 \left(\bigcap_i A_i\right)^c=\bigcup_i A_i^c.
\]
证明只须逐个结果判断: 不在任何 $A_i$ 中, 等价于在每个 $A_i^c$ 中; 不在全部 $A_i$ 中, 等价于至少在一个补集中. 因而这些等式既适用于有限族, 也适用于可数族.

在有限空间上, 所有子集都可以作为事件. 在连续空间上, 若要求概率满足合理的可数可加性, 通常不能对任意子集都赋予一致的长度或面积概率. 我们因此明确指定一个允许的事件族.

\begin{definition}
$\Omega$ 上的 $\sigma$ 代数是子集族 $\mathcal F$, 满足 $\Omega\in\mathcal F$, 若 $A\in\mathcal F$ 则 $A^c\in\mathcal F$, 若 $A_1,A_2,\ldots\in\mathcal F$ 则 $\bigcup_{n=1}^\infty A_n\in\mathcal F$. 概率空间是三元组 $(\Omega,\mathcal F,\Prob)$, 其中 $\Prob:\mathcal F\to[0,1]$ 满足
\[
 \Prob(\Omega)=1,\qquad
 \Prob\left(\bigcup_{n=1}^\infty A_n\right)=\sum_{n=1}^\infty\Prob(A_n)
 \quad\text{当 }A_n\text{ 两两互斥}.
\]
\end{definition}
补集与并集的封闭性, 通过 De Morgan 法则也给出对可数交集的封闭性. 本书有限、可数模型通常取 $\mathcal F$ 为全部子集. 对实数上的连续模型, 使用包含所有区间、并对上述操作封闭的 Borel 事件族, 必要时再补入零概率集的子集. 本书不证明不可测集合的存在, 也不要求读者构造它们; 这里说明事件必须属于 $\mathcal F$, 防止把有限模型的方便约定误当成对所有连续子集都成立的结论.

概率公理给出数学的一致性, 并不单独决定现实中各事件的概率. 长期频率、主观信念和市场价格可以帮助建立模型, 但各有前提. 重复试验的经验比例满足有限集合运算的一致性, 而把未来概率等同于过去频率还需稳定性假设. 若以合约价格解释概率, 需要同一计价单位及可以交易相应合约等条件. 有限次无套利比较解释有限可加性, 本身不推出可数可加性. 个人认为某个故事更可信, 也不等于它的事件集合更大.

\begin{example}
某报价系统给互斥事件 $A,B$ 的一单位支付合约分别定价 $0.2,0.3$, 却给 $A\cup B$ 的一单位支付合约定价 $0.7$. 假定可以买入和卖出这些合约, 且不计手续费. 说明这些价格为什么不一致.
\end{example}
\begin{solution}
买入 $A$ 和 $B$ 两份合约的总成本为 $0.5$, 卖出 $A\cup B$ 合约收到 $0.7$, 当场净收入为 $0.2$. 因 $A,B$ 互斥, 两份买入合约的总支付在每个结果上都恰好等于卖出合约的支付, 未来净支付为零. 因而留下确定的 $0.2$ 收入. 这解释了在允许双向交易的理想化系统中, 有限可加性如何来自相同支付必须有相同价格. 若合约交易有限制或交易成本不可忽略, 这种推理的假设就须重新检查.
\end{solution}

\originalsection{从公理导出运算法则}
\begin{proposition}
概率空间中有 $\Prob(\varnothing)=0$, $\Prob(A^c)=1-\Prob(A)$. 若 $A\subseteq B$, 则
\[
 \Prob(B)=\Prob(A)+\Prob(B\setminus A)\geq\Prob(A).
\]
此外,
\[
 \Prob(A\cup B)=\Prob(A)+\Prob(B)-\Prob(A\cap B).
\]
\end{proposition}
\begin{proof}
将 $\Omega$ 与无穷多个空集作为互斥事件代入可数可加性, 得 $1=1+\sum_{n\geq1}\Prob(\varnothing)$, 非负性迫使空集概率为零. 因而可数可加性包含有限可加性. $\Omega=A\cup A^c$ 是互斥分解, 得补集公式. 对 $A\subseteq B$, 使用 $B=A\cup(B\setminus A)$ 的互斥分解, 得单调性. 最后 $A\cup B=A\cup(B\setminus A)$, 而 $B=(B\setminus A)\cup(A\cap B)$, 两式相减即得并集公式.
\end{proof}
单调性意味着 $\Prob(A\cap B)\leq\Prob(A)$. 例如``某人从事某职业并且参加某社团''不可能比``某人从事某职业''概率更大, 无论附加细节怎样使故事听起来更贴合描述. 这是事件包含关系的后果.

\begin{proposition}[连续性与并集上界]
若 $A_n\subseteq A_{n+1}$, 则 $\Prob(\bigcup_n A_n)=\lim_n\Prob(A_n)$. 若 $A_{n+1}\subseteq A_n$, 则 $\Prob(\bigcap_n A_n)=\lim_n\Prob(A_n)$. 对任意可数事件族,
\[
 \Prob\left(\bigcup_n A_n\right)\leq\sum_n\Prob(A_n).
\]
\end{proposition}
\begin{proof}
递增情形令 $D_1=A_1$, $D_n=A_n\setminus A_{n-1}$, 则 $D_n$ 两两互斥, $A_n=\bigcup_{j\leq n}D_j$, 而 $\bigcup_nA_n=\bigcup_nD_n$. 可数可加性把左边变为级数之和, 有限部分和正是 $\Prob(A_n)$. 递减情形对补集使用递增结论并用补集公式.

一般并集令 $D_1=A_1$, $D_n=A_n\setminus\bigcup_{j<n}A_j$. 它们互斥且有相同并集, 又 $D_n\subseteq A_n$. 所以 $\Prob(\bigcup_nA_n)=\sum_n\Prob(D_n)\leq\sum_n\Prob(A_n)$.
\end{proof}
可数可加性中的``互斥''和``可数''都有作用. 非互斥事件要扣除重复部分. 不可数并集则不能照搬这条公理: 在 $[0,1]$ 的均匀模型中每个单点概率为零, 全部单点的并集却有概率一. 同样, 可数无限集合上不存在各点具有相同概率的均匀概率: 相同值若为正, 总和无限; 若为零, 总和为零.

\originalsection{等可能模型与计数概率}
在有限空间 $\Omega=\{\omega_1,\ldots,\omega_N\}$ 上, 任意非负数字 $p_i$ 满足 $\sum_i p_i=1$ 都确定概率
\[
 \Prob(A)=\sum_{\omega_i\in A}p_i.
\]
反之, 公理与有限可加性说明每个概率都由各单点的概率唯一确定. 若所有点等可能, $p_i=1/N$, 因而 $\Prob(A)=|A|/|\Omega|$. 均匀性是模型假设, 不能由``列出了所有结果''推出来.

\begin{example}
独立公平的两枚八面骰各标 $1,\ldots,8$. 求点数和为五的概率. 另求六次公平抛硬币恰有四次正面的概率.
\end{example}
\begin{solution}
骰子的有序结果有 $8^2=64$ 个, 和为五的有 $(1,4),(2,3),(3,2),(4,1)$ 四个, 所以概率为 $1/16$. 硬币序列有 $2^6$ 个, 正面位置的选择有 $\binom64$ 个, 所以概率为 $\binom64/2^6=15/64$. 这里用到的独立公平假设就是各完整序列等概率; 下一章会把独立性正式定义.
\end{solution}

\begin{example}
有 $m$ 个有标号的记录, 每个记录独立且均匀地取 $d$ 个可能日期之一. 求至少两个记录日期相同的概率. 当 $m\leq d$ 时, 给出易算的近似.
\end{example}
\begin{solution}
至少一次重复的直接分类很多, 所以先数无重复. 所有分配有 $d^m$ 个, 无重复分配有 $d(d-1)\cdots(d-m+1)$ 个, 因而
\[
 \Prob(\text{重复})=1-\prod_{j=0}^{m-1}\left(1-\frac jd\right).
\]
当 $m>d$ 时必有重复, 概率为一. 当各 $j/d$ 很小时, 对数展开 $\log(1-x)=-x+O(x^2)$ 给
\[
 \log\Prob(\text{无重复})
 =-\frac{m(m-1)}{2d}+O\left(\frac{m^3}{d^2}\right).
\]
因此在误差项小时, 无重复概率近似为 $\exp(-m(m-1)/(2d))$. 这解释了为什么发生碰撞所需的记录数量在 $\sqrt d$ 的尺度上, 而不在 $d$ 的尺度上. 用生日解释模型时, 独立性、各日期均匀和忽略闰日都只是简化假设.
\end{solution}

\begin{example}
每条记录独立且均匀地取得 $q$ 个代码之一. $q$ 条记录都不取得某个指定代码的概率是多少? 当 $q$ 增大时有何极限?
\end{example}
\begin{solution}
每个位置有 $q-1$ 个允许代码, 故合格序列有 $(q-1)^q$ 个, 总序列有 $q^q$ 个. 概率为 $(1-1/q)^q$. 因 $q\log(1-1/q)\to-1$, 其极限为 $\ee^{-1}$. 这和下一节错排极限相同, 但两种模型不同: 此处各记录允许重复代码, 错排则始终是双射.
\end{solution}

\originalsection{容斥原理与错排}
对三个事件, 加上单个事件概率, 减去两两交集, 再加回三重交集, 才能让每个结果恰好被计算一次. 一般公式保留了同样的结构.

\begin{theorem}[容斥原理]
任意有限事件族 $A_1,\ldots,A_n$ 满足
\[
 \Prob\left(\bigcup_{i=1}^nA_i\right)
 =\sum_{r=1}^n(-1)^{r+1}
 \sum_{1\leq i_1<\cdots<i_r\leq n}
 \Prob(A_{i_1}\cap\cdots\cap A_{i_r}).
\]
若各 $A_i$ 是有限集合, 将概率替换为元素个数也成立.
\end{theorem}
\begin{proof}
把样本空间按每个 $A_i$ 是否发生分成至多 $2^n$ 个互斥区域. 在一个恰好属于 $m\geq1$ 个事件的区域, 右边的单事件和计入它 $\binom m1$ 次, 二重交集和计入它 $\binom m2$ 次, 依次类推. 总系数为
\[
 \sum_{r=1}^m(-1)^{r+1}\binom mr
 =1-(1-1)^m=1.
\]
不属于任何事件的区域被计入零次. 因而右边对并集内每一区域恰好计入一次, 对外部区域不计入. 对这些互斥区域使用有限可加性便得概率公式. 同样逐个元素计数便得集合公式.
\end{proof}

\begin{example}
将 $n$ 个不同物品均匀随机地一一分配到 $n$ 个有标号的位置. 物品 $i$ 的原位置为 $i$. 求没有一个物品回到原位置的概率, 并求恰有 $k$ 个回到原位置的概率.
\end{example}
\begin{solution}
令 $A_i=\{\sigma(i)=i\}$. 指定 $r$ 个位置都固定之后, 剩下 $n-r$ 个位置可任意排列, 因而任意指定的 $r$ 重交集概率为 $(n-r)!/n!$. 容斥和中这样的交集有 $\binom nr$ 个, 其总贡献绝对值为
\[
 \binom nr\frac{(n-r)!}{n!}=\frac1{r!}.
\]
对并集取补, 得无固定点概率与错排数
\[
 \Prob(\text{无固定点})=\sum_{r=0}^n\frac{(-1)^r}{r!},\qquad
 D_n=n!\sum_{r=0}^n\frac{(-1)^r}{r!}.
\]
这里 $D_0=1$, $D_1=0$. 指定恰好哪 $k$ 个位置固定有 $\binom nk$ 种, 其余位置必须错排, 所以
\[
 \Prob(\text{恰有 }k\text{ 个固定点})=\frac{\binom nkD_{n-k}}{n!}.
\]
尤其 $k=n-1$ 不可能: 唯一剩余物品也必须回到原位置. 指数函数的 Taylor 级数给 $\sum_{r\geq0}(-1)^r/r!=\ee^{-1}$; 交错级数余项估计给
\[
 \left|\frac{D_n}{n!}-\ee^{-1}\right|\leq\frac1{(n+1)!}.
\]
故无固定点概率快速趋于 $\ee^{-1}$. 把每个固定事件当独立事件, 得到 $(1-1/n)^n$, 虽有同一极限, 却不是有限 $n$ 的正确公式.
\end{solution}

\originalsection{综合练习与解答}
\begin{exercise}
将十二名不同学生分给三个有名称的实验组, 人数分别为五、四、三. 再将十二人分成三个无名称的四人组. 两种分法各有多少种?
\end{exercise}
\begin{solution}
第一种为 $12!/(5!4!3!)$. 第二种先赋予组名, 得 $12!/(4!)^3$, 每个无名称分组有 $3!$ 种赋名方式, 所以为 $12!/((4!)^3 3!)$. 第一种人数不相同且组已有名称, 无须除以 $3!$.
\end{solution}

\begin{exercise}
四种颜色各有六个不同编号的球, 不计顺序随机取八个. 求颜色数量为 $3,3,1,1$ 的概率, 以及至多出现两种颜色的概率.
\end{exercise}
\begin{solution}
总数为 $\binom{24}{8}$. 第一个事件先选哪两种颜色各出现三次, 有 $\binom42$ 种, 然后在各色内选球, 所以概率为
\[
 \frac{\binom42\binom63^2\binom61^2}{\binom{24}{8}}.
\]
第二个事件不可能只出现一种颜色, 因为每色只有六球. 所以每个合格集合恰好对应唯一的颜色对, 合格数为 $\binom42\binom{12}{8}$, 概率为其除以 $\binom{24}{8}$. 若取球数改为不超过六, 只出现一种颜色的集合就会被多个颜色对重复计入, 必须另作修正.
\end{solution}

\begin{exercise}
求 $a+b+c=10$ 的非负整数解数量, 其中 $a\leq4,b\leq4$. 另说明在一副每个点数有四张、共有十种点数的牌中, 无序抽五张恰好形成一个三张同点数和一个两张同点数的概率.
\end{exercise}
\begin{solution}
不加限制的解数为 $\binom{12}{2}$. 违反 $a\leq4$ 的解令 $a'=a-5$, 得 $\binom72$ 个, $b$ 的违反数相同. 同时违反时令 $a'=a-5,b'=b-5$, 只有 $a'=b'=c=0$ 一个解. 容斥得 $66-21-21+1=25$.

牌的总手数为 $\binom{40}{5}$. 先指定三张组的点数有十种, 再指定两张组的不同点数有九种, 在各自四张牌内选择, 得概率
\[
 \frac{10\cdot9\binom43\binom42}{\binom{40}{5}}.
\]
三张组与两张组的角色不同, 不应再除以二. 同一原则可计算两个对子: 先无序选择两种点数, 再各选两张, 最后从其余八种点数的三十二张牌选单张, 得 $\binom{10}{2}\binom42^2\cdot32/\binom{40}{5}$.
\end{solution}

\begin{exercise}
四种颜色各有六张不同编号的卡片, 无序取四张. 求至多出现两种颜色的手数. 如果另规定一种获胜模式为``同色且编号恰为 $1,2,3,4$'', 这种模式有多少手?
\end{exercise}
\begin{solution}
先对六个颜色对分别数只使用该颜色对的四张集合, 得 $\binom42\binom{12}{4}$. 只用一种颜色的集合, 在包含这种颜色的三个颜色对中各计一次, 应从三次改为一次. 因而须减去 $2\cdot4\binom64$, 答案为
\[
 \binom42\binom{12}{4}-8\binom64.
\]
这类消重与上一道取八球的题不同: 取四张时单色手牌确实存在. 指定编号的获胜模式每色只有一手, 共有四手. 其概率分别是上述手数或四除以 $\binom{24}{4}$.
\end{solution}

\originalchapter{条件概率、Bayes 公式与独立性}
\sourcelecture{6--7}

一个试验可能已经完成, 但我们只得到部分信息. 掷骰之后有人告诉我们结果是偶数, 样本空间中的奇数结果就可以排除. 条件概率把这次信息更新写成一个新的概率模型. 它也让我们能够按时间顺序分解复杂试验, 或者把``某个原因产生某种观察''的概率反过来用于推断原因.

\originalsection{信息限制与归一化}
\begin{definition}
若 $F\in\mathcal F$ 且 $\Prob(F)>0$, 定义事件 $E$ 在给定 $F$ 下的条件概率为
\[
 \Prob(E\mid F)=\frac{\Prob(E\cap F)}{\Prob(F)}.
\]
竖线右边是已知发生的事件, 左边是仍要判断的事件.
\end{definition}
在有限均匀模型中, 这个比值是 $|E\cap F|/|F|$, 相当于把 $F$ 当成新的样本空间. 在非均匀模型中, 排除 $F$ 外部的结果之后, 内部结果的相对权重仍然保持, 统一除以 $\Prob(F)$ 使总概率重新成为一. 因而信息更新并不自动把剩余结果变成等可能.

\begin{proposition}
固定正概率事件 $F$, 映射 $E\mapsto\Prob(E\mid F)$ 是 $(\Omega,\mathcal F)$ 上的概率测度. 若把样本空间限制为 $F$, 则它也在事件族 $\{E\cap F:E\in\mathcal F\}$ 上定义一个概率空间.
\end{proposition}
\begin{proof}
单调性给 $0\leq\Prob(E\cap F)\leq\Prob(F)$, 所以条件概率在 $[0,1]$ 内. 又 $\Prob(\Omega\mid F)=\Prob(F)/\Prob(F)=1$. 对两两互斥的 $E_n$, 各 $E_n\cap F$ 也互斥, 且
\[
 \Prob\left(\bigcup_n E_n\mid F\right)
 =\frac{\Prob(\bigcup_n(E_n\cap F))}{\Prob(F)}
 =\sum_n\frac{\Prob(E_n\cap F)}{\Prob(F)}.
\]
所以可数可加性成立. 限制到 $F$ 后, 补集相对于 $F$ 取, 并集仍留在事件族内; 如果两个事件与 $F$ 的交相同, 它们给出相同条件概率, 因而定义无歧义.
\end{proof}
因此前一章的补集、并集、容斥公式都可以在条件概率模型中使用. 例如 $\Prob(E^c\mid F)=1-\Prob(E\mid F)$. 同时, $\Prob(F\mid F)=1$ 并不表示 $F=\Omega$, 只表示新模型集中在 $F$ 上.

\begin{example}
公平十面骰的点数为 $1,\ldots,10$. 已知点数大于四, 求它大于七的概率. 再求已知大于七时大于四的概率.
\end{example}
\begin{solution}
令 $E=\{8,9,10\}$, $F=\{5,6,7,8,9,10\}$. 有 $E\subseteq F$, 所以 $\Prob(E\mid F)=(3/10)/(6/10)=1/2$. 反过来 $\Prob(F\mid E)=1$. 条件概率一般不对称; 它的分母由已知的信息决定.
\end{solution}

当 $\Prob(F)=0$ 时, 上面的定义是 $0/0$, 没有定义. 不能直接宣布它为零或一. 以后对连续随机变量的条件分布, 会借助密度或条件期望建立另一种定义, 其适用条件和版本需要单独说明. 本章所有使用比值的地方都要求分母为正.

\originalsection{乘法公式与多步试验}
把条件概率定义重新排列可得
\[
 \Prob(E\cap F)=\Prob(F)\Prob(E\mid F).
\]
这条式子不要求 $E,F$ 独立. 它的意义是先达到 $F$, 再在 $F$ 已发生的模型里达到 $E$.

\begin{theorem}[链式乘法公式]
若 $E_1,\ldots,E_n$ 的每个前缀交集 $E_1\cap\cdots\cap E_j$ 在 $1\leq j<n$ 时具有正概率, 则
\[
 \Prob\left(\bigcap_{j=1}^n E_j\right)
 =\Prob(E_1)\prod_{j=2}^n
 \Prob\left(E_j\mid\bigcap_{i=1}^{j-1}E_i\right).
\]
\end{theorem}
\begin{proof}
记 $F_j=\bigcap_{i=1}^jE_i$. 定义给 $\Prob(F_j)=\Prob(F_{j-1})\Prob(E_j\mid F_{j-1})$. 从 $j=2$ 开始逐次代入, 中间的 $\Prob(F_j)$ 消去, 得所述乘积. 若某个前缀的概率为零, 完整交集也因包含于它而概率为零, 但后面的条件概率可能没有定义; 此时应在零概率前缀处停止计算, 不把未定义的项写入乘积.
\end{proof}

\begin{example}
袋中有五个红球、三个蓝球, 不放回依次取三球. 求依次出现红、蓝、红的概率, 并与恰有两红一蓝的概率比较.
\end{example}
\begin{solution}
链式公式给
\[
 \Prob(RBR)=\frac58\frac37\frac46=\frac5{28}.
\]
第二项分母为七、第三项分母为六, 因为前面的球已被取走; 第三项红球数为四, 因为第一步已取出一个红球. 另两种次序 $RRB,BRR$ 有相同概率, 所以恰有两红一蓝的概率为 $15/28$. 也可在无序空间验证为 $\binom52\binom31/\binom83=15/28$. 把每一步红球概率都写成 $5/8$ 则误用了独立性.
\end{solution}

\begin{example}
均匀排列 $n$ 个元素. 已知指定的前 $r$ 个元素都固定, 求另一个指定元素也固定的概率, 其中 $0\leq r<n$.
\end{example}
\begin{solution}
已有固定条件的排列数为 $(n-r)!$, 增加一个固定条件后为 $(n-r-1)!$. 两者之比为 $1/(n-r)$. 因而指定 $k$ 个元素同时固定的概率为
\[
 \frac1n\frac1{n-1}\cdots\frac1{n-k+1}=\frac{(n-k)!}{n!}.
\]
这重新得到错排计算中的交集概率, 并说明固定事件的条件概率会随已有固定数上升.
\end{solution}

\originalsection{全概率公式与 Bayes 公式}
当一个观察可能由几种情形产生时, 应分别数出各情形的贡献, 而不是把各条件概率不加权地平均.

\begin{theorem}[全概率公式]
设 $H_1,H_2,\ldots$ 是有限或可数个两两互斥的事件, 其并为 $\Omega$, 且各 $\Prob(H_i)>0$. 则对任意事件 $B$,
\[
 \Prob(B)=\sum_i\Prob(B\mid H_i)\Prob(H_i).
\]
若 $\Prob(B)>0$, 则
\[
 \Prob(H_j\mid B)=
 \frac{\Prob(B\mid H_j)\Prob(H_j)}
 {\sum_i\Prob(B\mid H_i)\Prob(H_i)}.
\]
\end{theorem}
\begin{proof}
$B=\bigcup_i(B\cap H_i)$ 是互斥分解. 可数可加性和两事件乘法公式给出第一式. 条件概率定义给
\[
 \Prob(H_j\mid B)=\frac{\Prob(H_j\cap B)}{\Prob(B)}
 =\frac{\Prob(B\mid H_j)\Prob(H_j)}{\Prob(B)},
\]
再用第一式展开分母便得第二式. 若划分中有零概率事件, 它们的交集贡献为零, 可以先删去, 无须对它们定义条件概率.
\end{proof}
第二式叫 Bayes 公式. $\Prob(H_j)$ 是观察之前的先验概率; $\Prob(B\mid H_j)$ 是情形 $H_j$ 下产生观察 $B$ 的似然; $\Prob(H_j\mid B)$ 是观察之后的后验概率. Bayes 公式没有额外创造一个关于世界的假设, 它从概率公理推出. 真正需要建模的是先验和似然, 以及观测事件到底是什么.

\begin{example}
一种自动质检系统用于识别瑕疵品. 产品中瑕疵率为 $1\%$. 对瑕疵品, 系统报警概率为 $0.95$; 对合格品, 系统误报概率为 $0.04$. 已知系统报警, 求产品确有瑕疵的概率.
\end{example}
\begin{solution}
令 $D$ 为瑕疵事件, $T$ 为报警. 划分 $D,D^c$ 给
\[
 \Prob(T)=0.95\cdot0.01+0.04\cdot0.99=0.0491.
\]
因而
\[
 \Prob(D\mid T)=\frac{0.0095}{0.0491}=\frac{95}{491}\approx0.1935.
\]
可以想象十万件产品: 约一千件有瑕疵, 其中九百五十件报警; 九万九千件合格, 其中三千九百六十件误报. 报警产品中只有九百五十件有瑕疵. 这解释了为什么较高的识别率也可能伴随较低的报警后瑕疵率. ``准确率''一词若未区分这两个条件概率及误报概率, 就不足以解题.
\end{solution}

\begin{example}
两条生产线供给同一种零件. 甲线占总供给的 $1/5$, 乙线占 $4/5$. 检测人员将甲线零件正确识别为甲的概率为 $0.9$, 将乙线零件误识别为甲的概率为 $0.1$. 给定检测报告为甲, 求真实来自甲线的概率.
\end{example}
\begin{solution}
真实来源必须按实际供给占比加权. 所求为
\[
 \frac{0.9/5}{0.9/5+0.1\cdot4/5}=\frac9{13}.
\]
报告的正确识别率 $0.9$ 是给定真实来源后的概率, 不能直接作为给定报告后的概率. 若零件进入检测流程的机会还依赖生产线, 应先以进入流程后的供给占比作为先验; 总产量的占比就不一定适用.
\end{solution}

从两事件 Bayes 公式还可看出, 当 $\Prob(A)>0$ 且 $\Prob(B)>0$ 时,
\[
 \frac{\Prob(A\mid B)}{\Prob(A)}=\frac{\Prob(B\mid A)}{\Prob(B)}.
\]
右边等于一表示观察没有改变 $A$ 的概率; 等于零表示观察与 $A$ 不相容, 后验为零; 它最大可到 $1/\Prob(A)$, 这时后验为一. 这些说法都依赖模型中的似然, 并不单凭观察看起来特别就成立.

\originalsection{胜算、连续更新与模型检查}
对 $0<\Prob(A)<1$, 定义 $A$ 的胜算为 $\Prob(A)/\Prob(A^c)$. 它是发生与不发生的概率之比, 不是概率本身. 胜算为 $r$ 时, 概率为 $r/(1+r)$.

\begin{proposition}[胜算更新]
若 $0<\Prob(A)<1$, $\Prob(B\mid A)>0$, $\Prob(B\mid A^c)>0$, 则
\[
 \frac{\Prob(A\mid B)}{\Prob(A^c\mid B)}
 =\frac{\Prob(A)}{\Prob(A^c)}
 \frac{\Prob(B\mid A)}{\Prob(B\mid A^c)}.
\]
\end{proposition}
\begin{proof}
分别对 $A,A^c$ 使用 Bayes 公式. 两个后验概率的分母同为 $\Prob(B)$, 相除后消去, 剩下所述乘积. 似然比的分母为正保证胜算比有意义. 若某个似然为零, 应直接用 Bayes 公式判断后验是否为零或一.
\end{proof}
例如先验胜算为 $2:3$, 观察的似然比为四, 则后验胜算为 $8:3$, 后验概率为 $8/11$. 有多项证据 $B_1,\ldots,B_m$ 时, 正确的似然比是
\[
 \frac{\Prob(B_1\cap\cdots\cap B_m\mid A)}
 {\Prob(B_1\cap\cdots\cap B_m\mid A^c)}.
\]
可以用链式公式将它写为一串\emph{有条件}的似然比. 只有证据在给定 $A$ 后相互独立, 且在给定 $A^c$ 后也相互独立时, 才能改成单项似然比的乘积.

\begin{example}
延续瑕疵品模型, 假设两次检测在给定产品质量后独立, 并具有相同报警率. 求两次都报警之后的瑕疵概率.
\end{example}
\begin{solution}
条件独立给 $\Prob(T_1T_2\mid D)=0.95^2$, $\Prob(T_1T_2\mid D^c)=0.04^2$. 因而
\[
 \Prob(D\mid T_1T_2)=
 \frac{0.01\cdot0.95^2}{0.01\cdot0.95^2+0.99\cdot0.04^2}
 =\frac{0.009025}{0.010609}\approx0.8507.
\]
假如第二次只是重复读取同一个传感器保存的同一结果, 那么 $T_2=T_1$, 两次报警事件等于一次报警事件, 后验仍为 $95/491$. 所以增加报告数量不必然增加独立证据.
\end{solution}
主观先验也可以满足公理并按 Bayes 更新, 但公式不能替输入数字提供事实依据. 还应核对证据的选择过程: 如果只挑选最支持某结论的观察来报告, ``收到这份报告''与``预先指定的观察发生''不是同一事件. 这类偏差应写进观测机制, 而不能靠事后套用公式消除.

\originalsection{主持人问题与信息产生方式}
条件概率需要指定信息怎样产生. 对下面的试验, 开门规则是模型的一部分, 并不是无关的故事细节.

\begin{example}
四个门中恰有一个藏有奖品, 位置均匀. 参加者先选门一. 知道位置的主持人从另外三个门中打开两个无奖品的门; 如果有多种可开的二门集合, 他均匀选择其中一个. 看见主持人打开门二和门三后, 坚持门一与换到门四, 获奖概率各是多少?
\end{example}
\begin{solution}
令 $J_i$ 表示奖品在门 $i$, $O$ 表示打开门二、三. 当奖品在门一时, 三个其余门都无奖品, 有三种可开的二门集合, 所以 $\Prob(O\mid J_1)=1/3$. 当奖品在门四时, 主持人只能开门二、三, 所以 $\Prob(O\mid J_4)=1$. 当奖品在门二或三时, $O$ 不可能发生. 因而
\[
 \Prob(O)=\frac14\frac13+\frac14=\frac13,
 \quad \Prob(J_1\mid O)=\frac14,
 \quad \Prob(J_4\mid O)=\frac34.
\]
换门更有利. 在主持人总会留下恰好一个其他门的规则下, 换门赢恰好对应初选错误, 也直接给 $3/4$. 主持人只公布了自己保证无奖品的门, 这并不是从所有门中无条件删除两个门.

若主持人不知道奖品位置, 只是从门二、三、四中均匀随机打开两个, 则他可能打开有奖品的门. 此时给定他碰巧打开门二、三且两门都空, $J_1$ 和 $J_4$ 对这个观察的似然都为 $1/3$, 后验各为 $1/2$. 两种试验留下相同的画面, 但信息产生机制不同, 条件概率也不同.
\end{solution}

\begin{example}
一个家庭恰有两个孩子, 区分长幼. 为说明计数假设, 把每个孩子的性别建模为独立的男或女, 各概率 $1/2$; 出生月份独立均匀为十二个月, 且与性别独立. 已知家庭中至少有一个男孩出生于一月, 求两个孩子都是男孩的概率. 若改成均匀选一个孩子, 并观察该孩子是生于一月的男孩, 答案是否相同?
\end{example}
\begin{solution}
每个孩子有二十四种等可能的性别月份组合, 有序家庭共有 $24^2=576$ 种. 令 $s$ 为``一月男孩''这个指定类别. 至少一个孩子为 $s$ 的家庭有 $576-23^2=47$ 种. 两个孩子都为男孩且至少一个为 $s$ 的家庭, 在 $12^2$ 个男孩月份对中排除 $11^2$ 个无一月对, 留下 $23$ 种. 所求为 $23/47$.

若均匀选取一个孩子后观察到其类别为 $s$, 还要把``选了长子还是幼子''纳入样本空间. 被选孩子的结果不限制另一个孩子, 另一个仍有一半概率为男孩, 所以答案为 $1/2$. 第一种机制中两个孩子都是 $s$ 的家庭只计一次; 第二种机制中它既可通过选长子报告, 又可通过选幼子报告, 权重不同. 实际家庭数据未必满足均匀月份或上述性别假设; 这里只建立一个明确可算的模型.
\end{solution}

\originalsection{独立事件与互斥事件}
\begin{definition}
事件 $A,B$ 独立, 是指
\[
 \Prob(A\cap B)=\Prob(A)\Prob(B).
\]
当 $\Prob(B)>0$ 时, 这等价于 $\Prob(A\mid B)=\Prob(A)$; 当 $\Prob(A)>0$ 时, 也等价于 $\Prob(B\mid A)=\Prob(B)$.
\end{definition}
等价关系只须代入条件概率定义并乘除正分母. 乘积定义还能处理零概率事件, 所以比条件概率说法更普遍. 独立是一种关于概率的对称关系, 不是关于物理因果的断言. 某事件用到同一个硬币结果, 仍可能与另一个事件独立; 两个看似无关的故事, 也可能在所建模型中不独立.

\begin{proposition}
若 $A,B$ 独立, 则 $A^c,B$, $A,B^c$, $A^c,B^c$ 也独立. 零概率事件及概率一事件与任意事件独立. 若 $A,B$ 互斥且两者概率都为正, 则它们不独立.
\end{proposition}
\begin{proof}
分解 $B=(A\cap B)\cup(A^c\cap B)$ 得
\[
 \Prob(A^c\cap B)=\Prob(B)-\Prob(A\cap B)
 =(1-\Prob(A))\Prob(B)=\Prob(A^c)\Prob(B).
\]
交换 $A,B$ 并再取补可得另外两式. 若 $\Prob(A)=0$, 则单调性给 $\Prob(A\cap B)=0$, 符合独立定义. 概率一事件的结论从其零概率补集得到. 最后, 互斥时交集概率为零, 而正概率的乘积大于零, 因而不能独立.
\end{proof}
互斥描述不能同时发生, 独立描述知道其中一个发生后不改变另一个的概率. 这两个概念在正概率情形下恰好不相容.

\begin{example}
公平抛两枚硬币. 令 $A$ 为第一枚正面, $B$ 为第二枚正面, $C$ 为两枚结果不同. 验证任意两事件独立, 并判断三事件是否相互独立.
\end{example}
\begin{solution}
空间是四个等概率结果 $HH,HT,TH,TT$. 三事件概率都为 $1/2$. $A\cap B=\{HH\}$, $A\cap C=\{HT\}$, $B\cap C=\{TH\}$, 概率都为 $1/4$, 等于各边缘概率乘积. 但是 $A\cap B\cap C=\varnothing$, 概率为零, 与 $(1/2)^3=1/8$ 不同. 知道第一枚正面时, 两枚不同的概率仍为一半; 知道两枚都正面时, 两枚不同的概率却为零.
\end{solution}

\originalsection{成对独立与相互独立}
\begin{definition}
有限事件族 $A_1,\ldots,A_n$ 称为相互独立, 若对每个非空指标子集 $I\subseteq\{1,\ldots,n\}$ 都有
\[
 \Prob\left(\bigcap_{i\in I}A_i\right)=\prod_{i\in I}\Prob(A_i).
\]
若只要求 $|I|=2$ 的这些等式, 则称成对独立. 无限事件族相互独立, 是指它的每个有限子族都相互独立.
\end{definition}
三个事件相互独立时, 要检查三个二重交集和一个三重交集, 单事件条件本身自动成立. 只检查全部事件的交集也不足以保证独立. 上例说明成对独立不推出相互独立, 因为联合信息可以揭示单条信息没有揭示的约束.

\begin{proposition}
相互独立事件族中, 将任意若干事件替换为其补集, 所得事件族仍相互独立. 特别地, 对每个 $\eps_i\in\{0,1\}$, 令 $A_i^{(1)}=A_i$, $A_i^{(0)}=A_i^c$, 则
\[
 \Prob\left(\bigcap_{i=1}^n A_i^{(\eps_i)}\right)
 =\prod_{i=1}^n\Prob(A_i^{(\eps_i)}).
\]
\end{proposition}
\begin{proof}
先只替换 $A_j$. 对任何含 $j$ 的指标子集 $I$, 令 $G=\bigcap_{i\in I\setminus\{j\}}A_i$. 独立条件和互斥分解给
\[
 \Prob(G\cap A_j^c)=\Prob(G)-\Prob(G\cap A_j)
 =\left(\prod_{i\in I\setminus\{j\}}\Prob(A_i)\right)(1-\Prob(A_j)).
\]
这恰为新事件的概率乘积. 不含 $j$ 的子集条件未改变, 所以替换一个后仍相互独立. 逐个替换即可得到一般结论; 空交集 $G=\Omega$ 的概率与空乘积都取一.
\end{proof}
这个命题使独立试验的``成功或失败''序列具有乘积概率. 若 $A_i$ 都有概率 $p$, 则任何指定含 $k$ 个成功的模式概率为 $p^k(1-p)^{n-k}$, 而这类模式有 $\binom nk$ 个, 因而恰好成功 $k$ 次的概率为 $\binom nkp^k(1-p)^{n-k}$. 后面会把它作为二项分布研究.

若 $I,J$ 不相交且 $\Prob(\bigcap_{j\in J}A_j)>0$, 相互独立还给
\[
 \Prob\left(\bigcap_{i\in I}A_i\,\middle|\,\bigcap_{j\in J}A_j\right)
 =\frac{\prod_{i\in I\cup J}\Prob(A_i)}{\prod_{j\in J}\Prob(A_j)}
 =\prod_{i\in I}\Prob(A_i).
\]
由此可见, 相互独立控制的不只是两两信息, 也控制任意一组事件联合发生后的信息更新.

\originalsection{随机次序中的独立与相关}
\begin{example}
将若干不同标签均匀随机排列. 设 $E_{a,b}$ 表示标签 $a$ 出现在 $b$ 之前. 当 $a,b,c,d$ 互不相同时, 验证 $E_{a,b}$ 与 $E_{c,d}$ 独立. 再计算 $\Prob(E_{a,b}\mid E_{a,c})$.
\end{example}
\begin{solution}
交换标签 $a,b$ 是排列空间的双射, 它翻转 $E_{a,b}$ 而保持 $E_{c,d}$; 交换 $c,d$ 同样翻转另一事件. 所以这两个事件的四种真假组合有相同排列数, 各概率为 $1/4$. 每个事件概率为 $1/2$, 因而独立.

对第二问, 只需考察 $a,b,c$ 的相对次序, 六种次序等概率. $E_{a,c}$ 对应 $abc,acb,bac$ 三种, 其中前两种同时满足 $E_{a,b}$, 因而条件概率为 $2/3$, 不等于 $1/2$. 共享一个标签的先后关系在这里不独立.
\end{solution}

\begin{exercise}
在均匀随机排列中, 已知 $a$ 出现在指定的 $r$ 个不同标签之前. 另取一个不在这些标签中且不同于 $a$ 的标签 $b$. 求 $a$ 也出现在 $b$ 之前的条件概率.
\end{exercise}
\begin{solution}
只看这 $r+2$ 个标签的相对次序即可. 每种相对次序等概率: 每个指定相对次序在全部排列中对应相同数目的位置选择和其余标签排列. 已知事件表示 $a$ 是 $r+1$ 个指定标签中最早的, 所以其概率为 $1/(r+1)$. 它与 $a$ 早于 $b$ 的交集表示 $a$ 是全部 $r+2$ 个标签中最早的, 概率为 $1/(r+2)$. 两者相除, 所求为 $(r+1)/(r+2)$. 增加已有的领先信息会提高 $a$ 领先新标签的条件概率.
\end{solution}

\originalsection{练习与解答}
\begin{exercise}
盒子甲有三白一黑四球, 盒子乙有一白三黑四球. 先以概率 $1/3$ 选甲、$2/3$ 选乙, 再不放回从选定盒子取两球. 求两球都白的概率, 以及两球都白后选了甲的概率.
\end{exercise}
\begin{solution}
给定甲时两白概率为 $(3/4)(2/3)=1/2$; 给定乙时两白不可能. 全概率给 $\Prob(WW)=(1/3)(1/2)+(2/3)0=1/6$. Bayes 给 $\Prob(\text{甲}\mid WW)=1$. 一次白球的观察却不能确定盒子: 其甲后验为 $((1/3)(3/4))/((1/3)(3/4)+(2/3)(1/4))=3/5$.
\end{solution}

\begin{exercise}
设 $A,B$ 独立且各有概率 $1/2$. 在已知 $A\cup B$ 发生的条件概率空间中, $A,B$ 是否仍独立?
\end{exercise}
\begin{solution}
原模型 $\Prob(A\cap B)=1/4$, $\Prob(A\cup B)=3/4$. 条件模型中 $\Prob(A\mid A\cup B)=\Prob(B\mid A\cup B)=2/3$, 而 $\Prob(A\cap B\mid A\cup B)=1/3$. 因 $1/3\ne4/9$, 不独立. 原模型独立并不保证经过任意共同筛选后仍独立; 筛选已排除了二者都不发生的结果.
\end{solution}

\begin{exercise}
构造三个事件, 它们的三重交集概率等于三个概率的乘积, 但它们并不成对独立.
\end{exercise}
\begin{solution}
公平抛两枚硬币, 令 $A=B$ 为第一枚正面事件, $C=\varnothing$. 有 $\Prob(A\cap B\cap C)=0=\Prob(A)\Prob(B)\Prob(C)$, 但 $\Prob(A\cap B)=1/2\ne1/4$. 这表明只检查全交集的乘积公式不够. 若要避免零概率事件, 可在八个等可能结果 $\{1,\ldots,8\}$ 上取
\[
 A=\{1,2,3,4\},\quad B=\{1,2,3,5\},\quad C=\{1,6,7,8\}.
\]
三事件各概率 $1/2$, 三重交集只有元素一, 概率 $1/8$, 但 $A\cap B$ 有三个元素, 概率 $3/8\ne1/4$.
\end{solution}

\originalchapter{离散随机变量、期望与方差}
\sourcelecture{8--11}

\originalsection{从实验结果到数值分布}
一次实验的结果可能是一串正反面、一个排列或若干人的选择. 我们往往只关心其中某个数值: 正面数、某张卡的位置、到场人数. 随机变量把复杂结果转换成这样的数值, 但并不改变底层实验.
\begin{definition}
随机变量是样本空间上的实值函数 $X:\Omega\to\R$, 并要求 $\{X\le t\}$ 对每个实数 $t$ 都是事件. 若存在有限或可数集 $D\subset\R$ 使 $\Prob(X\in D)=1$, 则 $X$ 是离散随机变量. 其概率质量函数和分布函数分别为
\[
p_X(x)=\Prob(X=x),\qquad F_X(t)=\Prob(X\le t)=\sum_{x\in D,\,x\le t}p_X(x).
\]
\end{definition}
有限或可数样本空间并取所有子集为事件时, 上述事件要求自动满足. 在更一般的空间上, 它保证我们写出的概率有意义. 质量函数非负, 总和为1; 反过来, 任一满足这两点的函数都可以成为某个离散变量的质量函数. 例如取 $\Omega=D$, 令 $X(x)=x$, 以 $p_X(x)$ 给点 $x$ 分配概率即可. 同一分布可以在不同样本空间实现, 所以知道分布并不等于知道所有实验细节.

\begin{example}
从 $m$ 个不同编号中等概率随机排列, 令 $X$ 为编号1的位置. 再独立掷三个均匀的四面骰, 令 $Y$ 为点数和. 求 $X$ 的分布和 $\Prob(Y=5)$.
\end{example}
\begin{solution}
固定位置 $k$ 后, 其余编号有 $(m-1)!$ 种排列, 因此 $p_X(k)=1/m$ ($1\le k\le m$). 对第二问, 写点数为 $1+u_i$. 需要 $u_1+u_2+u_3=2$, $u_i\ge0$. 此时 $u_i\le3$ 自动成立, 非负整数解有 $\binom{4}{2}=6$ 个, 总结果数为 $4^3$, 所以概率为 $3/32$.
\end{solution}

若 $D$ 为整数集, 分布函数在相邻整数之间保持不变, 在 $k$ 处增加 $p_X(k)$. 一般情况下仍有
\[
\Prob(a<X\le b)=F_X(b)-F_X(a).
\]
区间端点不能随意互换: $\Prob(X=a)$ 可能为正. 分布函数单调不减, 右连续, 且在正负无穷处分别趋于1与0. 右连续性来自事件 $\{X\le t+1/n\}$ 递减到 $\{X\le t\}$ 和概率的连续性.

\begin{definition}
事件 $A$ 的指示变量为 $\ind_A(\omega)=1$ ($\omega\in A$), 否则为0. 因而 $\sum_{i=1}^n\ind_{A_i}$ 数出发生的事件数.
\end{definition}
指示变量不是概率: 它随实验结果取0或1, 而概率是一个固定数. 两者的联系会在期望中出现. 对排列, 定义 $A_i$ 为编号 $i$ 恰在位置 $i$, 则固定点数为 $\sum_i\ind_{A_i}$; 即使这些事件不独立, 这个逐结果恒等式仍成立.

\originalsection{多事件独立与递推}
\begin{definition}
事件 $A_1,\ldots,A_n$ 相互独立, 是指每个非空指标子集 $J$ 都满足
\[
\Prob\left(\bigcap_{j\in J}A_j\right)=\prod_{j\in J}\Prob(A_j).
\]
离散随机变量相互独立, 是指联合质量函数等于边缘质量函数的乘积.
\end{definition}
只有每一对事件满足乘法规则, 称为两两独立. 三个以上事件时, 它弱于相互独立. 例如取两个独立均匀的0、1变量 $U,V$, 事件 $U=1$, $V=1$, $U\ne V$ 都有概率 $1/2$, 每一对交集有概率 $1/4$, 但三者交集为空. 独立事件的补事件也保留相互独立性: 将一个 $A$ 换成 $A^c$ 时, 利用 $\Prob(B\cap A^c)=\Prob(B)-\Prob(B\cap A)$ 得到所需乘积, 再逐个替换. 因而独立试验中, 已知若干试验的具体结果不会改变其余试验的分布, 只要所条件化的事件概率非零.

许多概率无需枚举全部样本点, 而可按第一步分解. 第一轮后问题仍具有同样的形式时, 这种分解给出递推关系.
\begin{example}
两个参与者每轮公平地转移1单位资金, 总资金为整数 $L\ge2$. 第一人初有 $i$ 单位, $0\le i\le L$. 当一人破产时停止. 求第一人先达到 $L$ 的概率及比赛轮数的期望.
\end{example}
\begin{solution}
先确认停止不是额外假设. 从任何未停止状态出发, 连续 $L$ 轮向上就足以到达边界, 此事件概率至少 $2^{-L}$. 每一块 $L$ 轮独立于之前的试验, 因此停止时间 $\tau$ 满足
\[
\Prob(\tau>kL)\le(1-2^{-L})^k.
\]
尾概率趋于0且可求和, 所以几乎必然停止, 并且 $\E\tau<\infty$.

令 $h_i$ 为胜率. 第一轮条件分解给出 $h_i=(h_{i-1}+h_{i+1})/2$, 边界 $h_0=0,h_L=1$. 所以相邻差 $h_{i+1}-h_i$ 相等, 由边界得到 $h_i=i/L$. 破产概率为 $(L-i)/L$.

令 $t_i$ 为剩余轮数的期望. 由于期望有限, 可写
\[
t_i=1+\tfrac12t_{i-1}+\tfrac12t_{i+1},\qquad t_0=t_L=0.
\]
试取二次函数以产生常数二阶差, 代入得到 $t_i=i(L-i)$. 这也是唯一解: 两个解之差满足前面的齐次递推且边界为0, 故相邻差恒定并全为0.
\end{solution}
如果上方没有资本边界, 一个初有 $i$ 单位的公平赌徒仍以概率1最终破产. 因为在 $0$ 与任意 $L>i$ 之间比赛时, 先破产概率为 $1-i/L$; 这一事件蕴含无上界游戏最终破产. 让 $L\to\infty$ 即得概率至少为1. 这不说明平均破产时间有限; 事实上它是无穷大. 有限边界停止时间不超过无上界破产时间, 前者期望 $i(L-i)$ 随 $L$ 无界.

\begin{example}
独立试验成功概率 $p\in(0,1)$, 求在出现 $m$ 次失败之前出现 $r$ 次成功的概率, $r,m\ge1$.
\end{example}
\begin{solution}
令 $H_{r,m}$ 表示该概率, 第一轮分解得到
\[
H_{r,m}=pH_{r-1,m}+(1-p)H_{r,m-1},\quad H_{0,m}=1,\quad H_{r,0}=0.
\]
也可以预先补足到 $r+m-1$ 轮. 在这些轮次中至少有 $r$ 次成功, 当且仅当第 $r$ 次成功早于第 $m$ 次失败. 因为总轮数不足以同时包含二者, 其中一方必先达到目标. 因而
\[
H_{r,m}=\sum_{k=r}^{r+m-1}\binom{r+m-1}{k}p^k(1-p)^{r+m-1-k}.
\]
公平情况下这就是 Pascal 三角形的一段尾和除以 $2^{r+m-1}$. 它也满足上面的递推, 将计数公式与逐步分析联系起来.
\end{solution}

\originalsection{期望、函数与收敛条件}
\begin{definition}
离散随机变量 $X$ 的有限期望定义为
\[
\E X=\sum_x xp_X(x),\qquad\text{条件是 }\sum_x|x|p_X(x)<\infty.
\]
此时称 $X$ 可积. 对非负变量也允许期望为 $+\infty$, 其值是有限部分和的上确界.
\end{definition}
绝对收敛条件保证枚举顺序不影响答案. 有正负项时, 不可用一个条件收敛级数定义概率论中的有限期望. 若正部分和负部分的期望都为无穷, 表达式 $+\infty-\infty$ 没有定义.

期望是按概率加权的平均值, 不必是可能出现的结果. 均匀骰子的期望为 $7/2$, 尽管没有这样的点数. 在可数样本空间中, 既可按样本点求和, 也可先把产生相同 $X$ 值的样本点合并:
\[
\sum_{\omega\in\Omega}X(\omega)\Prob(\{\omega\})
=\sum_xx\sum_{\omega:X(\omega)=x}\Prob(\{\omega\})
=\sum_xxp_X(x).
\]
对非负项, 合并由有限部分和的上确界解释; 对有符号项, 绝对收敛允许重排.
\begin{proposition}[函数的期望]
若 $\E|g(X)|<\infty$, 则 $\E g(X)=\sum_xg(x)p_X(x)$. 若 $g\ge0$, 公式在允许无穷的意义下仍成立.
\end{proposition}
\begin{proof}
设 $g(X)$ 取值 $y$. 它的质量函数为 $\sum_{x:g(x)=y}p_X(x)$. 将 $\sum_yy\Prob(g(X)=y)$ 依这一分组展开, 得到所述公式; 重排的依据如上.
\end{proof}
这避免了每次先求 $g(X)$ 的分布. 例如骰子平方的期望为 $\sum_{j=1}^6j^2/6=91/6$, 并不等于 $(7/2)^2$.

\begin{theorem}[期望的线性性]
若 $X_1,\ldots,X_n$ 可积, $a_1,\ldots,a_n,b$ 为实常数, 则
\[
\E\left(b+\sum_{i=1}^na_iX_i\right)=b+\sum_{i=1}^na_i\E X_i.
\]
不要求随机变量独立.
\end{theorem}
\begin{proof}
二变量情形取联合质量函数 $p(x,y)$, 有
\[
\sum_{x,y}|ax+by|p(x,y)\le |a|\E|X|+|b|\E|Y|<\infty.
\]
因此可以分拆和式并分别按另一个变量求和:
$\sum_{x,y}(ax+by)p(x,y)=a\sum_xxp_X(x)+b\sum_yyp_Y(y)$.
常数期望等于常数, 再有限次应用即得结论.
\end{proof}
特别地, $\E\ind_A=\Prob(A)$. 求一个计数变量的期望时, 往往无需知道整个分布.
\begin{example}
将 $n$ 个编号均匀随机排列, 固定点数为 $K$. 求 $\E K$.
\end{example}
\begin{solution}
$K=\sum_{i=1}^n\ind_{A_i}$, 每个编号保持原位的概率为 $(n-1)!/n!=1/n$. 所以 $\E K=\sum_i1/n=1$. 这不是说每次恰有一个固定点; 某些排列没有, 有些有很多.
\end{solution}

\begin{proposition}[整数变量的尾和公式]
若 $X$ 取非负整数值, 则
\[
\E X=\sum_{k=1}^{\infty}\Prob(X\ge k),
\]
两边可同时为无穷.
\end{proposition}
\begin{proof}
逐结果有 $X=\sum_{k\ge1}\ind_{\{X\ge k\}}$. 等价地, 直接展开
$\sum_{j\ge0}jp_X(j)=\sum_{j\ge0}\sum_{k=1}^jp_X(j)$.
所有项非负, 双重和等于各有限子和的上确界, 因而可交换求和顺序, 得到 $\sum_{k\ge1}\sum_{j\ge k}p_X(j)$.
\end{proof}
有限个变量的线性性不能无条件用于有符号无穷和. 非负项可按上述理由使用; 有符号项则需例如 $\sum_i\E|X_i|<\infty$ 来保证绝对可积.

\originalsection{均值之外的波动}
相同的均值可能对应很不同的结果. 恒付1与以概率 $1/M$ 付 $M$、否则付0都有均值1, 后者却随 $M$ 增大而越来越不稳定. 期望描述中心, 方差描述偏离中心的平方大小.
\begin{definition}
若 $\E X^2<\infty$, 定义
\[
\Var(X)=\E[(X-\E X)^2],\qquad
\operatorname{SD}(X)=\sqrt{\Var(X)}.
\]
\end{definition}
二阶矩有限保证一阶矩有限, 因为 $|x|\le1+x^2$. 平方消除了正负偏离相互抵消的可能, 也加重了极端值的影响. 方差的量纲是原量纲的平方, 标准差则与变量相同.
\begin{proposition}
二阶矩有限时,
\[
\Var(X)=\E X^2-(\E X)^2,\quad
\Var(aX+b)=a^2\Var(X),\quad
\operatorname{SD}(aX+b)=|a|\operatorname{SD}(X).
\]
且 $\Var(X)=0$ 当且仅当 $X=\E X$ 几乎必然成立.
\end{proposition}
\begin{proof}
设 $\mu=\E X$, 展开平方并使用线性性得
$\E(X-\mu)^2=\E X^2-2\mu\E X+\mu^2=\E X^2-\mu^2$.
又 $aX+b-\E(aX+b)=a(X-\mu)$, 得到缩放公式. 若非负变量 $(X-\mu)^2$ 的期望为0, 则每个正概率取值都只能为0, 故 $X=\mu$ 几乎必然; 反向直接成立.
\end{proof}
取绝对值不可省去: $a=-1$ 不会使标准差变成负数. 对均匀骰子, 方差为 $91/6-49/4=35/12$. 对上述稀有支付, 二阶矩为 $M$, 方差为 $M-1$. 这解释了为什么只有均值不能判断风险.

\begin{definition}
二阶矩有限的 $X,Y$ 的协方差为
\[
\Cov(X,Y)=\E[(X-\E X)(Y-\E Y)]=\E(XY)-\E X\E Y.
\]
\end{definition}
积可积, 因为 $2|XY|\le X^2+Y^2$. 协方差的符号描述中心化变量是否倾向于同向变化.
\begin{proposition}[和的方差]
\[
\Var\left(\sum_i a_iX_i\right)=\sum_i a_i^2\Var(X_i)
+2\sum_{i<j}a_ia_j\Cov(X_i,X_j).
\]
独立变量的协方差为0, 因而独立变量的方差可以相加.
\end{proposition}
\begin{proof}
将 $\sum_i a_i(X_i-\E X_i)$ 的平方展开并求期望即可. 独立时, 联合质量函数分解, 而绝对可积允许求和, 得 $\E XY=(\E X)(\E Y)$.
\end{proof}
零协方差不保证独立. 令 $X$ 均匀取 $-1,0,1$, $Y=X^2$. 则 $\E X=\E X^3=0$, 故协方差为0; 但知道 $Y=0$ 就确定 $X=0$.

\begin{example}
对 $n\ge2$ 个编号的随机排列, 求固定点数 $K$ 的方差.
\end{example}
\begin{solution}
设 $I_i=\ind_{A_i}$. 平方展开给出
\[
\E K^2=\sum_i\E I_i^2+\sum_{i\ne j}\E I_iI_j.
\]
$I_i^2=I_i$, 贡献为 $n/n=1$. 同时固定两个指定编号有 $(n-2)!$ 种排列, 概率为 $1/[n(n-1)]$, $n(n-1)$ 个有序非对角项再贡献1. 所以 $\E K^2=2$, $\Var(K)=1$. 若 $n=1$, $K\equiv1$, 方差为0, 不能使用含 $n-1$ 分母的公式.
\end{solution}

\originalsection{Bernoulli试验与二项分布}
\begin{definition}
取值0、1且 $\Prob(I=1)=p$ 的变量称为 Bernoulli 变量, $0\le p\le1$. 独立的 $n$ 个同参数 Bernoulli 变量之和服从二项分布, 记作 $X\sim\Bin(n,p)$, $n$ 为非负整数.
\end{definition}
令 $q=1-p$. 恰有 $k$ 次成功需指定 $k$ 个位置, 每一种序列的概率为 $p^kq^{n-k}$, 故
\[
\Prob(X=k)=\binom nkp^kq^{n-k},\qquad0\le k\le n.
\]
其他整数概率为0. 总和为 $(p+q)^n=1$. 在 $p=0$ 或1时直接理解为退化分布, 避免依赖 $0^0$ 记号. Bernoulli 变量有 $\E I=p$, $\Var(I)=p-p^2=pq$.

\begin{theorem}
$X\sim\Bin(n,p)$ 时, $\E X=np$, $\Var(X)=npq$. 若 $n\ge1$, $Y\sim\Bin(n-1,p)$, 则对整数 $r\ge1$,
\[
\E X^r=np\E(Y+1)^{r-1}.
\]
\end{theorem}
\begin{proof}
指示变量分解和线性性给出均值 $np$, 独立性给出方差 $npq$. 为说明直接求和如何工作, 使用
$k\binom nk=n\binom{n-1}{k-1}$:
\begin{align*}
\E X^r
&=\sum_{k=1}^nk^{r-1}k\binom nkp^kq^{n-k}\\
&=np\sum_{j=0}^{n-1}(j+1)^{r-1}\binom{n-1}{j}p^jq^{n-1-j}.
\end{align*}
这正是所述矩递推. $r=1$ 时剩余和为1; $r=2$ 时 $\E X^2=np[(n-1)p+1]$, 减去 $(np)^2$ 再得方差. 边界参数可直接检查.
\end{proof}
在均值计算中, 每项多出的 $k$ 不是障碍, 而是使 Pascal 三角形的第 $n$ 行转成第 $n-1$ 行的线索. 更一般地, 下降阶乘 $(X)_r=X(X-1)\cdots(X-r+1)$ 满足
\[
\E(X)_r=(n)_rp^r\quad(1\le r\le n).
\]
因为 $(X)_r$ 数出成功试验中的有序 $r$ 元组, 一共有 $(n)_r$ 个候选, 各个同时成功概率为 $p^r$. 这也可以用于高阶矩计算.

\begin{example}
一场活动有24个座位, 发出26份确认, 每人独立以概率 $9/10$ 到场. 求座位不足的概率. 另有12人独立公平选择两项方案, 求两方人数相同的概率.
\end{example}
\begin{solution}
到场人数 $A\sim\Bin(26,9/10)$, 座位不足是 $A\ge25$, 因此概率为
\[
\binom{26}{25}(9/10)^{25}(1/10)+(9/10)^{26}.
\]
也可数缺席者 $B=26-A\sim\Bin(26,1/10)$, 条件是 $B\le1$, 得到同式. 第二问为 $\binom{12}{6}2^{-12}$. 固定赞成者或固定反对者若另外存在, 应先扣除其人数, 再确定随机部分需要的成功数.
\end{solution}
应用前必须核对固定试验次数、相同成功概率和独立性. 不放回抽样通常不满足独立性, 下一章将处理; 不同人的到场概率不同也不再是普通二项分布, 尽管期望仍为各概率之和.

\originalsection{期望效用与投资组合}
期望的用途之一是长期平均, 后面的极限定理会说明这一联系的精确条件. 另一个用途是比较行动结果. 但最大化金钱期望与最大化个人满意度并非同一问题.
\begin{example}
某人有初始财富10. 方案甲确定再得2, 方案乙以各 $1/2$ 的概率再得0或4. 以 $u(w)=\sqrt w$ 表示财富效用, 比较两方案.
\end{example}
\begin{solution}
两方案最终财富期望都为12. 甲的期望效用为 $\sqrt{12}$, 乙为 $(\sqrt{10}+\sqrt{14})/2$. 平方比较或利用平方根的严格凹性可得后者较小. 直接计算:
\[
\left(\frac{\sqrt{10}+\sqrt{14}}2\right)^2=6+\frac{\sqrt{140}}2<12,
\]
因为 $140<144$. 因而按此效用偏好确定支付. 效用函数和财富定义必须先指定; 不能仅凭均值断言任何人都应作同样选择.
\end{solution}
这里使用通常约定: 效用越大越好. 若使用损失或成本函数, 应最小化其期望. 这两种方向不能混用.

投资组合中, 收益的线性性与波动的交叉项尤其清楚. 设两个资产的单位收益为 $R_1,R_2$, 权重 $w$ 与 $1-w$, 则组合收益为 $R=wR_1+(1-w)R_2$. 在二阶矩有限时,
\[
\E R=w\mu_1+(1-w)\mu_2,
\quad\Var(R)=w^2v_1+(1-w)^2v_2+2w(1-w)c,
\]
其中 $\mu_i=\E R_i$, $v_i=\Var(R_i)$, $c=\Cov(R_1,R_2)$. 对 $0\le w\le1$ 是不卖空的模型; 若允许实数权重, 则是在相应理想化借贷与卖空约束下计算.
\begin{example}
两个资产收益均值相同, 方差都为4. 比较各投一半时协方差为0、4、$-4$ 的组合方差.
\end{example}
\begin{solution}
等权组合方差为 $\tfrac14\cdot4+\tfrac14\cdot4+\tfrac12c=2+c/2$. 三种情形分别为2、4、0. 它们可以实现: 独立同分布的中心化收益实现 $c=0$, 相同收益实现 $c=4$, 相反中心化收益实现 $c=-4$. 第三种情形两部分波动逐结果抵消. 因此资产数量本身不能说明分散程度, 相关性同样重要.
\end{solution}
若 $v_1+v_2-2c=\Var(R_1-R_2)>0$, 将方差作为 $w$ 的二次函数求导得到无约束最小点
\[
w_* =\frac{v_2-c}{v_1+v_2-2c}.
\]
若要求 $0\le w\le1$, 应将此点截到区间内. 分母为0时 $R_1-R_2$ 几乎必然为常数, 各权重的方差相同, 不可除以0. 金融合约的无套利定价有时使用另一种风险中性概率下的期望; 它不是实际发生概率下的期望, 其条件和构造将在金融模型章节给出.

\originalsection{练习与完整解答}
\begin{exercise}
随机变量 $X$ 以概率 $1/4,1/2,1/4$ 分别取 $-2,1,4$. 求 $F_X(t)$、$\E X$、$\Var(X)$ 及 $\E(2X-1)^2$.
\end{exercise}
\begin{solution}
分布函数在 $t<-2$ 时为0, 在 $-2\le t<1$ 时为 $1/4$, 在 $1\le t<4$ 时为 $3/4$, 在 $t\ge4$ 时为1. $\E X=-1/2+1/2+1=1$, $\E X^2=1+1/2+4=11/2$, 故方差为 $9/2$. 最后一项为 $4\E X^2-4\E X+1=19$.
\end{solution}
\begin{exercise}
独立的 $n$ 次试验成功概率为 $p$. 令 $A$ 是成功数, $B$ 是失败数. 求二者的协方差, 并解释为何它们一般不独立.
\end{exercise}
\begin{solution}
$B=n-A$, 所以 $\Cov(A,B)=\Cov(A,n-A)=-\Var(A)=-np(1-p)$. 当 $n\ge1$ 且 $0<p<1$, 它为负, 因而不独立. 更直接地, $A$ 的数值完全确定 $B$. 在退化参数下两者都是常量, 可以独立.
\end{solution}
\begin{exercise}
均匀随机排列中, 令 $C$ 为位置1、2是否都保持原位的指示变量 ($n\ge2$). 求 $\E C$, 并比较 $\E(\ind_{A_1}\ind_{A_2})$ 与 $\E\ind_{A_1}\E\ind_{A_2}$.
\end{exercise}
\begin{solution}
$\E C=1/[n(n-1)]$, 而边缘期望乘积为 $1/n^2$. 前者较大, 表明一个编号固定后, 另一个固定的条件概率从 $1/n$ 增为 $1/(n-1)$. 此处不可以套用独立 Bernoulli 和的方差.
\end{solution}

\originalchapter{泊松过程与更多离散分布}
\sourcelecture{12--14}

\originalsection{稀有事件与泊松极限}
二项模型在固定次数内数成功. 若试验次数越来越多, 单次成功越来越罕见, 但总成功数的期望维持在有限水平, 会出现另一种稳定分布. 来电、粒子发射或某类罕见结果都可能启发这种模型, 但数学推导中的独立性和概率稳定性需要在具体应用中检查.
\begin{definition}
设 $\lambda\ge0$. 非负整数变量 $X$ 服从参数 $\lambda$ 的泊松分布, 记作 $X\sim\Pois(\lambda)$, 若
\[
\Prob(X=k)=\ee^{-\lambda}\frac{\lambda^k}{k!},\qquad k=0,1,2,\ldots.
\]
$\lambda=0$ 时定义 $X\equiv0$.
\end{definition}
指数函数的 Taylor 级数保证 $\sum_{k\ge0}\ee^{-\lambda}\lambda^k/k!=1$, 所以确实得到概率分布. $\lambda$ 在这里是给定观察窗口内的平均次数, 与下一节每单位时间的速率要区分.

\begin{theorem}[二项分布的泊松极限]
若 $p_n\in[0,1]$, $np_n\to\lambda\in[0,\infty)$, $X_n\sim\Bin(n,p_n)$, 则对每个固定非负整数 $k$,
\[
\Prob(X_n=k)\longrightarrow\ee^{-\lambda}\frac{\lambda^k}{k!}.
\]
\end{theorem}
\begin{proof}
先设 $\lambda>0$. 此时 $p_n\to0$, 且对充分大的 $n$,
\[
\binom nkp_n^k=\frac{(np_n)^k}{k!}\prod_{j=0}^{k-1}(1-j/n)\longrightarrow\frac{\lambda^k}{k!}.
\]
又 $\log(1-u)=-u+O(u^2)$ ($u\to0$), 而 $np_n^2=(np_n)p_n\to0$, 所以
\[
(n-k)\log(1-p_n)=-(n-k)p_n+O(np_n^2)\to-\lambda.
\]
指数化得到剩余因子趋于 $\ee^{-\lambda}$. 若 $\lambda=0$, $\Prob(X_n\ge1)\le\E X_n=np_n\to0$, 故在0处概率趋于1, 其他点趋于0.
\end{proof}
此处 $k$ 固定是条件的一部分. 不能把点态极限当作所有尾概率的统一误差界. 对任意固定整数阈值 $r$, 可将 $0,\ldots,r-1$ 的有限和取极限, 得到 $\Prob(X_n\ge r)$ 的极限. 由于极限质量总和为1, 也可先控制有限区间外的质量, 再比较其内每一点, 说明分布整体收敛. 但在有限 $n$ 的具体尾部估计中仍须判断近似误差.

公式各部分具有计数意义. 恰好 $k$ 次稀有成功需要选择 $k$ 个位置, 约有 $n^k/k!$ 种; 每组位置的成功概率约为 $(\lambda/n)^k$, 没有其他成功的概率约为 $\ee^{-\lambda}$. 这里 $k!$ 是有序位置转成无序位置的重复计数因子. $\ee^{-\lambda}$ 则来自 $(1-\lambda/n)^n$ 的极限, 与连续复利中 $(1+\lambda/n)^n\to\ee^\lambda$ 是同一指数结构.

\begin{proposition}
若 $X\sim\Pois(\lambda)$, 则 $\E X=\Var(X)=\lambda$, 且对整数 $r\ge1$, $\E(X)_r=\lambda^r$.
\end{proposition}
\begin{proof}
$\lambda=0$ 直接成立. 对 $\lambda>0$, 非负项允许平移求和:
\[
\E(X)_r=\sum_{k=r}^\infty\frac{k!}{(k-r)!}\ee^{-\lambda}\frac{\lambda^k}{k!}
=\lambda^r\ee^{-\lambda}\sum_{j=0}^\infty\frac{\lambda^j}{j!}=\lambda^r.
\]
$r=1$ 得均值, $r=2$ 结合 $X^2=(X)_2+X$ 得 $\E X^2=\lambda^2+\lambda$, 因而方差为 $\lambda$.
\end{proof}
从二项均值与方差猜测这个结果很自然, 但分布收敛本身不能直接保证矩收敛, 所以上述求和证明不能省略.

\begin{example}
一项设备独立进行 $10000$ 次检查, 每次出现指定罕见提示的概率为 $0.0003$. 求恰有两次提示的精确概率与泊松近似. 若另一系统每小时提示次数按参数 $3/2$ 的泊松分布建模, 求一小时内至少一次提示的概率.
\end{example}
\begin{solution}
第一问的精确值为 $\binom{10000}{2}(0.0003)^2(0.9997)^{9998}$. 平均数为3, 泊松近似为 $\ee^{-3}3^2/2!$. 第二问利用补事件为 $1-\ee^{-3/2}$. 两问都依赖所指定模型; 单有平均数并不能推出泊松分布.
\end{solution}

\originalsection{泊松过程的公理与计数分布}
只记录一个窗口的次数会丢失事件的时刻. 将截至时间 $t$ 的次数记作 $N(t)$, 就得到一族定义在同一样本空间上的变量. 一条实验结果对应整条非减阶梯路径.
\begin{definition}
速率 $\lambda>0$ 的齐次泊松过程是非减、右连续、非负整数值过程 $N(t)$ ($t\ge0$), 有 $N(0)=0$, 有限时间内次数有限, 并满足:
\begin{enumerate}
\item 不相交半开区间 $(s,t]$ 中的增量 $N(t)-N(s)$ 相互独立.
\item 增量分布只依赖区间长度, 称为平稳增量.
\item 当 $h\downarrow0$ 时, $\Prob(N(h)=1)=\lambda h+o(h)$, $\Prob(N(h)\ge2)=o(h)$.
\end{enumerate}
当 $\lambda=0$ 时, 约定 $N(t)\equiv0$ 为零速率的退化泊松过程. 其中 $f(h)=o(h)$ 表示 $f(h)/h\to0$.
\end{definition}
路径条件使“跳跃次数”有明确含义. 第三个条件排除短时间内多次事件与单次事件同阶的成簇机制. 平稳增量只说明分布不随起点改变, 独立增量则说明不同窗口间没有概率依赖, 二者不能互相替代.

\begin{theorem}
上述公理推出
\[
N(t)-N(s)\sim\Pois(\lambda(t-s)),\qquad0\le s<t.
\]
\end{theorem}
\begin{proof}
只需证 $s=0$. 固定 $t>0$, 等分为 $n$ 个长度 $h=t/n$ 的窗口. 令 $B_j$ 表示第 $j$ 个窗口至少出现一次事件的指示变量, 则它们独立同分布, 成功概率
\[
p_n=\Prob(N(h)\ge1)=\lambda t/n+o(1/n).
\]
所以 $\sum_jB_j\sim\Bin(n,p_n)$, $np_n\to\lambda t$. 若每个窗口至多发生一次, 则 $N(t)=\sum_jB_j$. 两者不相等的概率不超过
\[
\sum_{j=1}^n\Prob(\text{第 }j\text{ 个窗口至少两次})=n\,o(t/n)\longrightarrow0.
\]
因此每个 $k$ 处的概率差也趋于0, 应用泊松极限即得 $\Prob(N(t)=k)=\ee^{-\lambda t}(\lambda t)^k/k!$. 平稳增量给出一般区间版本.
\end{proof}
特别地, 无事件概率为 $\ee^{-\lambda t}$. 也可以从 $\Prob(N(h)=0)=1-\lambda h+o(h)$ 直接用独立分块得
\[
\Prob(N(t)=0)=\bigl[1-\lambda t/n+o(1/n)\bigr]^n\to\ee^{-\lambda t}.
\]
左边与 $n$ 无关, 因而等于该极限. 取对数时余项仍趋于0, 因为 $n\,o(1/n)=o(1)$.

\begin{example}
速率为每小时2的泊松过程, 求前半小时恰一次且后一小时恰两次的概率, 并求在已知前半小时无事件的条件下, 随后半小时仍无事件的概率.
\end{example}
\begin{solution}
两个窗口不相交, 均值分别为1、2, 联合概率为
$(\ee^{-1}\cdot1)(\ee^{-2}2^2/2)=2\ee^{-3}$.
条件化不改变未来窗口分布, 所以第二问为 $\ee^{-1}$. 若只知道总共三次, 则两个窗口次数不再独立, 因为其和被固定.
\end{solution}

\originalsection{指数间隔与过程构造}
令 $S_r$ 为第 $r$ 次事件的时刻, $T_1=S_1$, $T_r=S_r-S_{r-1}$ ($r\ge2$). 无事件概率立即给出
\[
\Prob(T_1>t)=\Prob(N(t)=0)=\ee^{-\lambda t},\qquad t\ge0.
\]
这引出第一个连续变量, 即指数等待时间.
\begin{definition}
$T\sim\Exp(\lambda)$ ($\lambda>0$), 是指 $T\ge0$ 且 $\Prob(T>t)=\ee^{-\lambda t}$ ($t\ge0$). 其分布函数为 $1-\ee^{-\lambda t}$, 密度为 $\lambda\ee^{-\lambda t}$ ($t>0$).
\end{definition}
密度就是区间概率的积分表示, 由分布函数求导得到. 后续章节会系统说明密度与连续变量. 此处直接用微积分可算
\[
\E T=\int_0^\infty t\lambda\ee^{-\lambda t}\dd t=1/\lambda,
\quad \E T^2=2/\lambda^2,\quad\Var(T)=1/\lambda^2.
\]
两次分部积分即可得到第二个积分, 边界项 $t^j\ee^{-\lambda t}$ ($j=1,2$) 在无穷处为0. 指数分布还有无记忆性:
\[
\Prob(T>s+t\mid T>s)=\frac{\ee^{-\lambda(s+t)}}{\ee^{-\lambda s}}=\ee^{-\lambda t}.
\]
已经等了多久, 不改变继续等待的分布. 这是模型性质, 不是所有实际等待都具有的规律.

在一个随机事件时刻重新开始过程, 不能直接把“确定时间的独立增量”当作证明. 我们改用构造来同时说明所有间隔独立, 并验证过程公理.
\begin{theorem}[指数间隔构造]
取独立同分布的 $T_j\sim\Exp(\lambda)$, 令 $S_0=0$, $S_r=\sum_{j=1}^rT_j$, 以及
\[
N(t)=\max\{r\ge0:S_r\le t\}.
\]
则这是速率 $\lambda$ 的泊松过程. 反过来, 满足前述公理的过程具有这一构造的分布, 因而连续间隔 $T_j$ 独立且均为 $\Exp(\lambda)$.
\end{theorem}
\begin{proof}
先证 $S_r\to\infty$ 几乎必然. 每个 $T_j\ge1/\lambda$ 的概率为 $\ee^{-1}$, 独立性保证从任意指标以后全都小于 $1/\lambda$ 的概率为 $\lim_m(1-\ee^{-1})^m=0$. 对起始指标作可数并, 可知几乎必然有无穷多个这样的间隔, 于是部分和发散. 因此 $N(t)$ 有限, 路径右连续且单次跳跃.

关键是计算固定 $t$ 内恰有 $m$ 个到达点的联合概率密度. 由变量替换 $s_j=t_1+\cdots+t_j$ (Jacobian为1), 前 $m$ 个到达时刻的密度为
$\lambda^m\ee^{-\lambda s_m}$, 在 $0<s_1<\cdots<s_m$ 上. 再要求 $T_{m+1}>t-s_m$, 其概率为 $\ee^{-\lambda(t-s_m)}$. 故在 $0<s_1<\cdots<s_m<t$ 中, 恰有这些 $m$ 次且没有更多次的密度为常数 $\lambda^m\ee^{-\lambda t}$.

取划分 $0=a_0<a_1<\cdots<a_\ell=t$, 指定每段事件数为 $k_i$, 总数 $m=\sum_i k_i$. 第 $i$ 段内 $k_i$ 个有序时刻的区域体积为 $(a_i-a_{i-1})^{k_i}/k_i!$: 长方体中去掉相等坐标的零体积部分后, 被 $k_i!$ 种严格次序等分. 因此所求联合概率为
\[
\lambda^m\ee^{-\lambda t}\prod_{i=1}^\ell\frac{(a_i-a_{i-1})^{k_i}}{k_i!}
=\prod_{i=1}^\ell\ee^{-\lambda(a_i-a_{i-1})}
\frac{[\lambda(a_i-a_{i-1})]^{k_i}}{k_i!}.
\]
这既给出泊松边缘分布, 又证明独立、平稳增量. 对任意不相交窗口, 把空隙也加入划分再对空隙次数求和即可. 泊松公式在 $h\downarrow0$ 时给 $\Prob(N(h)=1)=\lambda h+O(h^2)$, $\Prob(N(h)\ge2)=O(h^2)$, 所以公理全成立.

反过来, 由前一定理与独立增量, 公理已经确定所有有限个确定时刻的联合分布, 且与构造相同. 非减右连续路径在所有非负有理时刻的取值决定整条路径: 任意 $t$ 的取值是有理时刻从右趋近所得极限. $S_r=\inf\{t:N(t)\ge r\}$ 也由这些有理时刻的值决定. 因此到达时刻和间隔的联合分布与构造一致, 即间隔独立同分布为指数分布.
\end{proof}
这一证明没有把连续时刻的概率当作点概率. 到达某个精确时刻的概率为0; 计算的对象是时刻落在某个区域的积分.

\begin{proposition}[第 $r$ 次事件的等待时间]
对整数 $r\ge1$, $S_r$ 的分布满足
\[
\Prob(S_r>t)=\ee^{-\lambda t}\sum_{k=0}^{r-1}\frac{(\lambda t)^k}{k!},\qquad t\ge0,
\]
其密度与矩为
\[
f_{S_r}(t)=\frac{\lambda^rt^{r-1}\ee^{-\lambda t}}{(r-1)!}\ (t>0),
\quad\E S_r=r/\lambda,\quad\Var(S_r)=r/\lambda^2.
\]
\end{proposition}
\begin{proof}
$S_r>t$ 等价于 $N(t)\le r-1$, 对泊松质量函数求有限和得尾概率. 对其求导取负, 相邻项消去, 留下所述密度. 均值与方差来自 $r$ 个独立指数间隔之和.
\end{proof}
这称为整数形状参数的 Gamma 或 Erlang 分布. 等待第一次与等待第 $r$ 次不同; 后者一般没有无记忆性.

\originalsection{叠加、分裂与条件计数}
\begin{proposition}[泊松变量的叠加]
若 $X\sim\Pois(\alpha)$ 与 $Y\sim\Pois(\beta)$ 独立, 则 $X+Y\sim\Pois(\alpha+\beta)$. 若 $\alpha+\beta>0$, 则
\[
X\mid(X+Y=m)\sim\Bin\left(m,\frac{\alpha}{\alpha+\beta}\right).
\]
\end{proposition}
\begin{proof}
卷积给出
\[
\Prob(X+Y=m)=\ee^{-(\alpha+\beta)}\sum_{j=0}^m
\frac{\alpha^j\beta^{m-j}}{j!(m-j)!}
=\ee^{-(\alpha+\beta)}\frac{(\alpha+\beta)^m}{m!}.
\]
用联合概率除以该值, 即得条件二项概率. 参数为0时相应变量恒为0, 结论仍成立.
\end{proof}
两个独立泊松过程的和也是泊松过程, 速率相加. 每个窗口的次数按上述公式相加, 不同窗口依赖的原过程增量相互独立, 因而保留独立增量. 不允许省掉两个过程彼此独立的条件: 若把同一过程加给自己, 只会出现大小为2的跳跃.

\begin{theorem}[独立标记与泊松分裂]
在速率 $\lambda$ 的泊松过程每次事件上, 独立地以概率 $p\in[0,1]$ 标为甲, 否则标为乙; 标记独立于事件时刻. 则甲、乙形成彼此独立的泊松过程, 速率分别为 $\lambda p$ 与 $\lambda(1-p)$.
\end{theorem}
\begin{proof}
对长度 $t$ 的窗口, 设甲、乙次数为 $A,B$. 给定总次数 $a+b$, 甲数服从二项分布, 因而
\begin{align*}
\Prob(A=a,B=b)
&=\ee^{-\lambda t}\frac{(\lambda t)^{a+b}}{(a+b)!}
\binom{a+b}{a}p^a(1-p)^b\\
&=\left[\ee^{-\lambda pt}\frac{(\lambda pt)^a}{a!}\right]
\left[\ee^{-\lambda(1-p)t}\frac{[\lambda(1-p)t]^b}{b!}\right].
\end{align*}
这给出两个独立泊松变量. 多个不相交窗口的原次数独立, 窗口所用标记也独立, 所以所有窗口的甲乙次数联合概率进一步分解. 由此得到两个过程彼此独立及各自独立平稳增量. $p=0,1$ 直接解释为一个空过程.
\end{proof}
固定总数之后甲乙次数相互牵制; 不固定总数时却独立, 这并不矛盾. 泊松总数的随机性恰好补偿了条件下的负相关.

上述时刻区域的常密度还说明, 给定 $N(t)=m$, $m$ 个到达时刻的有序联合密度为 $m!/t^m$. 这正是 $m$ 个独立均匀 $[0,t]$ 点排序后的密度. 因而长度为 $s$ 的子窗口次数, 在给定总数 $m$ 时服从 $\Bin(m,s/t)$. 分成多段时则是多项计数, 而各段条件次数一般不独立.

\begin{example}
服务请求按每小时6的泊松过程到达, 每次独立以概率 $1/3$ 属于甲类. 求半小时内甲恰一次且乙恰两次的概率; 已知这一小时总共9次, 求前20分钟恰3次的概率.
\end{example}
\begin{solution}
半小时甲、乙次数独立, 参数分别为1、2, 所求概率为 $2\ee^{-3}$. 第二问是在固定总数下计数, 为
$\binom93(1/3)^3(2/3)^6$, 而不能使用参数2的无条件泊松概率.
\end{solution}

\originalsection{几何分布与无记忆性}
\begin{definition}
独立试验每次以概率 $p\in(0,1]$ 成功, 第一次成功所在轮次 $G$ 服从几何分布 $\Geom(p)$:
\[
\Prob(G=k)=q^{k-1}p,\qquad k=1,2,\ldots,\quad q=1-p.
\]
\end{definition}
前 $k-1$ 次失败且第 $k$ 次成功的乘积就是该式. $p>0$ 保证 $q^n\to0$, 因而几乎必然成功, 总概率 $p\sum_{j\ge0}q^j=1$. $p=0$ 时永不成功, 不能定义同样的有限值分布. 另一个常见约定统计成功前失败数 $F=G-1$, 其支持是 $0,1,\ldots$; 使用时务必声明约定.

尾概率为 $\Prob(G>k)=q^k$, 故对非负整数 $m,n$,
\[
\Prob(G>m+n\mid G>m)=q^n=\Prob(G>n).
\]
只要条件事件概率正, 几何分布有离散无记忆性. $p=1$ 时部分条件事件概率为0, 这些条件表达式不定义.
\begin{proposition}
$\E G=1/p$, $\E G^2=(2-p)/p^2$, $\Var(G)=q/p^2$.
\end{proposition}
\begin{proof}
尾和公式给 $\E G=\sum_{k\ge0}q^k=1/p$. 当 $0<p<1$, 级数 $\sum_{k\ge1}k^2q^{k-1}p$ 的相邻项比趋于 $q<1$, 故二阶矩有限. 首轮成功时 $G-1=0$; 首轮失败时, 后续等待与原 $G$ 同分布, 所以
\[
\E(G-1)^2=q\E G^2.
\]
展开左边得 $\E G^2-2/p+1=q\E G^2$, 从而 $\E G^2=(2-p)/p^2$. 减去 $1/p^2$ 得方差. $p=1$ 时 $G\equiv1$, 公式直接成立.
\end{proof}
先检查矩有限再在等式中移项很重要, 否则可能只是对无穷量做了非法相减. 同样首轮分解给 $\E(G-1)=q\E G$, 也可推出均值, 但仍需先确认可积.

几何与指数等待可通过时间网格连接. 令网格长度 $h\downarrow0$, 每格成功概率 $p_h=\lambda h+o(h)$, 独立. 对 $t\ge0$,
\[
\Prob(hG_h>t)=(1-p_h)^{\lfloor t/h\rfloor}\to\ee^{-\lambda t}.
\]
因而格点等待时间趋于指数等待时间. 这不是说 $G_h$ 本身趋于有限变量; 未缩放的轮数平均约为 $1/(\lambda h)$, 会越来越大.

\originalsection{负二项分布与第多次成功}
\begin{definition}
在同样的独立试验中, 第 $r$ 次成功所在轮次 $K$, $r\ge1$ 为整数, 称为本章约定的负二项变量, 参数为 $(r,p)$:
\[
\Prob(K=k)=\binom{k-1}{r-1}p^rq^{k-r},\qquad k\ge r.
\]
\end{definition}
第 $k$ 次必须成功, 前 $k-1$ 次恰有 $r-1$ 次成功, 故得到公式. 它与“前 $k$ 次恰有 $r$ 次成功”的二项事件不同: 后者不要求最后一轮成功. 有些文献把失败数 $K-r$ 称为负二项变量; 它的均值相应是 $rq/p$, 不能与轮数均值混淆.

\begin{proposition}
$K$ 可写为 $r$ 个独立 $\Geom(p)$ 变量之和, 故 $\E K=r/p$, $\Var(K)=rq/p^2$. 上述质量函数总和为1.
\end{proposition}
\begin{proof}
设 $G_j$ 为第 $j-1$ 次成功之后到第 $j$ 次成功所需轮数. 对任意正整数 $g_1,\ldots,g_r$, 指定这些间隔等于指定一个有限序列: $g_1-1$ 次失败后成功, 接着 $g_2-1$ 次失败后成功, 依此继续. 独立试验给该事件概率
\[
\prod_{j=1}^r q^{g_j-1}p.
\]
这直接证明间隔的联合质量函数分解, 无需预先调用随机时刻的独立性定理. 每个间隔有限几乎必然, 因而 $K=\sum_jG_j$ 有限几乎必然, 也保证质量函数总和为1. 线性性与独立方差相加得矩公式.
\end{proof}
还可从等待时间和次数的等价事件得到
\[
\Prob(K\le n)=\Prob(\Bin(n,p)\ge r),\qquad n\ge r.
\]
在泊松网格极限中, $hK_h$ 趋于第 $r$ 次泊松事件的等待时间: 上式右边由二项泊松极限趋于 $\Prob(\Pois(\lambda t)\ge r)$, 恰好是 $S_r$ 的分布函数.

\begin{example}
监测器每分钟独立以概率 $1/50$ 发出一次提示. 从开始计时起观察300分钟, 求平均提示数、前120分钟无提示概率、第5次提示恰在第200分钟的概率和等待第5次提示的平均分钟数.
\end{example}
\begin{solution}
总提示数为 $\Bin(300,1/50)$, 均值为6. 无提示概率为 $(49/50)^{120}$. 第5次恰在200分钟需前199分钟恰4次且第200分钟提示, 因而概率为
\[
\binom{199}{4}(1/50)^5(49/50)^{195}.
\]
等待平均为 $5/(1/50)=250$ 分钟. 总数恰5次的泊松近似为 $\ee^{-6}6^5/5!$; 整个300分钟无提示的指数模型近似为 $\ee^{-6}$. 所有“第几分钟”均按完整分钟轮次计数, 精确连续时刻的点概率则为0.
\end{solution}

\originalsection{不放回抽样与超几何分布}
二项分布把各次结果视为独立. 不放回抽样的组成会随抽取变化, 但可以直接用组合计数处理.
\begin{definition}
总体有 $M\ge1$ 个对象, 其中 $a$ 个为指定类别, $0\le a\le M$. 等概率抽取一个大小为 $n$ 的子集, $0\le n\le M$. 指定类别个数 $H$ 服从超几何分布:
\[
\Prob(H=k)=\frac{\binom ak\binom{M-a}{n-k}}{\binom Mn},
\quad \max(0,n-M+a)\le k\le\min(a,n).
\]
\end{definition}
分子同时选择 $k$ 个指定对象与 $n-k$ 个其他对象, 分母是所有子集数. 各 $k$ 分类互斥并穷尽全部抽样子集, 因而质量总和为1. 这也是组合恒等式
$\sum_k\binom ak\binom{M-a}{n-k}=\binom Mn$ 的概率解释.

\begin{proposition}
令 $p=a/M$. 则 $\E H=np$. 当 $M\ge2$ 时,
\[
\Var(H)=np(1-p)\frac{M-n}{M-1}.
\]
$M=1$ 时合法抽样均为确定结果, 方差为0.
\end{proposition}
\begin{proof}
把子集抽样实现为随机排列的前 $n$ 个对象, 令 $I_j$ 为第 $j$ 个对象属于指定类别的指示变量. 由对称性 $\E I_j=p$, 所以均值为 $np$. 对 $i\ne j$,
\[
\E I_iI_j=\frac{a(a-1)}{M(M-1)},\qquad
\Cov(I_i,I_j)=\frac{a(a-1)}{M(M-1)}-\frac{a^2}{M^2}
=-\frac{p(1-p)}{M-1}.
\]
将 $n$ 个方差与 $n(n-1)$ 个有序协方差相加:
\[
\Var(H)=np(1-p)-\frac{n(n-1)p(1-p)}{M-1}
=np(1-p)\frac{M-n}{M-1}.
\]
\end{proof}
均值与同参数二项分布相同, 方差却有有限总体修正因子. 当 $n=M$ 时全部抽完, 类别数确定, 方差为0. 当总体很大、抽样比例很小且 $a/M\to p$, 对固定 $n$ 各有限结果序列的概率趋于独立 Bernoulli 序列概率, 因而超几何分布趋于 $\Bin(n,p)$. 这里固定 $n$ 或另行控制抽样比例是近似成立的条件.

\begin{example}
20个零件中6个带指定标签, 不放回抽取4个. 求恰有2个带标签的概率、期望和方差.
\end{example}
\begin{solution}
概率为 $\binom62\binom{14}2/\binom{20}4$. $p=3/10$, 期望为 $6/5$, 方差为
$4(3/10)(7/10)(16/19)=336/475$. 若把各次错误地视为独立, 会得到方差 $21/25$, 漏掉抽样产生的负相关.
\end{solution}

\originalsection{模型核查与练习}
泊松模型的平均数与方差相等, 但反过来仅有这两个矩相等也不能确定分布. 若事件受共同环境影响, 可能成簇或随时间变速. 少量数据中月度直方图接近泊松, 不足以保证极罕见尾部或跨月独立性. 例如以各 $1/2$ 的概率选择整月速率0或2, 再条件于该速率生成泊松次数, 总均值仍为1, 但总方差为 $\E\Lambda+\Var(\Lambda)=2$. 直接计算亦可见 $\E X^2=\E(\Lambda^2+\Lambda)=3$, 方差为2. 混合环境改变了尾部, 同一潜在环境还可造成不同窗口间相关. 因此“许多稀有事件”是建模线索, 独立性才是上述推导可用的数学条件.

\begin{exercise}
$X\sim\Pois(2)$, $Y\sim\Pois(1)$ 独立. 写出 $\Prob(X>Y)$ 的收敛级数及 $\Prob(X=Y)$ 的级数, 并说明如何控制截断误差.
\end{exercise}
\begin{solution}
按 $X$ 的值分解得
\[
\Prob(X>Y)=\ee^{-3}\sum_{k=1}^\infty\frac{2^k}{k!}\sum_{j=0}^{k-1}\frac1{j!},
\quad
\Prob(X=Y)=\ee^{-3}\sum_{k=0}^\infty\frac{2^k}{(k!)^2}.
\]
这些是非负互斥事件概率的和, 故收敛且不超过1. 第一式截在 $k=M$ 时遗漏不超过 $\Prob(X>M)$. 若 $M+2>2$, 泊松尾中后续项比至多 $2/(M+2)$, 因而
\[
\Prob(X>M)\le\ee^{-2}\frac{2^{M+1}}{(M+1)!}\frac1{1-2/(M+2)}.
\]
第二式截断误差同样可由 $\Prob(X>M)$ 控制. 例如 $M=15$ 时该界小于 $5\times10^{-10}$, 所以可安全给出小数近似 $\Prob(X>Y)\approx0.605703$. 若将 $Y$ 参数改为 $1/2$, 同一推导给 $\ee^{-5/2}\sum_{k\ge1}(2^k/k!)\sum_{j=0}^{k-1}((1/2)^j/j!)$, 约为 $0.730988$.
\end{solution}

\begin{exercise}
独立公平轮次使整数状态每次增加或减少1, 初始状态为7, 到0或18时停止. 求到18的概率和平均停止轮数.
\end{exercise}
\begin{solution}
用上一章已证明的有限边界递推, 胜率为 $7/18$, 期望轮数为 $7(18-7)=77$. 若原量每步变化为 $d>0$, 先将全部边界与初始值按 $d$ 缩放, 只有它们位于同一可达格点时才可使用该状态模型.
\end{solution}

\begin{exercise}
$N(t)$ 是速率3的泊松过程. 求第3次事件在时间1之后的概率, 求已等了1单位且仍无事件后再等超过2单位的概率. 两个问题有何区别?
\end{exercise}
\begin{solution}
第一个事件等价于 $N(1)\le2$, 概率为 $\ee^{-3}(1+3+9/2)$. 第二问为 $\Prob(T_1>3\mid T_1>1)=\ee^{-6}$, 使用第一次等待的无记忆性. 第3次到达时间是三个指数间隔的和, 不能套用第一次等待的尾概率公式.
\end{solution}

\begin{exercise}
将一个速率4的泊松过程独立分成甲、乙两类, 甲概率为 $1/4$. 已知前2单位时间共8次, 求其中甲类恰3次的概率. 不给定总数时, 求同一窗口甲3次乙5次的概率.
\end{exercise}
\begin{solution}
条件概率为 $\binom83(1/4)^3(3/4)^5$. 无条件下两类独立, 均值分别为2、6, 联合概率为 $\ee^{-8}2^3 6^5/(3!5!)$. 前者为后者除以 $\Prob(N(2)=8)=\ee^{-8}8^8/8!$, 因而两种计算一致.
\end{solution}

\originalchapter{连续随机变量与常用连续分布}
\label{chap:continuous}
\sourcelecture{17--20}

离散分布把概率放在一个个点上. 连续模型则把概率分散到区间中: 精确命中某个点的概率可以为零, 而命中一段区间的概率仍然为正. 本章从这个区别出发, 建立密度、分布函数和积分期望的联系. 随后研究均匀、正态、指数、gamma、Cauchy 和 beta 分布. 这些分布分别描述位置、许多小扰动的总和、等待时间、累计等待时间、随机方向的投影和未知成功率.

\originalsection{密度、分布函数与积分期望}

\begin{definition}[概率密度]
若存在非负可积函数 $f_X:\R\to[0,\infty)$, 满足 $\int_{\R}f_X(x)\dd x=1$, 且对每个可测集合 $B$ 有
\[
 \Prob(X\in B)=\int_B f_X(x)\dd x,
\]
则称 $X$ 具有概率密度 $f_X$, 或称其分布绝对连续. 本章中的连续随机变量均指具有密度的随机变量.
\end{definition}

密度 $f_X(x)$ 的单位是随机变量单位的倒数. 它表示单位长度所承载的概率, 并非点 $x$ 本身的概率. 例如均匀分布在长度 $1/5$ 的区间上时, 密度为 $5$, 这不违反概率不超过 $1$. 对单点积分为零, 所以 $\Prob(X=x)=0$. 在 $f_X$ 连续的点,
\[
 \Prob(x<X\le x+h)=\int_x^{x+h}f_X(t)\dd t
 =f_X(x)h+o(h)\quad(h\downarrow0).
\]
密度在有限个点上的值可以改变而不改变分布. 因而写密度公式时, 支撑集端点取开区间还是闭区间没有概率上的差别.

\begin{definition}[分布函数]
任意实值随机变量的分布函数为 $F_X(t)=\Prob(X\le t)$. 若 $X$ 有密度, 则
\[
 F_X(t)=\int_{-\infty}^t f_X(x)\dd x.
\]
在密度的连续点有 $F_X'(t)=f_X(t)$.
\end{definition}

\begin{proposition}[区间端点]
令 $F_X(a-)=\lim_{t\uparrow a}F_X(t)=\Prob(X<a)$. 对任意分布和 $a<b$,
\begin{align*}
 \Prob(a<X\le b)&=F_X(b)-F_X(a),\\
 \Prob(a\le X\le b)&=F_X(b)-F_X(a-),\\
 \Prob(X=a)&=F_X(a)-F_X(a-).
\end{align*}
有密度时 $F_X$ 连续, 所以上述区间概率不受端点影响.
\end{proposition}
\begin{proof}
把事件 $\{X\le b\}$ 分别分成 $\{X\le a\}$ 与 $\{a<X\le b\}$, 或 $\{X<a\}$ 与 $\{a\le X\le b\}$, 即得前两式. $\{X=a\}$ 是 $\{X\le a\}$ 去掉 $\{X<a\}$ 后的部分. 对有密度的分布, 单点概率为零, 因此没有跳跃.
\end{proof}

\begin{example}
设 $f_X(x)=2x$ 在 $0<x<1$ 上, 在其余位置为零. 求分布函数、$\Prob(1/4<X<3/4)$ 和 $\Prob(X=1/2)$.
\end{example}
\begin{solution}
归一性由 $\int_0^1 2x\dd x=1$ 保证. 逐段积分得到
\[
 F_X(t)=\begin{cases}0,&t\le0,\\t^2,&0<t<1,\\1,&t\ge1.\end{cases}
\]
因此区间概率为 $(3/4)^2-(1/4)^2=1/2$, 点概率为零. 虽然 $f_X(1/2)=1$, 这个数描述附近区间的概率密度, 不表示点概率为 $1$.
\end{solution}

分布函数的定义适用于离散、连续和混合分布. 不能把上述积分式无条件推广到所有随机变量. 例如以概率 $1/3$ 取零, 以概率 $2/3$ 在 $(0,1)$ 均匀取值的变量 $M$ 满足
\[
 F_M(t)=\begin{cases}0,&t<0,\\1/3+2t/3,&0\le t<1,\\1,&t\ge1.\end{cases}
\]
这里 $\Prob(0\le M\le1/2)=2/3$, 而 $F_M(1/2)-F_M(0)=1/3$ 只算出 $\Prob(0<M\le1/2)$. 在跳点以外求导得到 $2/3$, 其积分仅为 $2/3$, 会漏掉零点的质量. 更一般地, 分布函数连续也不自动保证存在密度; 本课程不展开奇异连续分布的构造.

\begin{proposition}[连续型期望公式]
若 $X$ 有密度 $f_X$, $g$ 为可测函数, 并且 $\int_{\R}|g(x)|f_X(x)\dd x<\infty$, 则
\[
 \E[g(X)]=\int_{\R}g(x)f_X(x)\dd x.
\]
对于非负 $g$, 同一公式成立, 允许结果为 $+\infty$. 特别地, $\E[X]$ 为有限实数需要 $\int |x|f_X(x)\dd x<\infty$.
\end{proposition}
\begin{proof}
对指示函数 $g=\ind_B$, 公式就是密度定义. 对在有限个不交集合上取常值的非负阶梯函数, 由期望和积分的有限可加性得到公式. 把一般非负 $g$ 从下方以阶梯函数逼近, 期望与积分均定义为这些非负逼近的极限, 因而得到等式. 对一般 $g$, 写成 $g=g^+-g^-$; 绝对可积条件使两个非负部分的积分都有限, 可以相减. 这也说明不能把 $+\infty-\infty$ 作为期望.
\end{proof}

该论证使用非负积分的单调逼近性质. 对本章分段连续的密度和计算函数, 积分可用微积分中的通常积分或反常积分计算; 不需要先建立完整的测度论. 但对于任意可测集合的概率, 积分应理解为 Lebesgue 积分, 而非总能存在的 Riemann 积分.

积分的线性给出期望线性. 若 $\E[X^2]<\infty$, 则
\[
 \Var(X)=\E[(X-\E X)^2]=\E[X^2]-(\E X)^2.
\]
同样由展开可得 $\E[aX+b]=a\E X+b$ 和 $\Var(aX+b)=a^2\Var(X)$. 这些公式与离散型完全一致, 区别仅在如何计算期望.

\originalsection{均匀分布与可测集合}

\begin{definition}[均匀分布]
对 $\alpha<\beta$, 在 $[\alpha,\beta]$ 上的均匀分布具有密度
\[
 f_X(x)=\frac1{\beta-\alpha}\ind_{(\alpha,\beta)}(x).
\]
它的含义是相同长度的子区间具有相同概率.
\end{definition}

若 $U$ 在 $[0,1]$ 上均匀分布, 对 $r>-1$,
\[
 \E[U^r]=\int_0^1 u^r\dd u=\frac1{r+1}.
\]
因此 $\E U=1/2$, $\E U^2=1/3$, $\Var(U)=1/12$. 当 $X=\alpha+(\beta-\alpha)U$ 时, 对 $\alpha<x<\beta$,
\[
 \Prob(X\le x)=\Prob\left(U\le\frac{x-\alpha}{\beta-\alpha}\right)
 =\frac{x-\alpha}{\beta-\alpha}.
\]
区间外的分布函数分别为 $0$ 和 $1$, 所以 $X$ 正是所需均匀变量. 从仿射变换立即得到
\[
 \E X=\frac{\alpha+\beta}{2},\qquad
 \Var(X)=\frac{(\beta-\alpha)^2}{12}.
\]
均值取中点来自对称性, 方差随长度平方缩放则来自距离的平方.

\begin{example}
公交车严格每 $12$ 分钟到站一次, 乘客到站时刻相对于班次的相位在一个周期上均匀. 求等待时间 $W$ 的均值、方差和至少等待 $9$ 分钟的概率.
\end{example}
\begin{solution}
一个周期内, 等待时间与相位之和为 $12$. 将均匀相位反射后仍均匀, 所以 $W$ 在 $[0,12]$ 上均匀. 于是 $\E W=6$, $\Var(W)=12$, $\Prob(W\ge9)=3/12=1/4$. 单位分别是分钟、分钟平方和无单位概率. 这里固定周期的假设很重要; 随机到站的车辆不一定产生均匀等待时间.
\end{solution}

对 $U$ 均匀于 $[0,1]$, 有理数集合的概率为零: 它是可数个单点的并, 每个单点概率为零, 可数可加性遂给出结论. 这并不意味着每个子集都能赋予与长度一致的概率. 原课程通过如下构造说明为什么概率公理需要指定可测事件族.

\begin{proposition}[不能给所有子集赋予平移不变的长度]
在选择公理下, 不存在定义于 $[0,1)$ 所有子集的概率 $\mu$, 同时满足可数可加性和模 $1$ 平移不变性.
\end{proposition}
\begin{proof}
把 $[0,1)$ 看成首尾相接的圆. 定义 $x\sim y$ 当且仅当 $x-y$ 为有理数. 这是等价关系. 用选择公理从每个等价类选一个代表, 组成集合 $A$. 对有理数 $r\in[0,1)$, 记
\[
 A_r=\{(a+r)\bmod1:a\in A\}.
\]
这些集合两两不交. 事实上, 若 $(a+r)\bmod1=(a'+r')\bmod1$, 则 $a-a'$ 为有理数, 因而代表的唯一性给出 $a=a'$, 随后 $r=r'$. 它们的并覆盖整个圆: 每个 $x$ 与其代表 $a$ 相差有理数, 取差值模 $1$ 即可.

有理数是可数无限集. 若平移不变性成立, 每个 $A_r$ 的概率均为同一个 $c=\mu(A)$. 可数可加性要求 $1=\sum_r c$. 当 $c=0$ 时右侧为 $0$, 当 $c>0$ 时右侧无穷大, 两者都矛盾.
\end{proof}

通常概率论保留选择公理与可数可加性, 把事件限制为可测集合. 区间、单点、它们的可数并以及本书计算中的积分区域均属这一范围. 上面的 $A$ 不是可测集合, 所以无需为它指定概率. 本书不构造 Lebesgue 测度, 但也不把任意子集都默认为事件.

\originalsection{正态密度的归一化与标准化}

二项分布的柱形图集中在 $np$ 附近, 宽度约为 $\sqrt{np(1-p)}$. 当 $n$ 增大, 若横轴按这个宽度缩放, 柱形逐渐呈现钟形. 在证明这个极限以前, 先独立建立钟形密度本身.

\begin{lemma}[高斯积分]
\[
 \int_{-\infty}^{\infty}\ee^{-x^2/2}\dd x=\sqrt{2\pi}.
\]
\end{lemma}
\begin{proof}
记积分为 $I$. 它有限: 当 $|x|\ge1$ 时 $x^2/2\ge|x|/2$, 所以尾部受可积指数函数控制. 非负被积函数允许迭代积分, 从而
\[
 I^2=\int_{\R^2}\ee^{-(x^2+y^2)/2}\dd x\dd y.
\]
用 $x=r\cos\theta$, $y=r\sin\theta$ 换元. 在半径为 $r$ 的圆环上, 面积元为 $r\dd r\dd\theta$: 径向厚度 $\dd r$ 与切向长度 $r\dd\theta$ 相乘. 在有限圆盘上应用多元换元, 再让半径趋于无穷, 得
\[
 I^2=\int_0^{2\pi}\int_0^\infty
       \ee^{-r^2/2}r\dd r\dd\theta
 =2\pi[-\ee^{-r^2/2}]_0^\infty=2\pi.
\]
由于 $I>0$, 取正平方根得到结论.
\end{proof}

\begin{definition}[正态分布]
标准正态变量 $Z$ 的密度为
\[
 \phi(z)=\frac1{\sqrt{2\pi}}\ee^{-z^2/2},\qquad z\in\R.
\]
对 $\mu\in\R$, $\sigma>0$, 称 $X=\mu+\sigma Z$ 服从正态分布 $\Normal(\mu,\sigma^2)$. 其密度为
\[
 f_X(x)=\frac1{\sigma\sqrt{2\pi}}
       \ee^{-(x-\mu)^2/(2\sigma^2)}.
\]
\end{definition}

密度中的 $1/\sigma$ 保证拉伸后总面积仍为 $1$. 具体地,
$F_X(x)=F_Z((x-\mu)/\sigma)$, 对 $x$ 求导便得到密度式. 当 $\sigma=0$ 时 $X$ 恒等于 $\mu$, 分布退化为一个点, 不再有上述密度.

\begin{proposition}
标准正态变量满足 $\E Z=0$, $\Var(Z)=1$. 因而 $X\sim\Normal(\mu,\sigma^2)$ 的均值为 $\mu$, 方差为 $\sigma^2$.
\end{proposition}
\begin{proof}
先核对均值存在:
\[
 \E|Z|=\frac2{\sqrt{2\pi}}\int_0^\infty
 z\ee^{-z^2/2}\dd z=\sqrt{\frac2\pi}<\infty.
\]
由于 $z\phi(z)$ 是奇函数, 对称积分给出 $\E Z=0$. 又有 $\phi'(z)=-z\phi(z)$, 分部积分得
\[
 \E Z^2=\int_{\R}z^2\phi(z)\dd z
 =[-z\phi(z)]_{-\infty}^{\infty}+\int_{\R}\phi(z)\dd z=1.
\]
边界项为零, 因为指数衰减快于一次幂. 最后应用期望、方差的仿射变换公式.
\end{proof}

记 $\Phi(t)=\int_{-\infty}^t\phi(z)\dd z$. 这个积分没有初等函数表达式, 可以通过数值积分或表格求值. 对称性给出
\[
 \Phi(-t)=1-\Phi(t),\qquad
 \Prob(|Z|\le t)=2\Phi(t)-1\quad(t\ge0).
\]
常用值为 $\Phi(-1)\approx0.1587$, $\Phi(-2)\approx0.0228$, $\Phi(-3)\approx0.00135$. 因而落在均值一个、两个、三个标准差内的概率分别约为 $0.6827$, $0.9545$, $0.9973$.

\begin{example}
某理想化测量误差 $X$ 服从 $\Normal(2,9)$. 求 $\Prob(-1\le X\le8)$ 及 $\Prob(X>8)$.
\end{example}
\begin{solution}
这里标准差是 $3$, 而非 $9$. 令 $Z=(X-2)/3$, 则
\[
 \Prob(-1\le X\le8)=\Prob(-1\le Z\le2)
 =\Phi(2)-\Phi(-1)\approx0.8186,
\]
且 $\Prob(X>8)=1-\Phi(2)\approx0.0228$.
\end{solution}

\originalsection{二项分布的正态极限}

\begin{theorem}[De Moivre--Laplace 定理]
固定 $0<p<1$, 记 $q=1-p$, 令 $S_n\sim\Bin(n,p)$. 对任意有限实数 $a<b$,
\[
 \lim_{n\to\infty}\Prob\left(a\le\frac{S_n-np}{\sqrt{npq}}\le b\right)
 =\Phi(b)-\Phi(a).
\]
\end{theorem}

本定理把离散概率在标准化坐标中变成曲线下面积. 不能只把二项质量 $\Prob(S_n=k)$ 画得像钟形就当作证明. 以下证明利用相邻概率的比值确定钟形, 再由总概率为 $1$ 确定归一常数, 从而不依赖 Stirling 公式.

\begin{proof}
写 $\sigma_n=\sqrt{npq}$, $m=\lfloor np\rfloor$, $p_{n,k}=\Prob(S_n=k)$. 由二项质量公式,
\[
 \frac{p_{n,k+1}}{p_{n,k}}=\frac{n-k}{k+1}\frac pq.
\]
当 $k=m+\ell$, $|\ell|\le K\sqrt n$ 时, 把分子分母分别写成 $nq+O(1)-\ell$ 与 $np+O(1)+\ell$, 用
$\log(1+u)=u+O(u^2)$ 展开, 得到一致估计
\[
 \log\frac{p_{n,m+\ell+1}}{p_{n,m+\ell}}
 =-\frac{\ell}{npq}+O\left(\frac1n+\frac{\ell^2}{n^2}\right).
\]
这里 $O$ 表示绝对值被某个固定常数乘以括号内的量控制; 常数可依赖 $p,K$, 不依赖 $n,\ell$. 对 $\ell=0,\ldots,j-1$ 求和, 对负 $j$ 则反向求和, 有
\[
 \log\frac{p_{n,m+j}}{p_{n,m}}
 =-\frac{j^2}{2\sigma_n^2}+o(1),\qquad |j|\le K\sigma_n,
\]
其中误差在该范围一致趋零. 原因是一次求和的误差至多为
$O(|j|/n+|j|^3/n^2)=O(n^{-1/2})$, 而用 $j^2$ 替换 $j(j-1)$ 的误差也趋零.

令 $c_n=\sigma_n p_{n,m}$. 上式使有限标准化区间内的概率和成为 Riemann 和:
\[
 \Prob\left(a\le\frac{S_n-np}{\sigma_n}\le b\right)
 =c_n\left(\int_a^b\ee^{-t^2/2}\dd t+o(1)\right).
\]
网格间距为 $1/\sigma_n$, $m-np$ 有界, 因而中心的小偏移不影响极限. 先取区间 $[-1,1]$, 左侧不超过 $1$, 可知 $c_n$ 有界. 任取收敛子列, 设其极限为 $c$. 由于 $\E S_n=np$, $\Var(S_n)=npq$, 标准化变量 $W_n=(S_n-np)/\sigma_n$ 满足 $\E[W_n^2]=1$. 对 $K>0$, 逐点有 $\ind_{\{|W_n|>K\}}\le W_n^2/K^2$, 取期望即得尾界
\[
 \Prob\left(\left|\frac{S_n-np}{\sigma_n}\right|>K\right)\le K^{-2}.
\]
沿子列令 $n\to\infty$, 得
\[
 1-K^{-2}\le c\int_{-K}^K\ee^{-t^2/2}\dd t\le1.
\]
再令 $K\to\infty$, 用高斯积分得到 $c=1/\sqrt{2\pi}$. 每个收敛子列都只能有此极限, 而 $c_n$ 有界, 故整个序列趋于这个值. 代回 Riemann 和公式即得结论.
\end{proof}

有限 $n$ 时, 定理只提供近似, 不保证某个预定误差. 一个整数质量代表宽度为 $1$ 的小柱, 因而常使用连续性修正
\[
 \Prob(r\le S_n\le s)\approx
 \Phi\left(\frac{s+1/2-np}{\sqrt{npq}}\right)
 -\Phi\left(\frac{r-1/2-np}{\sqrt{npq}}\right).
\]
这来自把第 $k$ 个柱子对应到 $[k-1/2,k+1/2]$, 是改善近似的规则, 并非以上极限定理额外给出的误差界. 当 $p$ 很靠近 $0$ 或 $1$, 使 $np$ 或 $nq$ 很小时, 正态近似可能较差.

\begin{example}
独立抛掷公平硬币 $40000$ 次. 用正态近似估计正面数严格大于 $20200$ 的概率.
\end{example}
\begin{solution}
均值为 $20000$, 标准差为 $\sqrt{40000/4}=100$. 严格大于意味着至少 $20201$ 个正面. 使用连续性修正,
\[
 \Prob(S>20200)\approx1-\Phi\left(\frac{20200.5-20000}{100}\right)
 =1-\Phi(2.005)\approx0.0225.
\]
若不作修正, 使用 $1-\Phi(2)\approx0.0228$. 两个结果都表示两个标准差以上的右尾概率.
\end{solution}

\originalsection{指数等待时间与无记忆性}

\begin{definition}[指数分布]
对 $\lambda>0$, 若 $X$ 的密度为 $f_X(x)=\lambda\ee^{-\lambda x}$ 在 $x>0$ 上, 其余为零, 则记 $X\sim\Exp(\lambda)$. 参数 $\lambda$ 称为速率.
\end{definition}

积分给出
\[
 F_X(t)=\begin{cases}0,&t<0,\\1-\ee^{-\lambda t},&t\ge0,\end{cases}
 \qquad \Prob(X>t)=\ee^{-\lambda t}\quad(t\ge0).
\]
尾概率比密度更便于研究等待问题. 若时间单位为小时, 则 $\lambda$ 的单位为每小时, $\lambda t$ 无单位.

\begin{proposition}[指数矩]
对整数 $n\ge0$, $\E[X^n]=n!/\lambda^n$. 因而 $\E X=1/\lambda$, $\Var(X)=1/\lambda^2$.
\end{proposition}
\begin{proof}
记 $M_n=\int_0^\infty x^n\lambda\ee^{-\lambda x}\dd x$. 对 $n\ge1$ 分部积分,
\[
 M_n=[-x^n\ee^{-\lambda x}]_0^\infty
       +n\int_0^\infty x^{n-1}\ee^{-\lambda x}\dd x
     =\frac n\lambda M_{n-1}.
\]
边界项为零, 且 $M_0=1$, 归纳即得矩公式. 再用 $M_2-M_1^2$ 得方差.
\end{proof}

\begin{theorem}[无记忆性]
若 $X\sim\Exp(\lambda)$, 则对 $s,t\ge0$,
\[
 \Prob(X>s+t\mid X>s)=\Prob(X>t)=\ee^{-\lambda t}.
\]
即在已经等待 $s$ 且事件尚未发生的条件下, 剩余等待时间 $X-s$ 仍服从同一指数分布.
\end{theorem}
\begin{proof}
事件 $\{X>s+t\}$ 包含于 $\{X>s\}$, 因而条件概率为
$\ee^{-\lambda(s+t)}/\ee^{-\lambda s}=\ee^{-\lambda t}$.
\end{proof}

\begin{proposition}[无记忆分布的刻画]
设非负、有限的随机变量 $X$ 满足 $\Prob(X>0)=1$, 每个有限 $t$ 的尾概率均为正, 且满足上述无记忆等式. 则存在 $\lambda>0$ 使 $X\sim\Exp(\lambda)$.
\end{proposition}
\begin{proof}
令 $S(t)=\Prob(X>t)$. 无记忆性给出 $S(s+t)=S(s)S(t)$. 记 $h(t)=-\log S(t)$, 则 $h(s+t)=h(s)+h(t)$, 且 $h$ 非减. 从整数倍和分数倍先得 $h(r)=rh(1)$ 对非负有理数 $r$ 成立. 再用有理数从两侧夹逼任意实数, 单调性给出 $h(t)=th(1)$. 因为 $X$ 有限, $S(t)\to0$, 所以 $h(1)>0$. 取 $\lambda=h(1)$ 即可.
\end{proof}

几何分布是对应的离散模型: 若 $G$ 记录第一次成功的试验编号, 则
$\Prob(G>m+n\mid G>m)=(1-p)^n$. 这要求成功率固定, 各次试验独立. 如果成功率本身未知, 观察会更新成功率的分布, 预测就可能改变. 例如先以相同概率选择一枚成功率为 $1/2$ 或 $3/4$ 的硬币, 再重复独立抛这同一枚硬币. 连续 $k$ 次成功后, 下一次成功的条件概率为
\[
 \frac{(1/2)^{k+1}+(3/4)^{k+1}}{(1/2)^k+(3/4)^k}.
\]
它随 $k$ 改变. 条件于选定硬币的成功率时试验独立, 去掉这个条件后则不独立. 这就是等待模型的参数假设为何必须说清楚.

\begin{proposition}[独立指数变量的最小值]
若 $X_i\sim\Exp(\lambda_i)$, $i=1,\ldots,n$, 相互独立, 则
$T=\min_i X_i\sim\Exp(\lambda_1+\cdots+\lambda_n)$.
\end{proposition}
\begin{proof}
$T>t$ 等价于所有 $X_i>t$. 因独立性,
\[
 \Prob(T>t)=\prod_{i=1}^n\ee^{-\lambda_i t}
 =\ee^{-t\sum_i\lambda_i}.
\]
尾函数唯一确定分布.
\end{proof}

最小值的公式需要相互独立. 当两个等待时间完全相同, 最小值就是原变量, 速率不会加倍. 独立时, 多个钟同时运转会增加任一事件发生的总速率.

\begin{example}
有 $n$ 个元件, 每个寿命独立服从 $\Exp(\lambda)$. 设 $M$ 为最后一个元件失效的时刻. 求 $F_M$, $\E M$ 与 $\Var(M)$.
\end{example}
\begin{solution}
对 $t\ge0$, 所有元件都在 $t$ 前失效的概率为
\[
 F_M(t)=\prod_{i=1}^n\Prob(X_i\le t)=(1-\ee^{-\lambda t})^n.
\]
求矩时更便于使用相邻失效间隔. 第一次失效前有 $n$ 个指数钟, 等待 $D_1\sim\Exp(n\lambda)$. 第一个钟响后, 其余 $n-1$ 个钟由无记忆性恢复为独立的原速率钟, 所以 $D_2\sim\Exp((n-1)\lambda)$, 依此类推.

还需说明这些间隔独立, 不能只凭各自的边缘分布相加方差. 设第一次响的是第 $i$ 个钟, 时刻在 $t$ 附近, 其余钟的剩余寿命分别为 $r_j>0$. 用独立密度作变换 $x_i=t$, $x_j=t+r_j$, 联合密度为
\[
 \lambda\ee^{-n\lambda t}\prod_{j\ne i}\lambda\ee^{-\lambda r_j}.
\]
它分解为第一次时刻和身份的密度, 乘以独立剩余指数钟的密度. 因而剩余钟的联合分布不依赖已经发生的时刻与身份. 重复这一步, 所有 $D_i$ 相互独立. 于是
\[
 M=D_1+\cdots+D_n,\qquad
 \E M=\frac1\lambda\sum_{j=1}^n\frac1j,\qquad
 \Var(M)=\frac1{\lambda^2}\sum_{j=1}^n\frac1{j^2}.
\]
\end{solution}

在速率 $\lambda$ 的 Poisson 过程中, 时长 $t$ 内没有事件的概率是 $\Prob(N(t)=0)=\ee^{-\lambda t}$. 因而第一次事件的等待时间为指数分布. 独立增量与平移不变的计数规律, 或反过来的独立指数间隔, 将离散的计数与连续的时间连接起来. 在下一节可以直接检验: 独立指数间隔产生的事件数确实为 Poisson 分布, 而不是仅凭小时间区间的近似来断言.

\originalsection{Gamma 函数与累计等待时间}

\begin{definition}[Gamma 函数]
对 $\alpha>0$, 定义
\[
 \Gamma(\alpha)=\int_0^\infty x^{\alpha-1}\ee^{-x}\dd x.
\]
\end{definition}

在零附近, $x^{\alpha-1}$ 可积恰好需要 $\alpha>0$; 在无穷远处, 指数衰减保证积分有限. 分部积分给出
\[
 \Gamma(\alpha+1)=\alpha\Gamma(\alpha),\qquad
 \Gamma(1)=1,\qquad \Gamma(n)=(n-1)!\quad(n\ge1).
\]
在分部积分的边界, $x^\alpha\ee^{-x}$ 在零和无穷均趋零, 因而递推式对每个正实数 $\alpha$ 成立. 取 $x=z^2/2$ 并应用高斯积分还得到 $\Gamma(1/2)=\sqrt\pi$.

\begin{definition}[Gamma 分布]
若 $\alpha,\lambda>0$, 密度为
\[
 f_X(x)=\frac{\lambda^\alpha}{\Gamma(\alpha)}
            x^{\alpha-1}\ee^{-\lambda x}\ind_{(0,\infty)}(x),
\]
称 $X$ 服从 shape 参数为 $\alpha$、rate 参数为 $\lambda$ 的 gamma 分布, 记为 $\operatorname{Gamma}(\alpha,\lambda)$. 本书固定这一参数次序; 尺度参数为 $1/\lambda$.
\end{definition}

用 $u=\lambda x$ 换元, 密度的积分为 $\Gamma(\alpha)/\Gamma(\alpha)=1$. 这也说明速率改变只是在时间轴上缩放. 对 $r> -\alpha$,
\[
 \E[X^r]=\frac{\Gamma(\alpha+r)}{\lambda^r\Gamma(\alpha)}.
\]
这是直接将 $x^r$ 乘入密度后换元得到的; 条件 $r> -\alpha$ 控制零点附近的可积性. 特别地,
\[
 \E X=\frac\alpha\lambda,\qquad
 \E X^2=\frac{\alpha(\alpha+1)}{\lambda^2},\qquad
 \Var(X)=\frac\alpha{\lambda^2}.
\]
$\alpha=1$ 正是指数分布. 对非整数 $\alpha$ 分布仍有意义, 但不能解释成等待第 $\alpha$ 次事件.

\begin{proposition}[第 $n$ 次 Poisson 事件的时间]
在速率 $\lambda$ 的 Poisson 过程中, 第 $n$ 次事件的时刻 $T_n$ 服从 $\operatorname{Gamma}(n,\lambda)$.
\end{proposition}
\begin{proof}
$T_n>t$ 当且仅当 $N(t)\le n-1$. 因而对 $t\ge0$,
\[
 \Prob(T_n>t)=\ee^{-\lambda t}
       \sum_{k=0}^{n-1}\frac{(\lambda t)^k}{k!}.
\]
求导时相邻项相消, 留下
\[
 f_{T_n}(t)=-\frac{\dd}{\dd t}\Prob(T_n>t)
 =\lambda\ee^{-\lambda t}\frac{(\lambda t)^{n-1}}{(n-1)!}.
\]
这就是 $\operatorname{Gamma}(n,\lambda)$ 密度.
\end{proof}

同一密度也可以从离散等待的极限看出. 每单位时间做 $N$ 次成功率为 $\lambda/N$ 的独立试验, 第 $n$ 次成功发生在第 $k$ 次试验的概率为
\[
 \binom{k-1}{n-1}(\lambda/N)^n(1-\lambda/N)^{k-n}.
\]
固定 $n$ 和 $x>0$, 令 $k/N\to x$, 上式乘以 $N$ 趋于
$\lambda^n x^{n-1}\ee^{-\lambda x}/(n-1)!$. 乘 $N$ 是因为一个时间格的宽度为 $1/N$: 点质量除以格宽才对应密度. 严格的累计概率极限也可由
$\Prob(T_n>t)=\Prob(\Bin(\lfloor Nt\rfloor,\lambda/N)<n)$ 和二项到 Poisson 的极限得到.

现在反过来, 取独立指数间隔 $X_1,X_2,\ldots$, 并记 $S_n=X_1+\cdots+X_n$. 下一章将由卷积证明 $S_n$ 的上述 gamma 密度. 令 $N(t)$ 为 $S_n\le t$ 的最大 $n$, 则对 $n\ge1$,
\begin{align*}
 \Prob(N(t)=n)
 &=\int_0^t f_{S_n}(s)\Prob(X_{n+1}>t-s)\dd s\\
 &=\int_0^t\frac{\lambda^n s^{n-1}\ee^{-\lambda s}}{(n-1)!}
                 \ee^{-\lambda(t-s)}\dd s
 =\ee^{-\lambda t}\frac{(\lambda t)^n}{n!}.
\end{align*}
$n=0$ 时概率为 $\Prob(X_1>t)=\ee^{-\lambda t}$. 这些概率之和为 $1$, 也表明有限时间内不会产生无穷多个事件. 无记忆性使任意观察时刻的剩余间隔重新具有指数分布, 所以未来计数与过去独立; 这补足了从指数间隔到 Poisson 过程的联系.

整数 shape 还具有体积解释. 正数 $x_1,\ldots,x_{n-1}$ 满足 $x_1+\cdots+x_{n-1}\le t$ 的区域体积为 $t^{n-1}/(n-1)!$. 对 $n-1=1$ 它是区间长度, 更高维由对最后一个坐标积分归纳得到. 这解释了 gamma 密度中 $t^{n-1}$ 的来源, 而不把超平面的欧氏面积误认为同一体积.

\originalsection{Cauchy 分布与不存在的期望}

\begin{definition}[标准 Cauchy 分布]
标准 Cauchy 密度为
\[
 f_C(x)=\frac1{\pi(1+x^2)},\qquad x\in\R.
\]
\end{definition}

让角度 $\Theta$ 在 $(-\pi/2,\pi/2)$ 上均匀, 并令 $C=\tan\Theta$. 正切在这个区间严格递增, 因而
\[
 F_C(x)=\Prob(\Theta\le\arctan x)
 =\frac12+\frac1\pi\arctan x.
\]
求导得到 Cauchy 密度; $\arctan x$ 在两端趋于 $\pm\pi/2$, 同时核对总概率为 $1$. 几何上, 从 $(0,1)$ 发出的随机方向射线落在水平轴上的位置就是这个变量.

\begin{center}
\begin{tikzpicture}[scale=.85,>=Stealth]
 \draw[->] (-1.3,0)--(2.2,0) node[right] {$x$};
 \draw[dashed] (0,0)--(0,1);
 \draw[thick,->] (0,1)--(1.15,0);
 \fill (0,1) circle(2pt) node[above] {$(0,1)$};
 \fill (1.15,0) circle(2pt) node[below] {$C=\tan\Theta$};
 \draw (0,.75) arc[start angle=-90,end angle=-41,radius=.25];
 \node at (.17,.64) {$\Theta$};
 \node[left] at (0,.5) {$1$};
\end{tikzpicture}
\end{center}

对位置 $a\in\R$ 和尺度 $b>0$, $Y=a+bC$ 的密度为
\[
 f_Y(y)=\frac{b}{\pi((y-a)^2+b^2)}.
\]
$a$ 是对称中心和中位数, $b$ 是尺度, 不能称它们为均值与标准差.

\begin{proposition}
标准 Cauchy 变量没有有限期望, 其通常的期望也未定义; 二阶矩为无穷大. 更精确地, 对 $r\ge0$, $\E|C|^r<\infty$ 当且仅当 $r<1$.
\end{proposition}
\begin{proof}
在 $x\ge1$ 上,
\[
 \frac12x^{r-2}\le\frac{x^r}{1+x^2}\le x^{r-2}.
\]
所以尾积分收敛当且仅当 $r-2<-1$, 即 $r<1$. 在零附近, $r\ge0$ 保证可积. 特别地,
\[
 \int_0^R\frac{x}{\pi(1+x^2)}\dd x
 =\frac1{2\pi}\log(1+R^2)\longrightarrow\infty.
\]
正半轴上的期望为 $+\infty$, 负半轴上的为 $-\infty$, 因而不能相加定义为零. 对称截断积分为零只是主值, 不等于概率期望. 二阶矩由 $r=2$ 的情形发散.
\end{proof}

原课程还用 Brownian 运动解释 Cauchy 分布. Brownian 运动可理解为细网格随机游走的连续极限, 但其存在性与路径理论超出本章. 在这里明确引用它的半平面命中事实: 从 $(u,v)$、$v>0$ 出发, 首次命中水平轴的位置 $H$ 满足
\[
 \Prob(H\le x)=\frac12+\frac1\pi\arctan\frac{x-u}{v}.
\]
于是从 $(0,1)$ 出发的命中点为标准 Cauchy, 从 $(0,2)$ 出发的命中点具有 $2C$ 的分布. 此公式作为几何解释引用, 本书不以它作为其他分布结论的证明. Cauchy 分布具有稳定性: 两个独立标准 Cauchy 的和与 $2C$ 同分布. 下一章给出直接积分验证, 无须依赖 Brownian 运动. 此时若试图套用方差加法, 会写出 $4\Var(C)=2\Var(C)$; 错误在于 $\Var(C)$ 并非有限数, 所以不能作这种有限矩运算.

\originalsection{Beta 分布与未知成功率}

\begin{definition}[Beta 函数与分布]
对 $a,b>0$, 定义
\[
 B(a,b)=\int_0^1 x^{a-1}(1-x)^{b-1}\dd x.
\]
密度 $f(x)=x^{a-1}(1-x)^{b-1}/B(a,b)$ 在 $0<x<1$ 上, 其余为零的分布称为 $\operatorname{Beta}(a,b)$.
\end{definition}

条件 $a,b>0$ 分别保证两端的可积性. 两个参数都等于 $1$ 时是均匀分布. $a<1$ 或 $b<1$ 时密度可在端点无界, 但总概率仍然有限; 密度不需要处处有界.

\begin{proposition}[Beta--Gamma 恒等式]
\[
 B(a,b)=\frac{\Gamma(a)\Gamma(b)}{\Gamma(a+b)}\quad(a,b>0).
\]
\end{proposition}
\begin{proof}
将两个 gamma 积分相乘:
\[
 \Gamma(a)\Gamma(b)=\int_{u>0,v>0}
      u^{a-1}v^{b-1}\ee^{-(u+v)}\dd u\dd v.
\]
选择总量 $s=u+v$ 与比例 $r=u/(u+v)$, 是为了把指数中的总量与幂函数中的比例分离. 逆变换为 $u=sr$, $v=s(1-r)$, 区域变为 $s>0$, $0<r<1$. 其面积放大率为
\[
 \left|\det\begin{pmatrix}r&s\\1-r&-s\end{pmatrix}\right|=s.
\]
代入并分开非负积分,
\[
 \Gamma(a)\Gamma(b)
 =\left(\int_0^\infty s^{a+b-1}\ee^{-s}\dd s\right)
  \left(\int_0^1r^{a-1}(1-r)^{b-1}\dd r\right)
 =\Gamma(a+b)B(a,b).
\]
多元换元的条件及这个面积因子的作用, 将在下一章系统说明.
\end{proof}

对于 $P\sim\operatorname{Beta}(a,b)$, 把 $x^k$ 乘入密度, 得
\[
 \E[P^k]=\frac{B(a+k,b)}{B(a,b)}\quad(k\ge0).
\]
结合 gamma 递推式,
\[
 \E P=\frac a{a+b},\quad
 \E P^2=\frac{a(a+1)}{(a+b)(a+b+1)},\quad
 \Var(P)=\frac{ab}{(a+b)^2(a+b+1)}.
\]
这里均值由两个参数的比例决定, 而 $a+b$ 增大时分布更集中.

现在把成功率本身视为随机变量 $P$. 假设先从 $[0,1]$ 均匀选择 $P$, 再在给定 $P=p$ 的条件下独立做 $n$ 次成功率为 $p$ 的试验. 记成功数为 $H$. 对一个小区间 $\dd p$, 同时观察到 $P$ 落在该区间和 $H=h$ 的概率为
\[
 \binom nh p^h(1-p)^{n-h}\dd p.
\]
因此条件于 $H=h$ 的密度是把这式除以总概率:
\[
 f_{P\mid H=h}(p)=
 \frac{p^h(1-p)^{n-h}}{\int_0^1u^h(1-u)^{n-h}\dd u}
 =\frac{p^h(1-p)^{n-h}}{B(h+1,n-h+1)}.
\]
这就是 $\operatorname{Beta}(h+1,n-h+1)$. 条件事件 $H=h$ 的概率为正, 所以此处是普通条件概率, 不涉及对概率为零的事件 $P=p$ 直接相除.

\begin{example}
成功率先均匀选择, 随后进行 $8$ 次条件独立试验, 观察到 $5$ 次成功. 求下一次成功的预测概率和成功率的条件方差.
\end{example}
\begin{solution}
后验成功率为 $\operatorname{Beta}(6,4)$. 给定成功率时下一次成功概率为 $P$, 所以去掉条件后的预测概率为后验均值
\[
 \Prob(\text{下一次成功}\mid H=5)=\E[P\mid H=5]=\frac6{10}=\frac35.
\]
后验方差为 $6\cdot4/(10^2\cdot11)=6/275$. 它描述未知参数的剩余不确定性, 不是一次试验指示变量的方差.
\end{solution}

若先验改为 $\operatorname{Beta}(a,b)$, 用同一乘法得到后验
$\operatorname{Beta}(a+h,b+n-h)$, 预测成功概率为 $(a+h)/(a+b+n)$. 均匀先验仅是 $a=b=1$ 的特例, 不是所谓“不知道任何事情”唯一可能的数学表达.

\originalsection{习题与解答}

\begin{exercise}
设 $X$ 的密度为 $c(1-x)$, $0<x<1$, 其余为零. 求 $c$, $F_X$, 均值、方差和 $\Prob(X>1/2)$. 若以概率 $1/4$ 直接取 $0$, 以概率 $3/4$ 按此密度取值形成 $Y$, 求 $F_Y$ 与 $\Prob(0\le Y\le1/2)$.
\end{exercise}
\begin{solution}
归一条件为 $c\int_0^1(1-x)\dd x=c/2=1$, 所以 $c=2$. 在 $(0,1)$ 上 $F_X(t)=2t-t^2$, 两端外分别为 $0,1$. 由
\[
 \E X=2\int_0^1(x-x^2)\dd x=\frac13,\qquad
 \E X^2=2\int_0^1(x^2-x^3)\dd x=\frac16,
\]
得 $\Var(X)=1/18$. 右尾概率 $1-F_X(1/2)=1/4$. 对混合分布,
\[
 F_Y(t)=\begin{cases}0,&t<0,\\
 1/4+(3/4)(2t-t^2),&0\le t<1,\\1,&t\ge1.\end{cases}
\]
所以 $\Prob(0\le Y\le1/2)=F_Y(1/2)-F_Y(0-)=13/16$.
\end{solution}

\begin{exercise}
三个相互独立的指数钟的速率分别为 $1,2,4$. 求第一次响的时间的均值, 以及已经 $3$ 个时间单位没有钟响后, 再等待超过 $2$ 个时间单位的条件概率. 若只考虑速率为 $2$ 的钟, 已等 $3$ 单位后其条件总寿命的均值是多少?
\end{exercise}
\begin{solution}
最小等待时间 $T\sim\Exp(7)$, 所以 $\E T=1/7$. 无记忆性给出
$\Prob(T>5\mid T>3)=\ee^{-14}$. 对单独的速率 $2$ 的钟 $X$, 条件于 $X>3$, $X=3+(X-3)$, 剩余项的条件均值为 $1/2$, 所以 $\E[X\mid X>3]=7/2$. 无记忆性描述剩余寿命, 不表示总寿命的条件均值仍为 $1/2$.
\end{solution}

\begin{exercise}
事件按速率 $3$ 的 Poisson 过程发生. 设 $T$ 为第四次事件的时刻. 求 $\E T$, $\Var(T)$ 和 $\Prob(T>1)$. 再设未知成功率 $P$ 先验为 $\operatorname{Beta}(2,3)$, 观察到 $4$ 次成功、$2$ 次失败, 求后验和下一次成功的预测概率.
\end{exercise}
\begin{solution}
$T\sim\operatorname{Gamma}(4,3)$, 所以均值为 $4/3$, 方差为 $4/9$. $T>1$ 等价于 $N(1)\le3$, 故
\[
 \Prob(T>1)=\ee^{-3}\left(1+3+\frac{3^2}{2!}+\frac{3^3}{3!}\right)
 =13\ee^{-3}.
\]
后验密度正比于 $p^{2-1+4}(1-p)^{3-1+2}$, 因而为 $\operatorname{Beta}(6,5)$, 预测成功概率为 $6/11$.
\end{solution}

\begin{exercise}
若 $C$ 为标准 Cauchy 变量, 求 $\Prob(|C|\le1)$, $\Prob(C>t)$ ($t>0$), 并说明为什么由对称性不能断言 $\E C=0$.
\end{exercise}
\begin{solution}
$\arctan1=\pi/4$, 所以 $F_C(1)-F_C(-1)=1/2$. 右尾概率为 $1/2-\arctan(t)/\pi$, 当 $t\to\infty$ 时约为 $1/(\pi t)$, 显示比指数和正态尾部衰减慢得多. 正负半轴的期望分别发散到相反的无穷, 未满足绝对可积性. 对称性只能使对称截断的主值为零, 不能定义通常期望.
\end{solution}

\originalchapter{联合分布、独立性与随机变量的变换}
\label{chap:joint}
\sourcelecture{21--22}

一个随机变量的分布只描述它自身. 若要研究两个等待时间谁先结束、两次测量的和或成功率所占的比例, 就必须知道它们如何共同变化. 两个边缘分布相同, 并不意味着联合分布相同. 本章建立联合分布的语言, 由积分区域求事件概率, 再把变量变换与独立随机变量的卷积统一起来.

\originalsection{一个变量的函数}

求 $Y=g(X)$ 的分布, 最稳妥的起点是把事件 $\{Y\le y\}$ 改写成关于 $X$ 的事件. 若 $g$ 严格递增, 不等号方向保持; 若 $g$ 严格递减, 方向反转; 若 $g$ 不单调, 则可能有多个区间或多个逆像.

\begin{proposition}[单调变换]
设 $X$ 有密度 $f_X$, 取值在区间 $I$ 内, $g:I\to J$ 为严格单调的可微函数, 其导数不为零, 逆函数可微. 则 $Y=g(X)$ 在 $J$ 上的密度为
\[
 f_Y(y)=f_X(g^{-1}(y))\left|\frac{\dd}{\dd y}g^{-1}(y)\right|,
\]
在 $J$ 外为零.
\end{proposition}
\begin{proof}
当 $g$ 递增, $F_Y(y)=F_X(g^{-1}(y))$, 链式法则得到公式. 当 $g$ 递减, 由于 $X$ 无点质量,
\[
 F_Y(y)=\Prob(X\ge g^{-1}(y))=1-F_X(g^{-1}(y)).
\]
求导得到负的逆函数导数, 即其绝对值. 密度必须非负, 这正是绝对值的必要性. 若 $y$ 不在像区间内, 概率为零或一, 相应密度为零.
\end{proof}

例如 $Y=X^3$ 时, $F_Y(y)=F_X(\sqrt[3]y)$ 对所有 $y$ 成立. 对 $y\ne0$, $f_Y(y)=f_X(\sqrt[3]y)/(3|y|^{2/3})$. 零点处逆函数导数可能发散, 但单点上的密度值无需规定成有限数. 公式的微分版本需要导数条件, 分布函数版本只需要单调性.

\begin{example}
若 $Z\sim\Normal(0,1)$, 求 $Q=Z^2$ 的分布函数与密度.
\end{example}
\begin{solution}
对 $q<0$, $F_Q(q)=0$. 对 $q\ge0$,
\[
 \{Z^2\le q\}=\{-\sqrt q\le Z\le\sqrt q\},\qquad
 F_Q(q)=\Phi(\sqrt q)-\Phi(-\sqrt q)=2\Phi(\sqrt q)-1.
\]
对 $q>0$ 求导得到
\[
 f_Q(q)=\frac{\phi(\sqrt q)+\phi(-\sqrt q)}{2\sqrt q}
       =\frac{\ee^{-q/2}}{\sqrt{2\pi q}}.
\]
两个逆像 $\sqrt q$ 和 $-\sqrt q$ 都要计入. 由 $\Gamma(1/2)=\sqrt\pi$, 这是 $\operatorname{Gamma}(1/2,1/2)$ 密度. 虽然密度在 $0$ 附近无界, $\Prob(Q=0)=0$.
\end{solution}

更一般地, 如果 $g$ 在若干不交区间上分别严格单调, 对每个非临界值 $y$, 将各支的密度相加:
\[
 f_Y(y)=\sum_{x:g(x)=y}\frac{f_X(x)}{|g'(x)|}.
\]
这里要求各支换元合法, 且只计入位于 $X$ 支撑集内的逆像. 若 $g$ 在一个正概率区域上为常数, 则 $Y$ 可产生点质量, 不能只用这张密度公式.

\originalsection{联合质量、联合分布与边缘分布}

\begin{definition}[联合分布]
离散变量的联合质量为 $p_{X,Y}(x,y)=\Prob(X=x,Y=y)$. 任意两个实值变量的联合分布函数为
\[
 F_{X,Y}(a,b)=\Prob(X\le a,Y\le b).
\]
它表示以 $(a,b)$ 为右上角的左下无界区域的概率.
\end{definition}

离散情形将联合质量列成矩阵时, 各行各列的和分别给出 $X,Y$ 的边缘质量:
\[
 p_X(x)=\sum_y p_{X,Y}(x,y),\qquad
 p_Y(y)=\sum_x p_{X,Y}(x,y).
\]
两维分布函数同样包含边缘分布. 因事件 $\{X\le a,Y\le b\}$ 在 $b\to\infty$ 时递增到 $\{X\le a\}$,
\[
 F_X(a)=\lim_{b\to\infty}F_{X,Y}(a,b),\qquad
 F_Y(b)=\lim_{a\to\infty}F_{X,Y}(a,b).
\]
这一步用概率对递增事件的连续性, 不需要密度存在.

\begin{proposition}[矩形概率]
对 $a_1<a_2$, $b_1<b_2$,
\begin{align*}
 &\Prob(a_1<X\le a_2,\ b_1<Y\le b_2)\\
 &\quad=F_{X,Y}(a_2,b_2)-F_{X,Y}(a_1,b_2)
       -F_{X,Y}(a_2,b_1)+F_{X,Y}(a_1,b_1).
\end{align*}
\end{proposition}
\begin{proof}
从右上角的大左下区域去掉左侧与下侧两条区域, 两者的重叠被减了两次, 需加回一次. 这是事件的容斥公式. 半开边界使式子对有原子的分布也准确.
\end{proof}

\begin{definition}[联合密度]
若存在非负函数 $f_{X,Y}$, 积分为 $1$, 满足对可测区域 $A\subset\R^2$,
\[
 \Prob((X,Y)\in A)=\iint_A f_{X,Y}(x,y)\dd x\dd y,
\]
则称其为联合密度.
\end{definition}

此时
\[
 F_{X,Y}(a,b)=\int_{-\infty}^a\int_{-\infty}^b
                 f_{X,Y}(x,y)\dd y\dd x.
\]
在密度连续的区域, 逐次求导给出
$\partial^2F_{X,Y}/(\partial a\partial b)=f_{X,Y}(a,b)$.
这句话的前提是联合分布确实由密度表示. 如果已知 $F$ 足够光滑, 混合导数非负且可积, 并能由逐次积分恢复整个 $F$, 则可以反过来取得密度. 仅仅找到一个混合导数并不足以排除被导数漏掉的奇异质量.

例如令 $U$ 均匀于 $(0,1)$, 并取 $X=U,Y=U$. 两个边缘都有密度, 但所有联合概率落在对角线上. 对角线的平面面积为零, 任何普通二维密度在其上的积分都是零, 与此处概率为 $1$ 矛盾. 所以边缘连续不保证联合密度存在.

\begin{proposition}[边缘密度与联合期望]
若 $X,Y$ 有联合密度, 则
\[
 f_X(x)=\int_{\R}f_{X,Y}(x,y)\dd y,\qquad
 f_Y(y)=\int_{\R}f_{X,Y}(x,y)\dd x.
\]
若 $g\ge0$, 或 $\iint |g(x,y)|f_{X,Y}(x,y)\dd x\dd y<\infty$, 则
\[
 \E[g(X,Y)]=\iint_{\R^2}g(x,y)f_{X,Y}(x,y)\dd x\dd y.
\]
\end{proposition}
\begin{proof}
事件 $\{X\in B\}$ 对应区域 $B\times\R$. 对非负密度先对 $y$ 积分,
\[
 \Prob(X\in B)=\int_B\left(\int_{\R}f_{X,Y}(x,y)\dd y\right)\dd x,
\]
括号内函数就是边缘密度. $Y$ 的情形相同. 期望公式先对区域的指示函数成立, 再对有限阶梯函数成立, 最后用非负逼近或正负部分相减得到, 与上一章一维公式的证明相同. 非负函数可交换积分次序; 有正负号的函数则以绝对可积性保证交换.
\end{proof}

\begin{example}
在三角形 $0<y<x<1$ 上联合密度为常数 $c$, 其余为零. 求 $c$, 边缘密度、$\Prob(Y<1/2)$ 和 $\Prob(X+Y<1)$.
\end{example}
\begin{solution}
三角形面积为 $1/2$, 所以 $c=2$. 对固定 $x$, 可取的 $y$ 从 $0$ 到 $x$, 因而
$f_X(x)=2x$, $0<x<1$. 对固定 $y$, 可取的 $x$ 从 $y$ 到 $1$, 所以 $f_Y(y)=2(1-y)$, $0<y<1$. 它们各自积分为 $1$. 于是
\[
 \Prob(Y<1/2)=\int_0^{1/2}2(1-y)\dd y=\frac34.
\]
事件 $x+y<1$ 在原支撑集内还要求 $0<y<\min(x,1-x)$. 上界在 $x=1/2$ 改变, 所以必须分段:
\[
 \Prob(X+Y<1)=2\int_0^{1/2}x\dd x+
              2\int_{1/2}^1(1-x)\dd x=\frac12.
\]
\end{solution}

\originalsection{独立性及其判别}

\begin{definition}[随机变量独立]
若对任意可测 $A,B\subset\R$,
\[
 \Prob(X\in A,Y\in B)=\Prob(X\in A)\Prob(Y\in B),
\]
则称 $X,Y$ 独立. 多个变量相互独立要求任意有限多个这样的坐标事件的交概率均分解成各自概率的乘积. 独立同分布常缩写为 i.i.d.
\end{definition}

\begin{proposition}[质量与密度的分解]
离散变量独立当且仅当 $p_{X,Y}(x,y)=p_X(x)p_Y(y)$ 对所有可能取值成立. 有联合密度的变量独立当且仅当
\[
 f_{X,Y}(x,y)=f_X(x)f_Y(y)
\]
在除去一个面积为零的集合后成立.
\end{proposition}
\begin{proof}
离散情形, 必要性取单点事件即得, 充分性对 $x\in A,y\in B$ 的质量求和.

对连续情形, 若密度分解, 在 $A\times B$ 上积分便有
\[
 \int_A\int_B f_X(x)f_Y(y)\dd y\dd x
 =\left(\int_Af_X(x)\dd x\right)
  \left(\int_Bf_Y(y)\dd y\right),
\]
从而独立. 反过来, 独立性要求每个矩形的概率等于乘积密度在该矩形上的积分. 矩形概率决定整个联合分布: 开集可用可数个网格矩形组成, 再由递增并和补集延伸到可测事件. 因而实际联合分布与乘积密度表示同一分布, 密度的唯一性给出两者几乎处处相等. 在本书的分段连续区域内部, 也可直接对
$F_{X,Y}(a,b)=F_X(a)F_Y(b)$ 求混合导数得到所需等式.
\end{proof}

密度唯一性中的“几乎处处”不可省略: 单点或零面积边界上的密度可任意更改而不改变概率. 判断独立时还要核对整个支撑区域. 上一节三角形例子的边缘乘积 $4x(1-y)$ 在正方形中处处为正, 而真实联合密度在 $y>x$ 的一半正方形上为零, 所以不独立.

若 $X,Y$ 独立且可积, 则 $\E[XY]=\E X\E Y$. 绝对可积性也由
$\E|XY|=\E|X|\E|Y|<\infty$ 保证. 若二阶矩有限, 因而 $\Cov(X,Y)=0$. 反向通常不成立: 取 $U$ 在 $(-1,1)$ 上均匀, $X=U,Y=U^2$. 对称性给出 $\E U=\E U^3=0$, 所以协方差为零. 但 $Y$ 由 $X$ 完全决定, 例如
\[
 \Prob(|X|<1/2,Y<1/4)=1/2
 \ne(1/2)(1/2),
\]
所以二者不独立.

\begin{example}
某装置每次试验独立、等概率输出 $1,\ldots,8$ 中的一个数. 直到首次出现 $\{2,5,8\}$ 中的数时停止. 设停止编号为 $J$, 停止时输出为 $K$. 求联合分布并判断独立性.
\end{example}
\begin{solution}
对 $j\ge1$ 和 $k\in\{2,5,8\}$,
\[
 \Prob(J=j,K=k)=\left(\frac58\right)^{j-1}\frac18.
\]
求和得到 $\Prob(J=j)=(5/8)^{j-1}(3/8)$, 即几何分布, 而
\[
 \Prob(K=k)=\frac18\sum_{j=1}^\infty(5/8)^{j-1}=\frac13.
\]
联合质量等于两者乘积, 所以独立. 等待次数的随机性与最后出现哪一种合格输出无关.
\end{solution}

\originalsection{竞争指数钟与连续、离散混合的联合分布}

\begin{proposition}[首个时刻与首个类型]
设有限个变量 $X_i\sim\Exp(\lambda_i)$ ($1\leq i\leq n$, $\lambda_i>0$) 相互独立, $\Lambda=\sum_{i=1}^n\lambda_i$, $T=\min_iX_i$, $K$ 为首先达到最小值的编号. 则
\[
 T\sim\Exp(\Lambda),\qquad
 \Prob(K=i)=\frac{\lambda_i}{\Lambda},
\]
且 $T,K$ 独立.
\end{proposition}
\begin{proof}
平局概率为零, 因为联合密度在直线 $x_i=x_j$ 上的积分为零. 要使第 $i$ 个钟在时刻 $t$ 首先响, 它须在 $t$ 响, 其余所有钟均须大于 $t$. 因而时间与类型的联合密度质量为
\[
 \Prob(T\in\dd t,K=i)=\lambda_i\ee^{-\lambda_i t}
                 \prod_{j\ne i}\ee^{-\lambda_jt}\dd t
 =\lambda_i\ee^{-\Lambda t}\dd t.
\]
对类型求和得到 $T$ 的密度 $\Lambda\ee^{-\Lambda t}$; 对时间积分得到 $\lambda_i/\Lambda$. 原式恰等于两者乘积, 从而独立. 这里一个变量连续、一个离散, 应对时间积分、对类型求和, 不能寻找普通的二维面积密度.
\end{proof}

更严格地, 上面的联合密度可从 $X_i=t$, $X_j=t+r_j$ 的换元得到. 它等于
\[
 \lambda_i\ee^{-\Lambda t}\prod_{j\ne i}\lambda_j\ee^{-\lambda_jr_j},
 \qquad r_j>0.
\]
因此给定首个类型 $K=i$ 后, 各未响钟按原标签的剩余时间相互独立, 分别为 $\Exp(\lambda_j)$ ($j\ne i$), 且与首个时刻 $T$ 独立. 不能无条件说这些剩余时间独立于类型, 因为留下哪些钟依赖于 $K$. 若响过的第 $i$ 个钟也启动一个新的独立 $\Exp(\lambda_i)$ 钟, 则重新凑齐所有原标签; 这整个新/剩余时间向量的条件分布不再依赖 $T,K$, 每轮重新具有同样的竞争模型. 这就证明叠加独立 Poisson 过程后, 连续的相邻间隔和各事件类型相互独立; 间隔为 $\Exp(\Lambda)$, 类型概率为 $\lambda_i/\Lambda$.

\begin{example}
三条独立通知流的速率为每小时 $2,3,5$. 求首次通知的平均等待时间, 首次通知来自第三条流的概率, 第四次通知的时间分布, 以及第三条流第两次通知在全部通知中的编号分布.
\end{example}
\begin{solution}
总速率为 $10$, 首次等待均值为 $1/10$ 小时. 类型概率为 $2/10,3/10,5/10$, 所以首次来自第三条的概率为 $1/2$. 第四次通知时刻为四个独立 $\Exp(10)$ 间隔之和, 服从 $\operatorname{Gamma}(4,10)$.

各通知的类型独立, 第三条相当于成功率为 $1/2$ 的试验. 它的第二次通知在总序列中的编号 $L$ 是负二项等待变量:
\[
 \Prob(L=k)=\binom{k-1}{1}(1/2)^2(1/2)^{k-2},\qquad k\ge2.
\]
即使此前很长时间没有任何通知, 接下来第一条通知的等待仍为 $\Exp(10)$, 类型概率仍为以上三个数. 这依赖恒定速率与独立过程的建模假设.
\end{solution}

\originalsection{平面区域积分与随机几何}

\begin{example}
若 $X,Y$ 为独立标准正态变量, 求 $\Prob(X^2+Y^2\le r^2)$ ($r\ge0$), 并求距离 $R=\sqrt{X^2+Y^2}$ 的密度.
\end{example}
\begin{solution}
联合密度为 $(2\pi)^{-1}\ee^{-(x^2+y^2)/2}$. 事件是圆盘, 密度仅依赖距离, 所以极坐标最合适. 面积元中的径向因子不可省略:
\[
 \Prob(R\le r)=\int_0^{2\pi}\int_0^r
            \frac1{2\pi}\ee^{-s^2/2}s\dd s\dd\theta
 =1-\ee^{-r^2/2}.
\]
因此 $f_R(r)=r\ee^{-r^2/2}$ 对 $r>0$ 成立, 其余为零. 半径为 $1$ 时概率为 $1-\ee^{-1/2}\approx0.3935$, 比包含该圆盘的正方形 $[-1,1]^2$ 的概率约 $0.466$ 小.
\end{solution}

\begin{example}
平面上水平平行线的间距为 $d$. 随机投下一段长度 $\ell\le d$ 的线段. 假设较低端点的高度模 $d$ 为均匀变量 $U\in[0,d)$, 线段向上的方向角 $\Theta$ 均匀于 $[0,\pi]$, 且二者独立. 求线段与某条平行线相交的概率.
\end{example}
\begin{solution}
线段的竖直跨度为 $\ell\sin\Theta$. 因为 $\ell\le d$, 从较低端点到较高端点最多跨过一条线, 除去概率为零的端点情形. 穿越条件是
$U+\ell\sin\Theta\ge d$. 联合密度为 $1/(d\pi)$, 所以
\[
 \Prob(\text{穿越})=\frac1{d\pi}\int_0^\pi
            \int_{d-\ell\sin\theta}^d\dd u\dd\theta
 =\frac{\ell}{d\pi}\int_0^\pi\sin\theta\dd\theta
 =\frac{2\ell}{d\pi}.
\]
例如 $\ell=1,d=2$ 时概率为 $1/\pi$. 随机投掷的口头描述不足以指定模型; 这里的均匀位置、均匀角度与独立性是计算所需的明确假设.
\end{solution}

\originalsection{二维变量变换与 Jacobian}

一维换元中, 逆函数导数把新坐标的长度变为旧坐标的长度. 二维换元需要相应的面积放大率. 若旧坐标 $(x,y)$ 表示成新坐标 $(u,v)$ 的可微函数, 在一个很小的矩形上, 两条边分别近似变为向量
$(x_u,y_u)\dd u$ 与 $(x_v,y_v)\dd v$. 它们形成的平行四边形面积为
\[
 \left|\det\begin{pmatrix}x_u&x_v\\y_u&y_v\end{pmatrix}\right|
      \dd u\dd v.
\]
行列式的符号表示方向, 绝对值才表示面积.

\begin{theorem}[二维密度换元]
设 $(X,Y)$ 有联合密度 $f_{X,Y}$, 几乎所有概率集中在开区域 $D$. 设 $T:D\to G$ 是一一对应的连续可微变换, 逆变换 $(x,y)=T^{-1}(u,v)$ 也连续可微, Jacobian 行列式处处非零. 则 $(U,V)=T(X,Y)$ 在 $G$ 上的密度为
\[
 f_{U,V}(u,v)=f_{X,Y}(x(u,v),y(u,v))
       \left|\det\frac{\partial(x,y)}{\partial(u,v)}\right|.
\]
在 $G$ 外密度为零. 概率为零的边界可另外处理.
\end{theorem}
\begin{proof}
对 $A\subset G$, 事件 $(U,V)\in A$ 等价于 $(X,Y)\in T^{-1}(A)$. 在逆像区域积分, 再使用微积分中的多元换元定理,
\[
 \Prob((U,V)\in A)=\iint_{T^{-1}(A)}f_{X,Y}(x,y)\dd x\dd y
 =\iint_A f_{X,Y}(T^{-1}(u,v))|\det DT^{-1}(u,v)|\dd u\dd v.
\]
与密度定义比较即得公式. 所调用的多元换元定理是上述局部面积近似的积分版本; 它作为微积分工具使用, 本书不另建其分析学证明.
\end{proof}

计算必须同时写出逆变换、像区域和 Jacobian. 漏掉区域会给不可能的取值赋予密度. 若变换不是一一对应, 应拆成若干可逆分支并将各支贡献相加; 不能任选一个逆像. 极坐标在 $r>0$, $0<\theta<2\pi$ 上去掉一条射线后可逆, 而原点与射线是零面积集合, 这就是前面圆盘积分可合法使用极坐标的原因.

\begin{example}
若 $X,Y$ 独立且分别为 $\operatorname{Gamma}(a,\lambda)$、$\operatorname{Gamma}(b,\lambda)$, 求总量 $S=X+Y$ 和比例 $V=X/(X+Y)$ 的联合分布.
\end{example}
\begin{solution}
旧区域为 $x>0,y>0$. 逆变换为 $x=sv$, $y=s(1-v)$, 新区域为 $s>0$, $0<v<1$, 逆 Jacobian 绝对值为 $s$. 因而
\begin{align*}
 f_{S,V}(s,v)
 &=\frac{\lambda^{a+b}}{\Gamma(a)\Gamma(b)}
     (sv)^{a-1}[s(1-v)]^{b-1}\ee^{-\lambda s}s\\
 &=\left[\frac{\lambda^{a+b}s^{a+b-1}\ee^{-\lambda s}}{\Gamma(a+b)}\right]
   \left[\frac{v^{a-1}(1-v)^{b-1}}{B(a,b)}\right].
\end{align*}
两个括号都是归一密度. 所以 $S\sim\operatorname{Gamma}(a+b,\lambda)$,
$V\sim\operatorname{Beta}(a,b)$, 且总量与比例独立. 相同速率使指数项只依赖 $s$; 若速率不同, 一般会留下同时依赖 $s,v$ 的项, 不再得到上述独立性.
\end{solution}

\begin{example}
若 $X,Y$ 为独立标准正态变量, 求比值 $Q=X/Y$ 的分布.
\end{example}
\begin{solution}
$\Prob(Y=0)=0$, 所以比值几乎必然有定义. 取 $V=Y$, 逆变换为 $x=qv,y=v$, Jacobian 绝对值为 $|v|$. 分别在 $v>0$ 和 $v<0$ 上可逆, 求和得到
\[
 f_{Q,V}(q,v)=\frac{|v|}{2\pi}\ee^{-(1+q^2)v^2/2}.
\]
对 $v$ 积分,
\[
 f_Q(q)=\frac1\pi\int_0^\infty v\ee^{-(1+q^2)v^2/2}\dd v
       =\frac1{\pi(1+q^2)}.
\]
所以独立正态变量的比值为标准 Cauchy. 分子、分母各自有全部有限矩, 但除法可因分母接近零而产生很重的尾部.
\end{solution}

\originalsection{独立和的卷积与积分上下界}

\begin{theorem}[连续卷积]
若 $X,Y$ 独立且具有密度, 则 $S=X+Y$ 的密度可取为
\[
 f_S(s)=\int_{\R} f_X(s-y)f_Y(y)\dd y
       =\int_{\R} f_X(x)f_Y(s-x)\dd x.
\]
这个积分称为两个密度的卷积.
\end{theorem}
\begin{proof}
事件 $X+Y\le s$ 对应半平面 $x+y\le s$, 因而
\[
 F_S(s)=\int_{\R}\int_{-\infty}^{s-y}f_X(x)f_Y(y)\dd x\dd y
       =\int_{\R}F_X(s-y)f_Y(y)\dd y.
\]
在足够光滑且可交换求导与积分的情况下, 对 $s$ 求导即可得到卷积. 为避免无条件交换微分, 对任意 $a<b$ 直接作 $z=x+y$ 换元:
\begin{align*}
 \Prob(a<S\le b)
 &=\int_{\R}\int_{a-y}^{b-y}f_X(x)f_Y(y)\dd x\dd y\\
 &=\int_{\R}\int_a^b f_X(z-y)f_Y(y)\dd z\dd y\\
 &=\int_a^b\left(\int_{\R}f_X(z-y)f_Y(y)\dd y\right)\dd z.
\end{align*}
非负积分允许交换次序, 所以括号内确实为密度. 对全实线积分的值也为 $1$.
\end{proof}

卷积不是在直线 $x+y=s$ 上对二维密度作未经说明的面积积分; 那条直线面积为零. 实际上, 固定总和坐标 $s$ 后对另一个坐标积分, 得到的是每单位 $s$ 的概率.

若 $X$ 的支撑为 $[a,b]$, $Y$ 的支撑为 $[c,d]$, 卷积中非零项需满足 $c\le y\le d$ 及 $a\le s-y\le b$. 因而有效上下界为
\[
 \max(c,s-b)\le y\le\min(d,s-a).
\]
若下界大于上界, 密度为零. 当支撑是正半轴时, $s>0$ 的积分区间为 $0<y<s$; $s\le0$ 时为空.

\begin{example}
设 $X$ 在 $[0,2]$ 上均匀, $Y$ 在 $[0,1]$ 上均匀, 二者独立. 求 $S=X+Y$ 的密度.
\end{example}
\begin{solution}
联合密度的乘积为 $1/2$ 在相应矩形上. 有效区间是
$[0,1]\cap[s-2,s]$. 因此密度等于该区间长度的一半:
\[
 f_S(s)=\begin{cases}
  s/2,&0<s<1,\\
  1/2,&1\le s\le2,\\
  (3-s)/2,&2<s<3,\\
  0,&\text{其余}.
 \end{cases}
\]
梯形面积为 $1/4+1/2+1/4=1$. 若两个变量都在 $[0,1]$ 上均匀, 同样计算得到 $s$ ($0<s<1$) 和 $2-s$ ($1<s<2$) 组成的三角密度.
\end{solution}

\begin{center}
\begin{tikzpicture}[xscale=1.7,yscale=2.4,>=Stealth]
 \draw[->] (-.2,0)--(3.4,0) node[right] {$s$};
 \draw[->] (0,-.06)--(0,.75) node[above] {$f_S(s)$};
 \draw[very thick] (0,0)--(1,.5)--(2,.5)--(3,0);
 \draw[dashed] (1,0)--(1,.5) (2,0)--(2,.5);
 \foreach \x in {1,2,3}\node[below] at (\x,0) {$\x$};
 \node[left] at (0,.5) {$1/2$};
\end{tikzpicture}
\end{center}

\begin{proposition}[同速率 Gamma 的可加性]
若 $X\sim\operatorname{Gamma}(a,\lambda)$ 与
$Y\sim\operatorname{Gamma}(b,\lambda)$ 独立, 则
$X+Y\sim\operatorname{Gamma}(a+b,\lambda)$.
\end{proposition}
\begin{proof}
前面的总量与比例变换已证明此结论. 卷积还给出一条直接计算链: 对 $s>0$,
\begin{align*}
 f_{X+Y}(s)
 &=\frac{\lambda^{a+b}\ee^{-\lambda s}}{\Gamma(a)\Gamma(b)}
    \int_0^s(s-y)^{a-1}y^{b-1}\dd y\\
 &=\frac{\lambda^{a+b}\ee^{-\lambda s}s^{a+b-1}}{\Gamma(a)\Gamma(b)}
    \int_0^1(1-v)^{a-1}v^{b-1}\dd v\\
 &=\frac{\lambda^{a+b}s^{a+b-1}\ee^{-\lambda s}}{\Gamma(a+b)}.
\end{align*}
第二行使用 $y=sv$, 第三行使用 Beta--Gamma 恒等式. $s\le0$ 时密度为零. 常数和幂次均保留, 所以这是完整密度验证, 不仅是“正比于”的形状判断.
\end{proof}

取 $a=b=1$ 得两个独立指数变量的和为 $\operatorname{Gamma}(2,\lambda)$; 归纳得到 $n$ 个 i.i.d. 指数间隔的和为 $\operatorname{Gamma}(n,\lambda)$. 这完成了第4章独立指数间隔构造 Poisson 计数时所需的密度证明. 不同速率时通常不能简单相加形状参数. 例如 $X\sim\Exp(\lambda)$, $Y\sim\Exp(\mu)$, $\lambda\ne\mu$, 对 $s>0$,
\[
 f_{X+Y}(s)=\lambda\mu\ee^{-\lambda s}\int_0^s\ee^{(\lambda-\mu)y}\dd y
 =\frac{\lambda\mu}{\mu-\lambda}(\ee^{-\lambda s}-\ee^{-\mu s}).
\]
符号自动保证密度非负; 当 $\mu\to\lambda$ 时极限为 $\lambda^2s\ee^{-\lambda s}$.

\originalsection{正态和与旋转对称性}

\begin{theorem}[独立正态的线性组合]
若 $X_j\sim\Normal(\mu_j,\sigma_j^2)$ 相互独立, $a_j\in\R$, 且 $\sum_j a_j^2\sigma_j^2>0$, 则
\[
 \sum_{j=1}^n a_jX_j\sim
 \Normal\left(\sum_{j=1}^n a_j\mu_j,\ \sum_{j=1}^n a_j^2\sigma_j^2\right).
\]
若方差和为零, 则线性组合为常数.
\end{theorem}
\begin{proof}
先看独立标准正态 $Z_1,Z_2$. 对任意 $\theta$, 取旋转坐标
\[
 U=\cos\theta Z_1+\sin\theta Z_2,\qquad
 V=-\sin\theta Z_1+\cos\theta Z_2.
\]
旋转的一一对应性、Jacobian 绝对值为 $1$, 以及
$Z_1^2+Z_2^2=U^2+V^2$, 共同给出新密度
\[
 f_{U,V}(u,v)=\frac1{2\pi}\ee^{-(u^2+v^2)/2}=\phi(u)\phi(v).
\]
所以 $U,V$ 仍为独立标准正态. 给定 $b_1,b_2$ 不全为零, 写
$b_1=r\cos\theta$, $b_2=r\sin\theta$, $r=\sqrt{b_1^2+b_2^2}$, 则
$b_1Z_1+b_2Z_2=rU\sim\Normal(0,r^2)$.

现在把 $X_1,X_2$ 分别写成均值加标准差乘以标准正态变量, 平移均值并取 $b_j=a_j\sigma_j$, 就证明两个变量的公式. 对更多变量归纳: 已相加的一组变量的函数与其余变量独立, 因为独立性对坐标事件成立, 对由这些坐标决定的事件也成立. 每次使用两个变量的结果, 均值相加、方差相加, 得到所述结论.
\end{proof}

也可由卷积配方核对常数. 对均值为零、方差分别为 $u,v>0$ 的两个独立正态变量,
\[
 \frac{(s-y)^2}{u}+\frac{y^2}{v}
 =\frac{u+v}{uv}\left(y-\frac{vs}{u+v}\right)^2
   +\frac{s^2}{u+v}.
\]
卷积中先提取 $\ee^{-s^2/(2(u+v))}$, 剩余高斯积分为
$\sqrt{2\pi uv/(u+v)}$. 与原系数 $1/(2\pi\sqrt{uv})$ 相乘, 得
\[
 f_{X+Y}(s)=\frac1{\sqrt{2\pi(u+v)}}\ee^{-s^2/(2(u+v))}.
\]
旋转证明解释结构, 配方计算则直接检查密度. 二者都不需要用正态近似替代精确分布证明.

\begin{example}
独立测量 $X\sim\Normal(5,4)$, $Y\sim\Normal(1,9)$. 求 $D=2X-Y$ 的分布与 $\Prob(D>14)$.
\end{example}
\begin{solution}
均值为 $2\cdot5-1=9$, 方差为 $2^2\cdot4+(-1)^2\cdot9=25$. 所以
$D\sim\Normal(9,25)$, 右尾概率为 $1-\Phi((14-9)/5)=1-\Phi(1)\approx0.1587$.
负系数改变均值方向, 方差中的系数仍须平方.
\end{solution}

\originalsection{离散卷积与 Cauchy 和}

对独立离散变量,
\[
 \Prob(X+Y=k)=\sum_j\Prob(X=k-j)\Prob(Y=j).
\]
这是与连续卷积相同的分解: 固定总和, 枚举所有可能的组成方式. 例如同成功率的独立二项变量相加后为 $\Bin(m+n,p)$, 因为可以将两组独立 Bernoulli 试验连接成一组. 独立 $\Pois(\lambda)$ 与 $\Pois(\mu)$ 的和满足
\begin{align*}
 \Prob(X+Y=k)
 &=\ee^{-(\lambda+\mu)}\sum_{j=0}^k
      \frac{\lambda^{k-j}\mu^j}{(k-j)!j!}\\
 &=\ee^{-(\lambda+\mu)}\frac{(\lambda+\mu)^k}{k!},
\end{align*}
所以为 $\Pois(\lambda+\mu)$. 这里的可加性既可由计数解释, 也可由二项式展开核对.

$n$ 个独立几何等待变量的和 $G$ 记录第 $n$ 次成功的试验编号. 对 $k\ge n$, 第 $k$ 次须成功, 前 $k-1$ 次恰有 $n-1$ 次成功, 因而
\[
 \Prob(G=k)=\binom{k-1}{n-1}p^n(1-p)^{k-n}.
\]
这是负二项等待分布, 与 $n$ 个指数间隔之和的 gamma 分布相对应. 几何变量取值从 $1$ 开始是此公式的重要约定.

\begin{proposition}[标准 Cauchy 的稳定性]
若 $C_1,C_2$ 为独立标准 Cauchy 变量, 则 $C_1+C_2$ 与 $2C_1$ 同分布.
\end{proposition}
\begin{proof}
卷积为
\[
 f_{C_1+C_2}(s)=\frac1{\pi^2}\int_{\R}
     \frac{\dd y}{(1+y^2)(1+(s-y)^2)}.
\]
当 $s\ne0$, 记 $B=1/(s^2+4)$, $A=2/[s(s^2+4)]$. 直接通分得到
\[
 \frac1{(1+y^2)(1+(s-y)^2)}
 =\frac{Ay+B}{1+y^2}+\frac{A(s-y)+B}{1+(s-y)^2}.
\]
先在 $[-R,R]$ 上积分. 两个常数项的积分在 $R\to\infty$ 时各趋 $B\pi$. 第一个奇函数项的对称积分为零; 第二个线性项积分为
\[
 \frac A2\log\frac{1+(s+R)^2}{1+(s-R)^2}\longrightarrow0.
\]
不能把两个分别非绝对可积的线性项未经截断就拆成无穷区间积分; 此处共同截断避免了这一问题. 所以原积分为 $2\pi/(s^2+4)$.

当 $s=0$, 用 $y=\tan\theta$,
\[
 \int_{\R}\frac{\dd y}{(1+y^2)^2}
 =\int_{-\pi/2}^{\pi/2}\cos^2\theta\dd\theta=\frac\pi2,
\]
同样符合公式. 最终密度为 $2/[\pi(s^2+4)]$, 这正是 $2C_1$ 的密度.
\end{proof}

稳定性体现 Cauchy 的重尾特点. 两个变量的平均值仍是标准 Cauchy, 不能期待它像有限方差变量的平均那样集中. 上一章关于 Brownian 命中点的解释与这个精确积分结论一致, 但此处没有使用未建立的随机过程理论.

\originalsection{习题与解答}

\begin{exercise}
在区域 $0<x<y<2$ 上联合密度为 $c$, 其余为零. 求 $c$, 两个边缘密度和 $\Prob(X<1,Y>1)$. 判断 $X,Y$ 是否独立.
\end{exercise}
\begin{solution}
三角形面积为 $2$, 所以 $c=1/2$. 固定 $x$ 时 $y\in(x,2)$, 固定 $y$ 时 $x\in(0,y)$, 因而
\[
 f_X(x)=\frac{2-x}{2}\quad(0<x<2),\qquad
 f_Y(y)=\frac y2\quad(0<y<2).
\]
事件 $0<x<1<y<2$ 是面积为 $1$ 的矩形, 完全落在原区域内, 所以概率为 $1/2$. 两个边缘事件的概率分别为 $3/4,3/4$, 乘积 $9/16\ne1/2$, 所以不独立.
\end{solution}

\begin{exercise}
独立 $X\sim\Exp(2)$、$Y\sim\Exp(3)$. 求 $T=\min(X,Y)$, 首先达到最小值的变量为 $X$ 的概率, 以及 $S=X+Y$ 的密度. $T$ 与首个身份是否独立?
\end{exercise}
\begin{solution}
$T\sim\Exp(5)$, $X$ 首先结束的概率为 $2/5$. 联合时间与身份密度为 $2\ee^{-5t}$ 和 $3\ee^{-5t}$, 均分解为 $5\ee^{-5t}$ 乘以身份概率, 所以独立. 对总和, 支撑迫使 $0<y<s$, 因而
\[
 f_S(s)=\int_0^s2\ee^{-2(s-y)}3\ee^{-3y}\dd y
 =6(\ee^{-2s}-\ee^{-3s})\quad(s>0),
\]
其余为零. 总和一般不为指数分布, 最小值与总和不能混淆.
\end{solution}

\begin{exercise}
独立 $X\sim\operatorname{Gamma}(2,1)$、$Y\sim\operatorname{Gamma}(3,1)$. 求 $S=X+Y$, $V=X/(X+Y)$ 的分布、独立性, 以及 $\Prob(V<1/2)$.
\end{exercise}
\begin{solution}
由总量与比例变换, $S\sim\operatorname{Gamma}(5,1)$,
$V\sim\operatorname{Beta}(2,3)$, 且独立. 由于
$B(2,3)=\Gamma(2)\Gamma(3)/\Gamma(5)=1/12$, 比例密度为 $12v(1-v)^2$. 因而
\[
 \Prob(V<1/2)=12\left[\frac{v^2}{2}-\frac{2v^3}{3}+\frac{v^4}{4}\right]_0^{1/2}
 =\frac{11}{16}.
\]
这里不能先假设比例均匀; 仅当两个 shape 参数均为 $1$ 时比例才均匀.
\end{solution}

\begin{exercise}
独立 $U,V$ 均匀于 $(0,1)$. 求乘积 $W=UV$ 的分布函数与密度, 再求比值 $Q=U/V$ 的密度.
\end{exercise}
\begin{solution}
对 $0<w<1$, 固定 $v$. 当 $v\le w$, 所有 $u\in(0,1)$ 都满足 $uv\le w$; 当 $v>w$, 则需 $u\le w/v$. 所以
\[
 F_W(w)=\int_0^w1\dd v+\int_w^1\frac wv\dd v
 =w-w\log w,
\]
密度为 $-\log w$, $0<w<1$. 区间外分布函数分别为 $0,1$. 密度在零附近无界, 但 $\int_0^1-\log w\dd w=1$.

对比值, 取新变量 $(q,v)$, 逆变换为 $u=qv$, $v=v$, Jacobian 为 $v>0$. 旧支撑要求 $0<v<\min(1,1/q)$ 且 $q>0$, 因而
\[
 f_Q(q)=\int_0^{\min(1,1/q)}v\dd v
 =\begin{cases}1/2,&0<q\le1,\\1/(2q^2),&q>1,\\0,&q\le0.\end{cases}
\]
两段积分各为 $1/2$, 总概率为 $1$. 支撑不再是原来的单位正方形, 必须随变换重新确定.
\end{solution}

\originalchapter{期望的线性性、次序统计量与协方差}
\sourcelecture{23--24}
计算一个和的平均值, 通常不必先求这个和的完整分布. 关键是把复杂计数拆成容易求平均值的小量. 但平均值相加与方差相加有不同条件: 前者不要求独立, 后者必须检查交叉项. 本章先建立这一区别, 再讨论排序后的样本及相关性.

\originalsection{期望的线性性与尾概率}
\begin{theorem}[线性性]
若 $X_1,\ldots,X_n$ 可积, 即 $\E|X_i|<\infty$, 而 $a_1,\ldots,a_n$ 是实常数, 则
\[\E\left[\sum_{i=1}^n a_iX_i\right]=\sum_{i=1}^n a_i\E X_i.\]
这里允许各随机变量彼此依赖.
\end{theorem}
\begin{proof}
离散情形中, 设两变量联合质量为 $p(x,y)$. 因为
$\sum_{x,y}|ax+by|p(x,y)\le |a|\E|X|+|b|\E|Y|<\infty$, 可交换绝对收敛级数的求和顺序. 因而
\[
\begin{aligned}
\E(aX+bY)&=\sum_{x,y}(ax+by)p(x,y)\\
&=a\sum_x x\sum_y p(x,y)+b\sum_y y\sum_x p(x,y)\\
&=a\E X+b\E Y.
\end{aligned}
\]
若有联合密度, 同样把求和换成二重积分; 绝对可积保证交换次序合法. 对任意类型的变量, 期望是对同一概率空间的积分, 积分的有限线性性给出相同结论. 重复二变量结论即得有限和公式. 这里使用的是积分的线性性, 而不是联合分布的分解.
\end{proof}

\begin{definition}[指标变量]
事件 $A$ 的指标变量 $\ind_A$ 在 $A$ 发生时取1, 否则取0. 它的期望为 $\E\ind_A=\Prob(A)$.
\end{definition}
若 $N$ 统计一组事件发生的次数, 写成 $N=\sum_i\ind_{A_i}$ 后便得到 $\E N=\sum_i\Prob(A_i)$. 事件间是否独立, 对这一结论没有影响.
\begin{example}
把 $m$ 个编号包裹随机一一分配给 $m$ 个地址, 所有排列等可能. 求送对地址的包裹数 $N$ 的期望.
\end{example}
\begin{solution}
令 $I_i$ 表示第 $i$ 个包裹送到第 $i$ 个地址. 每个包裹的目的地在 $m$ 个地址中均匀分布, 故 $\E I_i=1/m$. 于是 $\E N=m(1/m)=1$. 这个平均数不随规模变化; 它并不表示每次恰好送对一个.
\end{solution}

\begin{theorem}[尾积分公式]
若 $X\ge0$, 则允许两边同为无穷的恒等式
\[\E X=\int_0^\infty\Prob(X>t)\dd t\]
成立. 若 $X$ 是非负整数变量, 则 $\E X=\sum_{k=0}^\infty\Prob(X>k)$.
\end{theorem}
\begin{proof}
每个固定 $x\ge0$ 都满足 $x=\int_0^\infty\ind_{\{t<x\}}\dd t$. 把 $x$ 换成 $X$, 非负积分可交换顺序, 得
$\E X=\int_0^\infty\E\ind_{\{X>t\}}\dd t$. 非负交换可由先限制到有限区间、再用单调逼近理解, 无须假设期望有限. 整数情形逐点有 $X=\sum_{k\ge0}\ind_{\{X>k\}}$, 用同样的非负交换得到级数式.
\end{proof}
几何上, 平均值是生存函数 $1-F_X(t)$ 下的面积. 这也解释逆分布抽样: 若 $U$ 在 $(0,1)$ 均匀分布, 定义 $Q(u)=\inf\{x:F_X(x)\ge u\}$, 则 $\Prob(Q(U)\le x)=F_X(x)$. 分布函数有跳跃或平坦部分时应使用这个广义逆, 不能假定普通反函数总存在.
\begin{example}
若等待时间 $T$ 满足 $\Prob(T>t)=\ee^{-\lambda t}$, $t\ge0$, 求其平均值.
\end{example}
\begin{solution}
直接积分得 $\E T=\int_0^\infty\ee^{-\lambda t}\dd t=1/\lambda$. 此处不需先求密度, 尾概率本身已经足够.
\end{solution}

\originalsection{排序后的样本}
\begin{definition}[次序统计量]
将独立同分布的 $X_1,\ldots,X_n$ 从小到大排列为 $X_{(1)}\le\cdots\le X_{(n)}$, 称这些变量为次序统计量. 若共同分布有密度, 两个样本相等的概率为0, 故排序几乎必然严格.
\end{definition}
设共同分布函数为 $F$, 密度为 $f$. 最大值不超过 $x$ 等价于所有样本不超过 $x$, 因而
\[\Prob(X_{(n)}\le x)=F(x)^n,\qquad
\Prob(X_{(1)}>x)=(1-F(x))^n.\]
对这些函数求导, 得最大值密度 $nf(x)F(x)^{n-1}$, 最小值密度 $nf(x)(1-F(x))^{n-1}$. 若样本均匀分布于 $(0,1)$, 最大值密度为 $nx^{n-1}$, 平均值为 $n/(n+1)$.

\begin{theorem}[联合排序密度与随机名次]
在 $x_1<\cdots<x_n$ 上, 排序向量的联合密度是
\[n!\prod_{i=1}^nf(x_i),\]
区域外为0. 记录各原始编号对应名次的排列, 在 $n!$ 个排列中均匀分布, 且与排序向量独立.
\end{theorem}
\begin{proof}
严格有序的一个小区域有 $n!$ 个原像, 分别对应原始坐标的各排列. 每个原像的密度均为 $\prod_i f(x_i)$, 坐标置换的体积因子为1, 所以总密度乘 $n!$. 更精确地, 对有序区域中的任一集合 $B$ 和任一固定排列 $\pi$, 联合概率为 $\int_B\prod_i f(x_i)\dd x_1\cdots\dd x_n$, 而排序向量落在 $B$ 的概率为其 $n!$ 倍. 所以联合概率等于 $(1/n!)\Prob((X_{(1)},\ldots,X_{(n)})\in B)$, 这正是均匀性与独立性.
\end{proof}

第 $k$ 小的密度为
\[f_{(k)}(x)=\frac{n!}{(k-1)!(n-k)!}F(x)^{k-1}(1-F(x))^{n-k}f(x).\]
理由是选出一个样本落在 $[x,x+\dd x]$, 选出 $k-1$ 个在其左边, 剩余的在右边; 有 $n\binom{n-1}{k-1}$ 种编号选择. 也可把联合排序密度的其余坐标积分掉, 左侧有序积分为 $F(x)^{k-1}/(k-1)!$, 右侧为 $(1-F(x))^{n-k}/(n-k)!$, 得到同一公式.
\begin{example}
独立抽取七个 $(0,1)$ 均匀样本. 已知第一个样本在七个中位列第五小, 求它的条件密度、平均值及众数.
\end{example}
\begin{solution}
名次排列与排序值独立, 所求分布是 $X_{(5)}$ 的分布, 密度为 $105x^4(1-x)^2$, $0<x<1$. 一般整数参数的贝塔积分
\[\int_0^1x^{a-1}(1-x)^{b-1}\dd x=\frac{(a-1)!(b-1)!}{(a+b-1)!}\]
由分部积分递推得到: 当 $b>1$ 时, 对 $x^{a-1}$ 积分, 将 $(1-x)^{b-1}$ 求导, 得原积分等于 $(b-1)/a$ 倍参数为 $(a+1,b-1)$ 的积分, 最后归结为 $\int_0^1x^{a+b-2}\dd x=1/(a+b-1)$. 因此平均值是相邻贝塔积分之比 $a/(a+b)$, 本题 $a=5,b=3$, 平均值为 $5/8$. 密度对数的导数是 $4/x-2/(1-x)$, 在 $x=2/3$ 为0, 左正右负, 所以众数为 $2/3$. 一般 $a,b>1$ 时众数为 $(a-1)/(a+b-2)$.
\end{solution}

\originalsection{协方差与和的方差}
\begin{proposition}[独立变量的乘积]
若 $X,Y$ 独立且 $g(X),h(Y)$ 可积, 则乘积也可积, 且
$\E[g(X)h(Y)]=\E g(X)\E h(Y)$.
\end{proposition}
\begin{proof}
离散情形把联合质量分解为 $p_X(x)p_Y(y)$, 连续情形把联合密度分解为 $f_X(x)f_Y(y)$. 先对绝对值求和或积分得到 $\E|g(X)h(Y)|=\E|g(X)|\E|h(Y)|<\infty$, 再交换求和或积分并分离两个因子. 一般独立分布也有相同的乘积分解.
\end{proof}
\begin{definition}[协方差]
若 $X,Y$ 有有限二阶矩, 定义
\[\Cov(X,Y)=\E[(X-\E X)(Y-\E Y)].\]
有限二阶矩使乘积可积, 因为 $2|ab|\le a^2+b^2$. 展开得
$\Cov(X,Y)=\E(XY)-\E X\E Y$.
\end{definition}
协方差为正意味着两个相对各自平均值的偏差, 同号的贡献超过异号的贡献. 它不衡量一切形式的依赖, 只衡量这种乘积平均.
\begin{proposition}[双线性与方差公式]
协方差对交换变量对称, $\Cov(X,X)=\Var X$, 且
\[
\Cov\left(\sum_i a_iX_i,\sum_j b_jY_j\right)
=\sum_{i,j}a_ib_j\Cov(X_i,Y_j).
\]
因此
\[\Var\left(\sum_iX_i\right)=\sum_i\Var X_i+2\sum_{i<j}\Cov(X_i,X_j).\]
\end{proposition}
\begin{proof}
减去各自期望后, 第一个和变成 $\sum_i a_i(X_i-\E X_i)$, 第二个和相同处理. 相乘并用有限线性性便得双线性式. 令两组变量相同且系数全为1, 对角项是方差, 每一非对角项出现两次. 独立时乘积命题说明各交叉协方差为0; 两两不相关其实已经足够使方差相加.
\end{proof}
\begin{example}
独立的 $X_1,\ldots,X_5$ 方差均为 $v>0$. 令 $A=X_1+X_2+X_3+X_4$, $B=X_3+X_4+X_5$. 求协方差与相关系数.
\end{example}
\begin{solution}
展开协方差后, 仅重复出现的 $X_3,X_4$ 贡献 $v$, 所以 $\Cov(A,B)=2v$. 又 $\Var A=4v$, $\Var B=3v$, 相关系数为 $2/\sqrt{12}=1/\sqrt3$. 一般独立同方差变量的两个子集和, 协方差等于交集大小乘共同方差.
\end{solution}
\begin{example}
继续包裹随机分配问题, 设 $m\ge2$. 求送对数 $N$ 的方差.
\end{example}
\begin{solution}
$\Var I_i=(1/m)(1-1/m)$. 对 $i\ne j$, 同时送对这两个包裹的排列有 $(m-2)!$ 个, 故 $\E(I_iI_j)=1/[m(m-1)]$, 从而
\[\Cov(I_i,I_j)=\frac1{m(m-1)}-\frac1{m^2}=\frac1{m^2(m-1)}.\]
代入和的方差公式,
\[\Var N=m\frac1m\left(1-\frac1m\right)+m(m-1)\frac1{m^2(m-1)}=1.\]
这些指标并不独立, 忽略协方差会少算 $1/m$. $m=1$ 时 $N$ 恒等于1, 方差为0, 不属于上述计算范围.
\end{solution}

\originalsection{相关系数及其局限}
\begin{definition}
当 $0<\Var X,\Var Y<\infty$ 时, 定义
$\rho(X,Y)=\Cov(X,Y)/\sqrt{\Var X\Var Y}$. 常量变量的方差为0, 其相关系数不定义.
\end{definition}
\begin{theorem}
$-1\le\rho(X,Y)\le1$. 等号 $|\rho|=1$ 当且仅当存在 $a\ne0,b\in\R$, 使 $Y=aX+b$ 几乎必然成立. 此时 $\rho$ 等于 $a$ 的符号.
\end{theorem}
\begin{proof}
令 $U=(X-\E X)/\sqrt{\Var X}$, $V=(Y-\E Y)/\sqrt{\Var Y}$. 二者均值为0、二阶矩为1, 且 $\E UV=\rho$. 于是 $\E(U-V)^2=2-2\rho\ge0$, $\E(U+V)^2=2+2\rho\ge0$, 得上下界. 若 $\rho=1$, 非负变量 $(U-V)^2$ 平均为0, 故 $U=V$ 几乎必然; 若 $\rho=-1$ 则 $U=-V$. 两种情况都转成所述仿射关系. 反向代入协方差定义即可验证.
\end{proof}
对 $a,c\ne0$, $\rho(aX+b,cY+d)=\operatorname{sgn}(ac)\rho(X,Y)$. 正比例改变单位不会改变相关系数, 负比例会翻转方向. 独立且方差有限、非零的变量一定不相关, 逆命题不成立.
\begin{example}
令 $X$ 在 $(-1,1)$ 均匀分布, $Y=X^2$. 证明它们不相关但不独立.
\end{example}
\begin{solution}
对称性给 $\E X=\E X^3=0$, 所以 $\Cov(X,Y)=0$. 两变量方差均为正. 但事件 $A=\{|X|<1/2\}$ 与 $B=\{Y<1/4\}$ 相同, 概率为 $1/2$. 因而 $\Prob(A\cap B)=1/2\ne1/4=\Prob(A)\Prob(B)$, 它们不独立. 相关为0只排除了线性共同变化, 没有排除平方关系.
\end{solution}

\originalsection{无限平均值与有限决策}
有限线性性中的可积条件不能随意删除. $\infty-\infty$ 没有数值意义, 无限序列中每一步改善也不保证极限对象改善. 例如起初持有编号1的凭证, 第 $n$ 步交出编号 $n$ 的凭证并收到编号 $2n,2n+1$ 的凭证. 每一步净增一个, 完成 $n$ 步时持有 $n+1$ 个; 但固定编号的凭证最终都会交出, 所以逐编号极限为空集. 计数与这一极限不交换, 不能把有限阶段的净收益直接延伸到无穷阶段.

同样, 对每个有限 $n$, 先等待 $n$ 天再以 $1/n$ 的概率遭遇无限等待的策略, 都有无限期望损失. 把 $1/n$ 降为 $1/(n+1)$ 的好处乘以无穷, 再据此无限推迟, 并不是两个有限期望之间的比较. 若损失或效用无界, 应先说明采用哪一种有限截断或优化准则.

无限总量的搬移也会产生类似错觉. 若第 $n$ 堆有 $r^n$ 单位, $r>1$, 把比例 $q$ 移向左边, 剩余移向右边, 则内部第 $n$ 堆的新量为 $q r^{n+1}+(1-q)r^{n-1}$. 例如 $r=4,q=1/2$, 每堆变为原来的 $17/8$ 倍. 这没有从有限总量制造财富: 初始总量已经无穷, 边界流入流出和截断总量必须另算.
\begin{exercise}
令 $K$ 满足 $\Prob(K=k)=2^{-(k+1)}$, $k=0,1,\ldots$. 两个信封的金额为 $A=3^K$ 与 $2A$, 随机等可能打开一个, 看到金额 $V$. 计算看到 $V=1$ 和看到其他可能金额时换到另一个信封的条件平均收益, 并解释为何不能推出无条件平均收益为正.
\end{exercise}
\begin{solution}
幂 $3^k$ 与 $2\cdot3^j$ 不会相等. 所以看到 $V=3^k$ 就知道自己拿的是小信封, 换后收益为 $3^k$; 看到 $V=2\cdot3^k$ 就知道是大信封, 收益为 $-3^k$. 特别地 $V=1$ 时收益为1. 无条件收益的正、负部分期望分别为
$\frac12\sum_{k\ge0}3^k2^{-(k+1)}=\infty$, 因此收益的期望未定义, 不能通过把正、负无穷相减比较策略. 已知金额后的判断必须使用实际条件分布, 不能不加检查地说另一个金额以各半概率是当前金额的两倍或一半.
\end{solution}

\originalchapter{条件期望与分布的变换函数}
\sourcelecture{23, 25--26}
条件分布描述已获得信息以后还剩下的随机性. 条件期望则把这整个分布压缩成一个预测值. 因为信息本身也是随机获得的, 条件期望是随机变量, 而不是一个固定数字. 本章从能直接计算的联合分布入手, 再用指数函数把分布编码, 为独立和与极限定理准备工具.

\originalsection{条件分布与零概率条件}
若 $(X,Y)$ 离散, 当 $p_Y(y)>0$ 时,
\[p_{X\mid Y}(x\mid y)=\frac{p_{X,Y}(x,y)}{p_Y(y)}.\]
这就是把联合概率中固定 $y$ 的一列归一化. 连续情形不能使用 $\Prob(X\in A,Y=y)/\Prob(Y=y)$, 因为分母为0.
\begin{definition}[条件密度]
若 $(X,Y)$ 有联合密度 $f(x,y)$, 令 $f_Y(y)=\int_\R f(x,y)\dd x$. 在 $0<f_Y(y)<\infty$ 的点, 定义
\[f_{X\mid Y}(x\mid y)=\frac{f(x,y)}{f_Y(y)}.\]
它作为 $x$ 的函数非负且积分为1. 在其他 $y$ 处可任意指定一个概率密度; 这些点对 $Y$ 的分布而言总概率为0.
\end{definition}
条件密度的依据是下述重建公式: 对任意可测集合 $A,B$,
\[\Prob(X\in A,Y\in B)=\int_B\left(\int_Af_{X\mid Y}(x\mid y)\dd x\right)f_Y(y)\dd y.\]
将定义代入, 右边恢复联合密度积分. 所以可以先按 $f_Y$ 抽取 $Y$, 再按该 $y$ 对应的条件密度抽取 $X$, 得到原联合分布. 也得到边缘密度的混合公式
$f_X(x)=\int_\R f_{X\mid Y}(x\mid y)f_Y(y)\dd y$, 几乎处处成立.

在适当连续条件下, 上述定义还可从窄带条件的极限得到. 令 $g_A(v)=\int_A f(x,v)\dd x$. 若 $g_A$ 与 $f_Y$ 在 $y$ 连续且 $f_Y(y)>0$, 则
\[
\Prob(X\in A\mid |Y-y|<\eps)
=\frac{\int_{y-\eps}^{y+\eps}g_A(v)\dd v}{\int_{y-\eps}^{y+\eps}f_Y(v)\dd v}
\longrightarrow\frac{g_A(y)}{f_Y(y)}.
\]
分子分母同时除以 $2\eps$, 再用连续函数的局部平均趋于函数值即可. 没有这些局部条件时, 不能保证每个指定 $y$ 都有这种极限.
\begin{example}
在三角形 $0<x<y<1$ 上均匀取点 $(X,Y)$. 求 $X$ 在给定 $Y=y$ 时的分布.
\end{example}
\begin{solution}
面积为 $1/2$, 联合密度为2. 因而 $f_Y(y)=\int_0^y2\dd x=2y$, $0<y<1$, 条件密度为 $1/y$, $0<x<y$. 所以给定 $Y=y$, $X$ 均匀分布于 $(0,y)$, 而不是 $(0,1)$. 条件会改变取值范围, 必须在除以边缘密度以前就确认支持集.
\end{solution}

同一个零概率集合, 不同的邻域收缩方式可能给不同答案. 例如在矩形 $0<x<1,-1<y<1$ 上均匀取点. 按 $|Y|<\eps$ 收缩, 极限的 $X$ 是 $(0,1)$ 均匀分布. 若观察变量换成 $T=Y/X$, 则 $T=0$ 仍对应 $Y=0$, 但按 $|T|<\eps$ 收缩时, 对固定 $x$ 的带宽为 $2\eps x$. 对 $0<\eps<1$, 条件密度正比于 $x$, 归一化后为 $2x$. 零概率条件的含义由观察变量及其条件分布决定, 不能只凭集合相同就认定条件分布相同.

\originalsection{条件期望的对象与基本性质}
\begin{definition}
设 $X$ 可积. 离散情形定义
$m(y)=\sum_x xp_{X\mid Y}(x\mid y)$, 连续联合密度情形定义
$m(y)=\int_\R xf_{X\mid Y}(x\mid y)\dd x$.
在这些公式有意义的点上, $m(y)=\E[X\mid Y=y]$ 是一个数. 定义随机变量
$\E[X\mid Y]=m(Y)$; 它只依赖观察到的 $Y$.
\end{definition}
绝对可积保证条件绝对期望几乎处处有限. 例如连续情形
$\int\bigl(\int |x|f_{X\mid Y}(x\mid y)\dd x\bigr)f_Y(y)\dd y=\E|X|<\infty$.
因此可能失效的条件值只落在一个 $Y$ 概率为0的集合. 此外 $|m(y)|\le\E[|X|\mid Y=y]$, 故 $m(Y)$ 也可积.

\begin{theorem}[全期望与拉出已知因子]
若 $X$ 可积, 则 $\E[\E(X\mid Y)]=\E X$. 若 $h(Y)$ 有界, 则
\[\E[h(Y)X\mid Y]=h(Y)\E[X\mid Y],\qquad
\E[h(Y)X]=\E[h(Y)\E(X\mid Y)].\]
同样公式对无界 $h$ 也成立, 只要 $h(Y)X$ 可积.
\end{theorem}
\begin{proof}
连续情形, 全期望是
\[
\int_\R m(y)f_Y(y)\dd y
=\int_\R\int_\R xf(x,y)\dd x\dd y=\E X.
\]
可交换积分是因为 $\iint |x|f(x,y)\dd x\dd y<\infty$. 离散情形把积分换成求和, 由绝对收敛得到相同等式. 条件给定 $y$ 后, $h(y)$ 是常数, 所以可从关于 $x$ 的期望中提出. 对该等式再用全期望即得第二式. 无界情形先把 $h$ 截断为 $h\ind_{\{|h|\le n\}}$; 由 $|hX|$ 可积和条件积分的绝对值估计, 两边均可通过可积控制取极限. 特别地, 这也保证 $h(Y)m(Y)$ 可积.
\end{proof}
这里的全期望不是在未知条件之间简单平均, 而是以条件值本身的概率为权重. 条件均值虽可能随 $y$ 很大或很小, 对所有条件再平均仍恢复原平均值.

若不具备联合密度, 例如 $X=Y$, 可用下述刻画: $m(Y)$ 可积且对每个有界函数 $h$ 满足
$\E[h(Y)m(Y)]=\E[h(Y)X]$. 这称为关于 $Y$ 所含信息的条件期望. 它的存在在一般概率空间中需要测度论的条件期望存在定理, 本书调用该定理而不证明它; 对离散与联合密度情形, 上面的构造已直接建立存在. 刻画保证几乎处处唯一: 两个候选之差为 $d(Y)$, 取 $h(Y)=\operatorname{sgn}(d(Y))$, 便得 $\E|d(Y)|=0$.

\begin{proposition}[塔式性质]
若 $X$ 可积, $Z=g(Y)$, 即 $Z$ 提供的信息包含于 $Y$ 提供的信息, 则
\[\E[\E(X\mid Y)\mid Z]=\E(X\mid Z).\]
\end{proposition}
\begin{proof}
左边记作 $W$, 是 $Z$ 的函数. 对任意有界 $h(Z)$, 连用两次刻画,
\[\E[h(Z)W]=\E[h(Z)\E(X\mid Y)]=\E[h(g(Y))X]=\E[h(Z)X].\]
唯一性给出结论. 从更细的信息回到更粗的信息后, 得到与直接按粗信息预测相同的结果. 条件顺序不能任意交换; 这里需要明确的信息包含关系.
\end{proof}
条件期望也满足线性性和非负性, 证明分别在条件积分中用线性性及被积函数非负, 或使用刻画及唯一性. 若 $X$ 与 $Y$ 独立, 条件分布不随 $y$ 改变, 所以 $\E[X\mid Y]=\E X$. 若 $X=h(Y)$, 则已经知道 $X$, 所以 $\E[X\mid Y]=X$.

\begin{example}
独立进行30次成功概率为 $p$ 的试验, $S$ 为总成功数, $K$ 为前12次成功数. 求 $\E[S\mid K]$ 和 $\Var(S\mid K)$.
\end{example}
\begin{solution}
写 $S=K+L$, 后18次成功数 $L\sim\Bin(18,p)$ 与 $K$ 独立. 给定 $K=k$, 前段数已成为常数, 因而 $\E[S\mid K]=K+18p$, $\Var(S\mid K)=18p(1-p)$. 再平均得到 $\E S=12p+18p=30p$.
\end{solution}
\begin{example}
一批20件产品中有5件带标记, 不放回依次抽取6件. 已知前两件均带标记, 求6件中带标记数的条件期望.
\end{example}
\begin{solution}
剩余18件中有3件带标记. 后四个抽取位置的每个指标在给定条件下期望为 $3/18$. 无须假设这些指标独立, 线性性给条件期望 $2+4(3/18)=8/3$.
\end{solution}

\originalsection{条件方差与最佳平方预测}
\begin{definition}
若 $\E X^2<\infty$, 定义
\[\Var(X\mid Y)=\E[(X-\E[X\mid Y])^2\mid Y]
=\E[X^2\mid Y]-(\E[X\mid Y])^2.\]
这通常是 $Y$ 的随机函数. 平方恒等式是在每个条件分布中展开, 并非无条件方差的记号改写.
\end{definition}
令 $m=\E[X\mid Y]$. 条件积分的柯西不等式给 $m^2\le\E[X^2\mid Y]$, 故 $m$ 有有限二阶矩.
\begin{theorem}[全方差]
若 $\E X^2<\infty$, 则
\[\Var X=\Var(\E[X\mid Y])+\E[\Var(X\mid Y)].\]
\end{theorem}
\begin{proof}
由全期望,
$\E\Var(X\mid Y)=\E X^2-\E m^2$.
又 $\E m=\E X$, 所以
$\Var m=\E m^2-(\E X)^2$.
相加正好为 $\E X^2-(\E X)^2$. 第一项描述不同条件均值之间的变化, 第二项描述每个条件内部仍未消除的变化.
\end{proof}
\begin{example}
在本章三角形取点模型中, 求 $\E X$、$\Var X$.
\end{example}
\begin{solution}
给定 $Y=y$ 后, $X$ 均匀于 $(0,y)$, 所以条件均值为 $y/2$, 条件方差为 $y^2/12$. 从 $f_Y(y)=2y$ 得 $\E Y=2/3$, $\E Y^2=1/2$, $\Var Y=1/18$. 因此
\[\E X=\frac13,\qquad \Var X=\frac14\frac1{18}+\frac1{12}\frac12=\frac1{18}.\]
用边缘密度 $f_X(x)=2(1-x)$ 直接计算二阶矩为 $1/6$, 减去 $1/9$ 也得到 $1/18$.
\end{solution}
\begin{theorem}[平方损失下的最佳预测]
若 $\E X^2<\infty$, 对任何具有有限二阶矩的预测 $a(Y)$,
\[\E[(X-a(Y))^2]=\E[(X-m(Y))^2]+\E[(m(Y)-a(Y))^2].\]
因此条件期望是依赖 $Y$ 的最佳平方预测, 且最优预测几乎处处唯一.
\end{theorem}
\begin{proof}
把误差分解为 $X-m+(m-a)$, 平方展开. 交叉项期望为0, 因为 $m-a$ 是 $Y$ 的函数, 而 $\E[X-m\mid Y]=0$. 两个因子的乘积可积由二阶矩及柯西不等式保证; 可以先截断 $m-a$, 再用可积控制得到拉出因子的等式. 剩下两个平方项均非负, 第二项仅当 $a=m$ 几乎必然时为0.
\end{proof}
只允许常数预测时同样有
$\E(X-c)^2=\Var X+(\E X-c)^2$, 最优常数为 $\E X$.
若 $X,Y$ 独立、均值为0, 令 $Z=X+Y$, 方差分别为 $\sigma_X^2,\sigma_Y^2$, 则
\[\Cov(X,Z)=\sigma_X^2,\quad
\rho(X,Z)=\frac{\sigma_X}{\sqrt{\sigma_X^2+\sigma_Y^2}},\quad
\E[Z\mid X]=X,\quad\Var(Z\mid X)=\sigma_Y^2,\]
其中相关系数要求 $\sigma_X>0$ 及分母非零. 全方差将 $\Var Z$ 分成已观察到的 $X$ 的变化和未观察到的 $Y$ 的变化.

\originalsection{矩母函数及微分条件}
\begin{definition}
实随机变量 $X$ 的矩母函数为 $M_X(t)=\E\ee^{tX}$, $t\in\R$. 允许值为 $+\infty$. 其有限定义域是 $D_X=\{t:M_X(t)<\infty\}$, 且总有 $M_X(0)=1$.
\end{definition}
它将分布变成指数函数的加权平均. 对 $t>0$, $\Prob(X\ge b)>0$ 时有 $M_X(t)\ge\ee^{tb}\Prob(X\ge b)$. 对负尾也有类似估计, 所以两侧都有非零质量的变量, 在正负两个方向上均会使矩母函数至少按某个指数下界增长.
\begin{theorem}[矩的生成]
若存在 $a>0$ 使 $M_X(a),M_X(-a)<\infty$, 则对 $|t|<a$, $M_X$ 可任意次微分, 且
\[M_X^{(k)}(t)=\E[X^k\ee^{tX}],\qquad M_X^{(k)}(0)=\E X^k.\]
并且对 $|t|<a$,
$M_X(t)=\sum_{k=0}^\infty(\E X^k)t^k/k!$.
\end{theorem}
\begin{proof}
由 $\ee^{a|X|}\le\ee^{aX}+\ee^{-aX}$ 知 $\E\ee^{a|X|}<\infty$. 对任意 $0<b<a$, 多项式增长慢于指数增长, 所以每个整数 $k\ge0$ 存在常数 $C$ 使
$|x|^k\ee^{b|x|}\le C\ee^{a|x|}$.
取 $t$ 的小邻域包含于 $[-b,b]$, 差商由中值定理被 $|X|\ee^{b|X|}$ 控制, 因而可把导数移入期望. 对已取得的导数重复论证, 得所有阶的公式. 最后指数级数的绝对值和为 $\ee^{|t||X|}$, 它可积, 故可逐项求期望. 这说明交换步骤成立的原因; 仅有某些有限矩, 并不自动保证存在矩母函数的邻域.
\end{proof}
从而 $\E X=M'_X(0)$, $\Var X=M''_X(0)-(M'_X(0))^2$.
\begin{proposition}[和与仿射变换]
独立变量 $X,Y$ 在各自矩母函数有限的 $t$ 处满足
\[M_{X+Y}(t)=M_X(t)M_Y(t).\]
另有 $M_{cX+d}(t)=\ee^{dt}M_X(ct)$. 对独立同分布的 $X_1,\ldots,X_n$, 和的矩母函数为 $M_X(t)^n$.
\end{proposition}
\begin{proof}
$\ee^{t(X+Y)}=\ee^{tX}\ee^{tY}$, 独立乘积期望公式给出第一式. 第二式直接从 $\ee^{t(cX+d)}=\ee^{dt}\ee^{ctX}$ 得到. 对独立和重复第一式即可.
\end{proof}

\originalsection{常用分布与存在性}
若 $B$ 是成功概率 $p$ 的伯努利变量, 直接计算得
$M_B(t)=1-p+p\ee^t$. 二项变量为 $n$ 个独立伯努利变量之和, 故
\[M_{\Bin(n,p)}(t)=(1-p+p\ee^t)^n,\qquad t\in\R.\]
对 $X\sim\Pois(\lambda)$,
\[M_X(t)=\sum_{k=0}^\infty\ee^{tk}\ee^{-\lambda}\frac{\lambda^k}{k!}
=\ee^{-\lambda}\sum_{k=0}^\infty\frac{(\lambda\ee^t)^k}{k!}
=\exp\{\lambda(\ee^t-1)\},\qquad t\in\R.\]
这也恢复 $\E X=\Var X=\lambda$: 一阶导数为 $M_X(t)\lambda\ee^t$, 二阶导数为 $M_X(t)[\lambda^2\ee^{2t}+\lambda\ee^t]$.

若 $G\sim\Normal(0,1)$, 配方得到
\[
M_G(t)=\frac1{\sqrt{2\pi}}\int_\R\exp\left(tx-\frac{x^2}2\right)\dd x
=\ee^{t^2/2}\frac1{\sqrt{2\pi}}\int_\R\ee^{-(x-t)^2/2}\dd x
=\ee^{t^2/2}.
\]
因此 $X=\mu+\sigma G$ 的矩母函数为 $\exp(\mu t+\sigma^2t^2/2)$, 在所有实数上有限.

对率为 $\lambda>0$ 的指数变量,
\[M_X(t)=\lambda\int_0^\infty\ee^{-(\lambda-t)x}\dd x
=\frac{\lambda}{\lambda-t},\qquad t<\lambda.\]
$t=\lambda$ 时被积函数为常数, $t>\lambda$ 时增长, 两种情况积分均为无穷. 不能把有理表达式延伸到 $t>\lambda$ 当作期望, 那里甚至会得到负值.

伽马分布采用形状参数 $\alpha>0$ 与率参数 $\lambda>0$, 密度为
$f(x)=\lambda^\alpha x^{\alpha-1}\ee^{-\lambda x}/\Gamma(\alpha)$, $x>0$, 其中 $\Gamma(\alpha)=\int_0^\infty u^{\alpha-1}\ee^{-u}\dd u$. 代换 $u=(\lambda-t)x$ 得
\[M_X(t)=\left(\frac{\lambda}{\lambda-t}\right)^\alpha,\qquad t<\lambda.\]
$t\ge\lambda$ 时积分发散. 当 $\alpha=n$ 为正整数时, 该函数也是 $n$ 个独立同率指数变量之和的矩母函数. 非整数形状不能解释为非整数个变量之和, 必须用密度积分.

对首次成功试验次数 $X\sim\Geom(p)$, 支持集为 $1,2,\ldots$, 有
\[M_X(t)=\sum_{k\ge1}p(1-p)^{k-1}\ee^{tk}
=\frac{p\ee^t}{1-(1-p)\ee^t},\]
当 $0<p<1$ 时要求 $t<-\log(1-p)$; $p=1$ 时 $X=1$, 矩母函数为 $\ee^t$.

存在性不能省略. 标准柯西变量的密度为 $1/[\pi(1+x^2)]$. 当 $t>0$ 时, 正半轴上的指数增长使积分发散; 当 $t<0$ 时负半轴同理, 所以只有 $M_X(0)=1$ 有限. 另有所有整数矩都有限而矩母函数仍不在0的邻域有限的例子: 若 $X=\ee^G$, $G$ 标准正态, 则 $\E X^k=\ee^{k^2/2}$, 但 $t>0$ 时
\[\E\ee^{tX}=\frac1{\sqrt{2\pi}}\int_\R\exp(t\ee^g-g^2/2)\dd g=\infty,\]
因为在充分大的 $g$ 上指数中的 $t\ee^g-g^2/2$ 趋于无穷, 被积函数最终大于1. 因此把形式上的矩级数当作必然收敛的矩母函数会出错.

\begin{theorem}[矩母函数唯一性, 调用版]
若 $X,Y$ 的矩母函数在同一开区间 $(-a,a)$ 上有限且相等, 则 $X,Y$ 同分布.
\end{theorem}
这是概率变换的唯一性定理; 完整证明需要拉普拉斯变换的唯一性, 超出本章初等积分范围, 此处明确作为工具调用. 它要求整个邻域的函数相等, 不能仅以 $M(0)=1$, 若干导数相等, 或没有指数可积条件的全部矩相等代替. 前面的计算因此真正识别了独立和的分布: 独立泊松变量之和仍为泊松分布, 参数相加; 独立正态变量之和仍为正态分布, 均值与方差分别相加; 同率独立伽马变量的形状参数相加.
\begin{exercise}
独立的 $A\sim\Pois(2)$, $B\sim\Pois(5)$, 令 $D=A-B$. 求 $D$ 的矩母函数、均值、方差, 并判断它是否为泊松变量.
\end{exercise}
\begin{solution}
由独立及变号公式,
$M_D(t)=\exp\{2(\ee^t-1)+5(\ee^{-t}-1)\}$.
令括号内为 $L(t)$, 则 $M'_D=M_DL'$, $M''_D=M_D[(L')^2+L'']$. 在0处 $L'=-3,L''=7$, 因而均值为 $-3$, 方差为7. 它不是泊松变量, 因为存在负值的正概率, 例如 $A=0,B=1$ 的概率为 $5\ee^{-7}>0$. 变换公式必须结合支持集理解.
\end{solution}

\originalsection{特征函数与极限工具}
\begin{definition}
实随机变量 $X$ 的特征函数为
\[\varphi_X(t)=\E\ee^{\ii tX}=\E\cos(tX)+\ii\E\sin(tX),\qquad t\in\R.\]
复数期望按实部、虚部分别定义. 因 $|\ee^{\ii tX}|=1$, 它对所有分布及所有实 $t$ 都存在, 且 $|\varphi_X(t)|\le1$, $\varphi_X(0)=1$.
\end{definition}
特征函数把增长的实指数换成有界的振荡函数. 若 $X$ 有密度, $\varphi_X(t)=\int\ee^{\ii tx}f_X(x)\dd x$ 是密度的傅里叶变换; 无密度时定义仍有效. 若 $X$ 整数值, 则 $\varphi_X(t+2\pi)=\varphi_X(t)$, 特别地在 $2\pi$ 的整数倍处为1, 因为对每个整数 $k$ 都有 $\ee^{2\pi\ii k}=1$.

独立性给 $\varphi_{X+Y}(t)=\varphi_X(t)\varphi_Y(t)$, 仿射变换给 $\varphi_{cX+d}(t)=\ee^{\ii dt}\varphi_X(ct)$, 证明与矩母函数完全相同, 且不需要指数可积条件. 它在 $t$ 上连续: 若 $t_n\to t$, 被积函数逐点收敛并被常数1控制, 所以可交换极限与期望.

若 $\E|X|^k<\infty$, 则可微到 $k$ 阶, 且
\[\varphi_X^{(k)}(t)=\E[(\ii X)^k\ee^{\ii tX}],\qquad
\E X^k=\ii^{-k}\varphi_X^{(k)}(0).\]
差商由中值积分公式被 $|X|$ 控制, 重复微分时由 $|X|^j$ 控制, $j\le k$, 这些量均可积. 正确的还原因子是 $\ii^{-k}$, 例如二阶导数满足 $\varphi_X''(0)=-\E X^2$.

对正态变量, $\varphi_{\Normal(\mu,\sigma^2)}(t)=\exp(\ii\mu t-\sigma^2t^2/2)$. 不能直接把实积分中的配方平移到复数路径而不给理由. 一个实变量推导是令 $G$ 标准正态, 由可积的一阶矩与分部积分,
\[
\varphi'_G(t)=\int_\R\ii x\ee^{\ii tx}\frac{\ee^{-x^2/2}}{\sqrt{2\pi}}\dd x
=-t\varphi_G(t).
\]
这里写 $x f_G(x)=-f'_G(x)$, 分部积分的端点项因高斯衰减而为0. 配合 $\varphi_G(0)=1$, 微分方程给 $\varphi_G(t)=\ee^{-t^2/2}$, 再用仿射变换得到一般式.

\begin{definition}[依分布收敛]
若对 $F_X$ 的每个连续点 $x$, 都有 $F_{X_n}(x)\to F_X(x)$, 则称 $X_n$ 依分布收敛到 $X$, 记作 $X_n\Rightarrow X$. 不要求这些变量定义在同一概率空间.
\end{definition}
连续点的限制不可删掉. 例如常量 $X_n=1/n$ 收敛到常量0, 在 $x=0$ 却有 $F_{X_n}(0)=0$ 而 $F_X(0)=1$.
\begin{theorem}[特征函数唯一性与连续性, 调用版]
相同的特征函数确定相同的分布. 若对每个实数 $t$, $\varphi_{X_n}(t)\to\varphi_X(t)$, 则 $X_n\Rightarrow X$. 更一般地, 若逐点极限为 $\varphi$, 且 $\varphi$ 在0连续, 则 $\varphi$ 是某个概率分布的特征函数, 并有向该分布的依分布收敛.
\end{theorem}
这是勒维连续性定理与特征函数唯一性定理. 本章把它们作为明确的外部分析工具, 不把逐点收敛的计算误称为定理本身的证明. 0处连续条件防止概率质量逃向无穷: 例如 $X_n$ 在 $(-n,n)$ 均匀分布, $t\ne0$ 时 $\varphi_{X_n}(t)=\sin(nt)/(nt)\to0$, 而在0处始终为1. 极限函数在0不连续, 不对应概率分布的特征函数.

矩母函数也有连续性定理: 若某个共同邻域 $(-a,a)$ 内各 $M_{X_n}$ 有限, 且逐点收敛于该邻域内有限的 $M_X$, 则 $X_n\Rightarrow X$. 原课程所用的所有实数上有限并收敛的版本是它的较强条件特例. 这一结论同样作为变换连续性工具调用, 不对一般只有形式矩级数的情形使用.

\begin{example}
令 $X_n\sim\Bin(n,\lambda/n)$, 其中固定 $\lambda>0$, 只取 $n>\lambda$. 用变换说明其极限分布.
\end{example}
\begin{solution}
对每个实数 $t$,
\[M_{X_n}(t)=\left(1+\frac{\lambda(\ee^t-1)}n\right)^n
\longrightarrow\exp\{\lambda(\ee^t-1)\}.
\]
右边是 $\Pois(\lambda)$ 的矩母函数, 所有函数均在整个实轴有限. 因而可调用矩母函数连续性定理, 得 $X_n\Rightarrow\Pois(\lambda)$. 只有函数极限计算加上连续性定理, 才完成分布极限的论证.
\end{solution}
下一章将把独立和标准化后计算特征函数极限. 有限方差足以产生二阶局部展开, 无须矩母函数存在. 所以中心极限定理的一般有限方差版本使用特征函数更合适; 若另有指数可积条件, 矩母函数也能完成对应证明. 指数变换还可结合 $\Prob(X\ge b)\le\ee^{-tb}M_X(t)$ 研究大偏差, 但这类估计需要选取 $t>0$ 且矩母函数有限.

\originalchapter{大数定律、中心极限定理与凸性}
\sourcelecture{27,30--31}

\originalsection{平均值的偏差与收敛的含义}
设 $X_1,X_2,\ldots$ 独立同分布, 写 $S_n=\sum_{k=1}^nX_k$, $A_n=S_n/n$. 期望描述分布的重心, 而实际观察到的是随机的 $A_n$. 本章先控制平均值偏离重心的概率, 再描述偏差的形状, 最后区分一次大样本实验与一条无限样本路径.

\begin{theorem}[Markov与Chebyshev不等式]
若 $Y\geq0$, $a>0$, 则 $\Prob(Y\geq a)\leq\E Y/a$. 若 $X$ 有有限均值 $\mu$ 和方差 $\sigma^2$, 则对 $r>0$,
\[\Prob(|X-\mu|\geq r)\leq\frac{\sigma^2}{r^2}.\]
\end{theorem}
\begin{proof}
逐个样本点有 $Y\geq a\ind_{\{Y\geq a\}}$, 取期望并除以 $a$ 即得第一式. 对非负变量 $(X-\mu)^2$ 应用第一式, 门槛取 $r^2$, 就得到第二式. 第一式在期望无限时仍成立, 但不提供有用上界.
\end{proof}
这些估计只用很少的信息, 因而通常比精确分布粗. 不过, 一旦方差随样本量趋于零, 粗估计已足以证明收敛.

\begin{definition}[三种收敛]
在同一概率空间上, $Y_n$ \emph{依概率收敛}于 $Y$, 记为 $Y_n\xrightarrow{\Prob}Y$, 是指对每个 $\eps>0$, $\Prob(|Y_n-Y|>\eps)\to0$. \emph{几乎处处收敛} $Y_n\to Y$ 是指 $\Prob(\{\omega:Y_n(\omega)\to Y(\omega)\})=1$. \emph{依分布收敛} $Y_n\Rightarrow Y$ 是指在 $F_Y$ 的每个连续点 $x$, $F_{Y_n}(x)\to F_Y(x)$. 最后一种收敛只比较分布, 不要求变量定义在同一概率空间.
\end{definition}
\begin{proposition}
几乎处处收敛蕴含依概率收敛, 依概率收敛蕴含依分布收敛. 若极限是常数 $c$, 则依分布收敛于 $c$ 等价于依概率收敛于 $c$.
\end{proposition}
\begin{proof}
若几乎处处收敛, 固定 $\eps>0$, 令 $B_N=\bigcup_{n\geq N}\{|Y_n-Y|>\eps\}$. 这些事件递减, 其交集是偏差超过 $\eps$ 无限次的事件, 概率为零. 概率的从上连续性给出 $\Prob(B_N)\to0$, 从而每个 $n\geq N$ 的偏差概率都不超过它.

若依概率收敛, 对 $\delta>0$ 有
\[F_Y(x-\delta)-\Prob(|Y_n-Y|>\delta)\leq F_{Y_n}(x)\leq F_Y(x+\delta)+\Prob(|Y_n-Y|>\delta).\]
先让 $n\to\infty$, 再让 $\delta\downarrow0$, 在连续点得到结论. 最后, 若 $Y_n\Rightarrow c$, 则 $c-\eps$ 和 $c+\eps$ 都是常数分布的连续点,
\[\Prob(|Y_n-c|>\eps)\leq F_{Y_n}(c-\eps)+1-F_{Y_n}(c+\eps)\to0.\]
另一方向已证明.
\end{proof}
\begin{example}
在 $[0,1)$ 上取均匀概率. 对每个 $k\geq0$, 按次序列出 $2^k$ 个变量
$\ind_{[j/2^k,(j+1)/2^k)}$, $0\leq j<2^k$, 将各层串成序列. 每个变量等于1的概率为 $2^{-k}\to0$, 所以序列依概率趋于0. 但每个样本点在每层恰有一次取值1, 也在所有足够大层中取到0, 所以没有任何样本点上的收敛. 这说明小概率偏差仍可能不断出现.

再令 $Y$ 以相等概率取 $\pm1$, $Y_n=(-1)^nY$. 所有 $Y_n$ 与 $Y$ 同分布, 因而依分布收敛于 $Y$; 奇数 $n$ 时 $\Prob(|Y_n-Y|>1)=1$, 所以不依概率收敛于 $Y$.
\end{example}

\originalsection{弱大数定律及有限均值的边界}
\begin{theorem}[有限方差弱大数定律]
若 $X_k$ 独立同分布, $\E X_k=\mu$, $\Var(X_k)=\sigma^2<\infty$, 则
\[\Prob(|A_n-\mu|\geq\eps)\leq\frac{\sigma^2}{n\eps^2},\qquad A_n\xrightarrow{\Prob}\mu.\]
\end{theorem}
\begin{proof}
期望的线性性给出 $\E A_n=\mu$. 独立性使不同项的协方差为零, 故 $\Var(A_n)=n\sigma^2/n^2=\sigma^2/n$. 再应用Chebyshev不等式. 对任意固定的正 $\eps$, 上界趋于零.
\end{proof}
\begin{example}
独立检测中每次发现缺陷的概率为 $p$. 设 $X_k$ 是第 $k$ 次是否发现缺陷的指示变量. 由 $p(1-p)\leq1/4$, 要保证缺陷比例与 $p$ 的距离小于 $0.02$ 的概率至少 $0.95$, 充分条件是
\[\frac1{4n(0.02)^2}\leq0.05,\quad n\geq12500.\]
这是对所有 $p$ 都有效的保守保证, 不是精确所需的最小样本量.
\end{example}
有限方差并非弱大数定律的必要条件. 方差证明失效时, 可以先删去极端尾部, 再把删去部分的平均影响控制住.
\begin{theorem}[有限均值弱大数定律]
若 $X_k$ 独立同分布且 $\E|X_1|<\infty$, 则 $A_n\xrightarrow{\Prob}\E X_1$.
\end{theorem}
\begin{proof}
对固定 $K>0$, 令 $U_k=X_k\ind_{\{|X_k|\leq K\}}$, $V_k=X_k-U_k$, $m_K=\E U_1$. 有界的 $U_k$ 满足有限方差弱大数定律. 记 $r_K=\E(|X_1|\ind_{\{|X_1|>K\}})$, 可积性意味着 $r_K\to0$; 例如按 $j\leq|X_1|<j+1$ 分层, 可积的非负级数尾和趋于零. 又有 $|m_K-\mu|\leq r_K$.

给定 $\eps>0$, 取 $K$ 使 $r_K<\eps/3$. 三角不等式与Markov不等式给出
\[\Prob(|A_n-\mu|>\eps)\leq\Prob\left(\left|\frac1n\sum U_k-m_K\right|>\frac\eps3\right)+\frac{3r_K}{\eps}.\]
固定 $K$ 让 $n\to\infty$, 第一项趋于零; 再让 $K\to\infty$, 得到所需结论. 此处顺序不能倒置: 方差界的常数可以依赖于 $K$.
\end{proof}
\begin{example}
密度 $f(x)=\alpha x^{-\alpha-1}\ind_{\{x\geq1\}}$, $1<\alpha<2$, 满足
$\E X=\alpha/(\alpha-1)<\infty$, 而 $\E X^2=\infty$. 因而有限均值版本适用, 有限方差的概率上界不适用. 标准Cauchy变量的密度为 $1/[\pi(1+x^2)]$, 绝对均值无限; 其独立样本平均仍有同一Cauchy分布, 所以并不向某个常数集中. 后面的特征函数可直接核对这一点.
\end{example}

\originalsection{特征函数与中心极限定理}
大数定律说 $A_n-\mu$ 变小, 却没有说明放大后的偏差是什么形状. 独立和的标准差是 $\sigma\sqrt n$, 所以自然考察 $(S_n-n\mu)/(\sigma\sqrt n)$.
\begin{definition}
随机变量 $X$ 的特征函数是 $\varphi_X(t)=\E\ee^{\ii tX}$, $t\in\R$. 因 $|\ee^{\ii tX}|=1$, 它总是存在. 独立的 $X,Y$ 满足 $\varphi_{X+Y}(t)=\varphi_X(t)\varphi_Y(t)$; 而 $\varphi_{aX+b}(t)=\ee^{\ii tb}\varphi_X(at)$.
\end{definition}
两个公式分别由独立性分解期望和代入指数得到. 若 $\E|X|^m<\infty$, 在零点求导给出 $\varphi_X^{(m)}(0)=\ii^m\E X^m$, 因而 $\E X^m=\ii^{-m}\varphi_X^{(m)}(0)$. 求导交换可以先限制到 $|X|\leq K$, 再用 $|X|^m$ 的可积尾部控制差商. 这里的 $\ii^{-m}$ 不能误写为 $\ii^m$.

以下分析工具单独声明. 本章证明特征函数的极限计算, 但不把该工具的引用算作它的本地证明.
\begin{theorem}[L\'evy连续性定理, 引用的分析工具]
若概率分布的特征函数 $\varphi_n(t)$ 对每个实数 $t$ 趋于 $\varphi(t)$, 且 $\varphi$ 在零点连续, 则 $\varphi$ 是某个概率分布的特征函数, 对应分布依分布收敛到该分布. 特别地, 若极限已经是变量 $Y$ 的特征函数, 则 $Y_n\Rightarrow Y$.
\end{theorem}
这一定理涉及分布的紧性与Fourier变换的唯一性, 超出本书所设微积分先修范围, 本章不证明. 下面所有以它完成的证明都明确保留这一依赖; 不需要这个工具的有限均值弱大数定律已在上节独立证明.

为了看清讲义中的另一条弱大数定律路线, 若 $\E|X|<\infty$, 则
\[\varphi_X(h)=1+\ii h\mu+o(h).\]
确切地说, $(\ee^{\ii hX}-1)/h\to\ii X$, 且绝对值不超过 $|X|$; 在有界部分一致取极限, 尾部由 $\E(|X|\ind_{|X|>K})$ 控制, 就可交换极限. 因而
\[\varphi_{A_n}(t)=\varphi_X(t/n)^n\longrightarrow\ee^{\ii t\mu}.\]
这里使用 $n\log(1+z_n)=nz_n+O(n|z_n|^2)$, 其中 $z_n=\ii t\mu/n+o(1/n)$; 复对数只在1附近选取连续分支, 不需要全局定义. 由L\'evy定理和常数极限的等价性, 又得到弱大数定律. Cauchy特征函数 $\ee^{-|t|}$ 则给出 $\varphi_{A_n}(t)=\ee^{-|t|}$, 解释了平均不集中的现象; 此处使用Cauchy特征函数的已知计算.

\begin{lemma}[有限二阶矩的展开]
若 $\E Y=0$, $\E Y^2=1$, 则 $\varphi_Y(h)=1-h^2/2+o(h^2)$.
\end{lemma}
\begin{proof}
令 $r(u)=\ee^{\ii u}-1-\ii u+u^2/2$. 微积分的Taylor公式给出 $r(u)/u^2\to0$. 同时
$|\ee^{\ii u}-1-\ii u|\leq u^2/2$, 可由二次积分余项得到, 所以 $|r(u)|\leq u^2$. 对 $|Y|\leq K$ 的部分, $r(hY)/h^2$ 一致趋于零; 剩余部分的期望绝对值不超过 $\E(Y^2\ind_{|Y|>K})$, 它随 $K\to\infty$ 趋于零. 故 $\E r(hY)=o(h^2)$, 取期望并使用两个矩条件即得结论.
\end{proof}
\begin{theorem}[独立同分布中心极限定理]
若 $X_k$ 独立同分布, $\E X_k=\mu$, $0<\Var(X_k)=\sigma^2<\infty$, 则
\[B_n=\frac{S_n-n\mu}{\sigma\sqrt n}\Rightarrow Z,\qquad Z\sim\Normal(0,1).\]
因此对固定实数 $a<b$, $\Prob(a\leq B_n\leq b)\to\Phi(b)-\Phi(a)$, 其中 $\Phi(x)=\int_{-\infty}^x\ee^{-u^2/2}\dd u/\sqrt{2\pi}$.
\end{theorem}
\begin{proof}[证明(调用L\'evy连续性定理)]
令 $Y_k=(X_k-\mu)/\sigma$. 它们独立同分布, 均值零, 二阶矩一. 引理与独立性给出
\[\varphi_{B_n}(t)=\left[\varphi_Y\left(\frac{t}{\sqrt n}\right)\right]^n
=\left[1-\frac{t^2}{2n}+o\left(\frac1n\right)\right]^n\to\ee^{-t^2/2}.\]
最后一步仍可用1附近的复对数: 对固定 $t$, $n\log\varphi_Y(t/\sqrt n)\to-t^2/2$. 标准正态的特征函数为 $\ee^{-t^2/2}$: 对高斯密度积分求导并分部积分得到 $\varphi'_Z(t)=-t\varphi_Z(t)$, 而 $\varphi_Z(0)=1$, 解此常微分方程即得公式. 高斯尾部保证求导、分部积分均合法. 现在调用L\'evy定理得到依分布收敛. 因正态分布连续, 两个端点的分布函数收敛给出区间概率; 若使用闭区间, 可用任意小的端点扩张和缩小夹逼, 再让扩张量趋于零.
\end{proof}
对Bernoulli变量 $X_k$ 取 $\Prob(X_k=1)=p$, $0<p<1$, 则 $\mu=p$, $\sigma^2=p(1-p)$, 所以CLT立即给出De Moivre--Laplace结论:
\[\Prob\left(a\leq\frac{S_n-np}{\sqrt{np(1-p)}}\leq b\right)\longrightarrow\Phi(b)-\Phi(a).\]
它把二项计数的正态近似置于一般独立和的理论之中.

若 $\sigma=0$, 每个 $X_k=\mu$ 几乎处处, 和没有波动, 无须也不能除以 $\sigma$. CLT不声称 $B_n$ 依概率趋于一个正态变量, 也不自动提供给定 $n$ 的误差界.

讲义还给出矩母函数路线. 若存在 $\delta>0$ 使 $M_Y(t)=\E\ee^{tY}$ 在 $|t|<\delta$ 有限, 可用两侧指数可积性控制求导, 得 $M_Y(t)=1+t^2/2+o(t^2)$. 因而 $M_{B_n}(t)\to\ee^{t^2/2}$ 对每个固定 $t$ 成立(当 $n$ 足够大时函数存在). 矩母函数连续性定理的适用版本是: 在零点某个共同开邻域中, 矩母函数趋于某分布的有限矩母函数, 则分布收敛. 此定理同样是外部分析工具. 有限方差本身不保证矩母函数在零点邻域存在, 所以一般CLT必须使用特征函数路线.

\begin{example}
每个随机计分 $X$ 以概率 $1/4,1/2,1/4$ 取 $0,2,4$. 有 $\E X=2$, $\E X^2=6$, $\Var(X)=2$. 对 $n=5000$ 次独立计分, $\E S_n=10000$, 标准差为100. 因而
\[\Prob(9800\leq S_n\leq10100)\approx\Phi(1)-\Phi(-2)\approx0.8186.\]
这是近似值. 若希望改善整数端点附近的近似, 可以使用格点跨度为2所对应的半格修正, 但CLT本身不证明该修正的误差更小.
\end{example}
\begin{example}
某计数模型每小时产生独立的 $\Pois(3)$ 次事件. $m$ 小时总数 $N_m\sim\Pois(3m)$, 均值和方差均为 $3m$. 把它写成 $m$ 个独立小时计数的和, 得
$\Prob(N_m\leq3m+a\sqrt{3m})\approx\Phi(a)$. 例如 $m=300$ 时标准差为30, 超过960次的概率约为 $1-\Phi(2)\approx0.0228$. 近似适用于所声明的独立、恒定速率模型, 不能仅凭事件数量大就断定现实数据服从正态.
\end{example}
非同分布或依赖变量也有CLT, 但需要额外条件以排除某一项支配总波动等情形. 本章不陈述这些推广, 因此“很多小因素相加”只是寻找模型的线索, 不是可以代替条件的证明.

\originalsection{强大数定律与乘法过程}
弱大数定律分别控制每个大的 $n$; 强大数定律要求在概率为一的同一批路径上, 此后所有平均值最终都接近 $\mu$. 仅有 $O(1/n)$ 的偏差概率不可求和, 因而不能直接排除无限次偏差. 第31讲通过四阶矩把概率界改成可求和的 $O(1/n^2)$.
\begin{lemma}[第一Borel--Cantelli引理]
若 $\sum_n\Prob(E_n)<\infty$, 则 $E_n$ 无限次发生的概率为零, 无须独立性.
\end{lemma}
\begin{proof}
对任意 $N$, $\Prob(\bigcup_{n\geq N}E_n)\leq\sum_{n\geq N}\Prob(E_n)\to0$. 无限次发生的事件包含在每个这样的并集中, 所以其概率为零.
\end{proof}
\begin{theorem}[本章证明的强大数定律版本]
若 $X_k$ 独立同分布且 $\E|X_1|^4<\infty$, 则 $A_n\to\mu=\E X_1$ 几乎处处.
\end{theorem}
\begin{proof}
中心化 $Y_k=X_k-\mu$. 由 $|x-y|^4\leq8(|x|^4+|y|^4)$, $K=\E Y_1^4<\infty$. 记 $v=\E Y_1^2$, 则 $v^2\leq K$, 因 $\Var(Y_1^2)=K-v^2\geq0$.

展开 $T_n^4=(\sum Y_k)^4$. 指标出现方式为 $1+1+1+1$, $2+1+1$, $3+1$, $2+2$, $4$. 前三种都有一个只出现一次的独立因子, 均值零, 所以期望为零. 每对不同指标产生 $6Y_i^2Y_j^2$, 因为从四个位置中选两个放 $i$ 有 $\binom42=6$ 种. 因而
\[\E T_n^4=nK+6\binom n2v^2=nK+3n(n-1)v^2\leq3Kn^2.\]
Markov不等式给出
$\Prob(|T_n/n|>\eps)\leq3K/(n^2\eps^4)$, 其和有限. Borel--Cantelli表明, 对固定 $\eps$, 偏差超过 $\eps$ 只发生有限次. 对 $\eps=1,1/2,1/3,\ldots$ 取可数个概率为一的事件之交, 概率仍为一. 在交集上对每个正 $\eps$ 选 $1/j<\eps$, 从而 $T_n/n\to0$.
\end{proof}
\begin{remark}
一般独立同分布强大数定律只需 $\E|X_1|<\infty$. 本章并未证明这个最一般版本; 上面的完整本地证明范围是有限四阶矩. 这足以覆盖本节所有有界乘数的例子, 也与源讲义的证明条件一致. 强大数定律蕴含弱大数定律, 正是本章第一节已证明的几乎处处收敛蕴含依概率收敛.
\end{remark}
\begin{example}
设某抽象乘法过程每期乘数 $R_k$ 独立, 以概率 $0.54$ 取 $1.2$, 以概率 $0.46$ 取 $0.8$. 总量 $W_n=W_0\prod_{k=1}^nR_k$, $W_0>0$. 一方面
\[\E R_1=1.016,\qquad \E W_n=W_0(1.016)^n\to\infty.\]
另一方面, 令 $X_k=\log R_k$, 它有界, 而
\[\mu_{\log}=0.54\log1.2+0.46\log0.8\approx-0.00419<0.\]
强大数定律给出 $n^{-1}\log(W_n/W_0)\to\mu_{\log}$ 几乎处处. 在这些路径上最终 $\log(W_n/W_0)<n\mu_{\log}/2\to-\infty$, 因此 $W_n\to0$ 几乎处处. 期望增长与路径衰减并不矛盾: 少数极大值可以支撑期望. 这里没有交换极限和期望的可积控制.
\end{example}

\originalsection{Jensen不等式及等号}
\begin{definition}
区间 $I$ 上的实值函数 $g$ 称为凸函数, 若对 $x,y\in I$, $0\leq\lambda\leq1$,
$g(\lambda x+(1-\lambda)y)\leq\lambda g(x)+(1-\lambda)g(y)$. 若 $x\ne y$, $0<\lambda<1$ 时不等号严格, 则称严格凸. $g$ 凹是指 $-g$ 凸.
\end{definition}
图形上, 凸函数位于两点连线之下. 若二阶可微, $g''\geq0$ 等价于凸性: 导数递增使割线斜率递增, 从而满足凸性; 反过来凸性使左右差商夹住导数并使导数递增, 再求导即得 $g''\geq0$.
\begin{lemma}[支撑直线]
若 $g$ 在区间 $I$ 凸, $m$ 是 $I$ 的内点, 则存在有限实数 $a$, 使对所有 $x\in I$,
$g(x)\geq g(m)+a(x-m)$.
\end{lemma}
\begin{proof}
对 $x<m<y$, 凸性给出
\[\frac{g(m)-g(x)}{m-x}\leq\frac{g(y)-g(m)}{y-m}.\]
左侧所有斜率的上确界不超过右侧所有斜率的下确界. 两者有限: 取固定的左右内点, 用它们的割线斜率给这些界限以有限约束. 在两者之间取 $a$. 对 $x<m$, 左侧斜率不超过 $a$, 移项得所需式; 对 $y>m$, 右侧斜率不小于 $a$, 也得所需式. 在 $m$ 取等. 若可微, 可以取 $a=g'(m)$; 不可微时仍有支撑线.
\end{proof}
\begin{theorem}[Jensen不等式]
设 $X\in I$ 几乎处处, $\E|X|<\infty$, $m=\E X$ 是区间 $I$ 的内点, $g:I\to\R$ 凸. 则 $\E g(X)$ 在扩展实数意义下存在, 且
\[\E g(X)\geq g(\E X).\]
取任一在 $m$ 的支撑斜率 $a$, 当期望有限时, 等号成立当且仅当
$g(X)=g(m)+a(X-m)$ 几乎处处. 特别地, 若 $g$ 严格凸, 等号成立当且仅当 $X=m$ 几乎处处. 凹函数在其期望有意义时不等号反向.
\end{theorem}
\begin{proof}
支撑线 $L(x)=g(m)+a(x-m)$ 满足 $g(X)-L(X)\geq0$. $L(X)$ 可积, 所以 $g(X)$ 的负部可积, 期望不会出现 $\infty-\infty$. 取期望得
$\E g(X)\geq\E L(X)=g(m)$. 等号等价于非负变量 $g(X)-L(X)$ 的期望为零, 而非负变量期望为零等价于它几乎处处为零: 若 $\Prob(H>0)>0$, 某个 $j$ 满足 $\Prob(H\geq1/j)>0$, 则 $\E H>0$.

若严格凸且某 $x\ne m$ 也落在支撑线上, 对 $m,x$ 间的内点, 严格凸性要求函数低于这条线, 与支撑性质矛盾. 因而接触集合只有 $m$, 等号条件随之成立. 对 $-g$ 应用结果得到凹函数版本.
\end{proof}
若均值是 $I$ 中包含的端点, 则 $X$ 被限制在端点一侧; $\E(X-m)=0$ 和固定符号使 $X=m$ 几乎处处, Jensen直接取等. 这处理了上述内点条件之外的边界, 但不能给未定义的端点值强行赋值.

\begin{example}
对正的可积变量 $R$, 若 $\E|\log R|<\infty$, 严格凹的 $\log$ 给出
$\E\log R\leq\log\E R$, 非常数时严格. 这解释乘法过程为何必须看平均对数增长率. 对同样均值的两个结果, Jensen只保证确定值 $\E R$ 的对数效用不低于随机结果; 它并不比较任意两个不同随机分布的风险高低.
\end{example}
\begin{example}
设报酬函数 $g(x)=0.01C+0.25\max(x-C,0)$, $C>0$. 它是两条仿射函数的最大值, 因而凸. 确定结果 $X=C$ 的报酬为 $0.01C$. 若 $X$ 以相等概率取 $0.6C,1.4C$, 均值仍为 $C$, 但平均报酬为
$[0.01C+(0.01C+0.10C)]/2=0.06C$.

这个模型展示凸报酬会使执行者从波动获益. 如果每期另有独立的失败概率 $q>0$, 第一次失败前的期数具有几何尾 $\Prob(T>n)=(1-q)^n\to0$, 所以最终失败几乎必然. 执行者此前所得报酬与委托者最终损失是不同随机量. 此处讨论模型的激励结构, 不由Jensen推断所有更波动的方案都更好, 也不对实际投资提出建议.
\end{example}

\originalsection{习题与解答}
\begin{exercise}
独立变量 $X_k$ 均值为 $\mu$, 方差不超过常数 $C$, 不要求同分布. 证明 $n^{-1}\sum_{k=1}^nX_k\xrightarrow{\Prob}\mu$.
\end{exercise}
\begin{solution}
平均值期望仍为 $\mu$, 独立性给出方差不超过 $nC/n^2=C/n$. Chebyshev给出偏差超过 $\eps$ 的概率不超过 $C/(n\eps^2)\to0$. 这说明相同均值与统一方差界已足以完成这一弱定律版本.
\end{solution}
\begin{exercise}
令 $X_k$ 独立同分布, 密度 $\ee^{-x}\ind_{x>0}$. 给出 $S_n$ 的CLT标准化, 并估计 $n=400$ 时 $\Prob(S_n>440)$. 说明 $S_n/n$ 与标准化变量的极限为何不同.
\end{exercise}
\begin{solution}
分部积分得到 $\E X=1$, $\E X^2=2$, 方差1. 所以 $(S_n-n)/\sqrt n\Rightarrow\Normal(0,1)$. 在 $n=400$ 时门槛标准化为 $(440-400)/20=2$, 所求近似为 $1-\Phi(2)\approx0.0228$. 平均值的偏差尺度为 $1/\sqrt n$ 并趋于零, 所以 $S_n/n\xrightarrow{\Prob}1$; CLT把这个变小的偏差乘以 $\sqrt n$, 保留非退化的形状.
\end{solution}
\begin{exercise}
若 $X$ 在 $[-2,2]$ 取值且 $\E X=0$, 证明 $\E\ee^X\geq1$, 并确定等号. 再说明 $\E|X|\geq|\E X|$ 的等号是否必须要求 $X$ 为常数.
\end{exercise}
\begin{solution}
指数严格凸且 $X$ 有界, Jensen给出 $\E\ee^X\geq\ee^{\E X}=1$, 等号当且仅当 $X=0$ 几乎处处. 绝对值只是凸, 在非负半轴仿射. 例如以相等概率取1和3, 有 $\E|X|=|\E X|=2$ 却非常数. 一般该等号成立当且仅当 $X$ 几乎处处同号(允许零): 写 $X=X^+-X^-$, 等号等价于 $\E X^+=0$ 或 $\E X^-=0$.
\end{solution}

\originalchapter{有限状态马尔可夫链}
\sourcelecture{32}

\originalsection{状态、历史与转移矩阵}
独立试验每次重新开始, 下一次的分布不受过去影响. 另一类系统有持续的状态: 机器是否运行、队列所处级别、天气类别. 马尔可夫模型允许下一步依赖现在, 同时假设现在的状态已包含预测下一步所需的全部过去信息.
\begin{definition}
设状态空间 $S=\{0,\ldots,d-1\}$ 有限. 序列 $X_0,X_1,\ldots$ 是时间齐次马尔可夫链, 若存在矩阵 $P=(p_{ij})$, 使每个概率为正的历史满足
\[\Prob(X_{n+1}=j\mid X_0=i_0,\ldots,X_n=i)=p_{ij}.\]
矩阵的每个元素非负, 每行和为1, 称为转移矩阵. 初始分布记为行向量 $\alpha=(\alpha_0,\ldots,\alpha_{d-1})$.
\end{definition}
齐次是指 $p_{ij}$ 不依赖时刻 $n$. 条件概率只在条件事件概率为正时使用; 不必给不可能的历史赋予额外意义. 给定初始分布和转移矩阵, 一条有限路径的概率由乘法公式确定:
\[\Prob(X_0=i_0,\ldots,X_n=i_n)=\alpha_{i_0}\prod_{r=0}^{n-1}p_{i_ri_{r+1}}.\]
这是把抽象定义转换为可计算模型的第一步.

\begin{example}
某设备每天记录为正常状态0或维护状态1. 正常后下一天进入维护的概率为 $1/4$, 维护后下一天恢复正常的概率为 $1/2$. 则
\[P=\begin{pmatrix}3/4&1/4\\1/2&1/2\end{pmatrix}.\]
从正常开始, 连续两天仍正常的路径概率为 $(3/4)^2=9/16$. 两天后正常的概率还要加上先维护再恢复的路径, 得 $9/16+(1/4)(1/2)=11/16$. 单一路径概率与末端事件概率不是同一个量.
\end{example}
\begin{center}
\begin{tikzpicture}[>=Stealth,every node/.style={font=\small}]
\node[draw,circle] (a) at (0,0) {正常};
\node[draw,circle] (b) at (4,0) {维护};
\draw[->] (a) edge[bend left=20] node[above] {$1/4$} (b);
\draw[->] (b) edge[bend left=20] node[below] {$1/2$} (a);
\draw[->] (a) edge[loop left] node[left] {$3/4$} (a);
\draw[->] (b) edge[loop right] node[right] {$1/2$} (b);
\end{tikzpicture}
\end{center}

\begin{proposition}[Chapman--Kolmogorov公式]
写 $p_{ij}^{(n)}=\Prob(X_n=j\mid X_0=i)$, 其中从 $i$ 出发的链按相同转移规则定义. 则 $P^{(0)}=I$, $P^{(n)}=P^n$, 且
\[p_{ij}^{(m+n)}=\sum_{k\in S}p_{ik}^{(m)}p_{kj}^{(n)}.\]
时刻 $n$ 的分布是 $\alpha P^n$.
\end{proposition}
\begin{proof}
按中间状态 $X_m=k$ 分解事件. 路径概率的乘法式说明, 给定 $X_m=k$ 和任何此前历史, 之后 $n$ 步到 $j$ 的概率只依赖 $k$, 等于 $p_{kj}^{(n)}$. 全概率公式给出所列求和. 取 $n=1$ 归纳得到矩阵幂. 再按初始状态求和, 第 $j$ 个概率为 $\sum_i\alpha_i(P^n)_{ij}$.
\end{proof}
例如设备模型有
\[P^2=\begin{pmatrix}11/16&5/16\\5/8&3/8\end{pmatrix}.\]
行表示起点, 列表示终点. 采用行分布向量时, 必须在右侧乘 $P$, 不能把行列约定混用.

三种特殊矩阵帮助理解定义. 若所有行等于同一分布 $q$, 则 $X_1,X_2,\ldots$ 独立同分布为 $q$, 并且与 $X_0$ 独立; 只有初始分布也为 $q$ 时, 整个 $X_0,X_1,\ldots$ 才独立同分布. 若 $P=I$, 状态永不改变, 却可以有随机的起点. 若每个元素是0或1, 每行恰有一个1, 以函数 $f(i)$ 表示其所在列, 则 $X_{n+1}=f(X_n)$, 初始状态确定之后路径完全确定.

\originalsection{停留时间与状态选择}
\begin{proposition}
从状态 $i$ 出发, 令 $T=\min\{n\geq1:X_n\ne i\}$. 若 $p_{ii}<1$, 则
\[\Prob(T=k)=p_{ii}^{k-1}(1-p_{ii}),\quad k\geq1,\qquad \E T=\frac1{1-p_{ii}}.\]
若 $p_{ii}=1$, 则 $T=\infty$ 几乎处处.
\end{proposition}
\begin{proof}
$T=k$ 要求前 $k-1$ 步都留在 $i$, 第 $k$ 步离开, 路径乘法给出概率. 尾和公式给出 $\E T=\sum_{n\geq0}\Prob(T>n)=\sum_{n\geq0}p_{ii}^n$. 最后一种情况每一步离开概率为零, 可数并集仍概率零.
\end{proof}
在设备模型中, 从正常到首次维护平均等待4天, 从维护到首次正常平均等待2天. $T$ 从下一天计起, 因而不把今天额外加一次.

几何停留意味着已经在状态内待了很久, 也不会改变下一步离开的概率. 这是模型的可检验限制. 若一台机器的故障概率随累计运行时间上升, 仅用“正常/故障”两个状态往往不能描述历史影响. 可以把正常状态再细分为运行年龄级别, 或记录必要的累计变量. 扩充状态能保留所需记忆, 但任意扩充不自动保证马尔可夫性, 仍要检验条件概率是否真的只依赖新状态. 源讲义用关系状态讨论同样的建模问题; 此处以设备停留时间训练相同的判断.

\originalsection{平稳分布与长期分布}
\begin{definition}
概率行向量 $\pi$ 若满足 $\pi P=\pi$, 称为平稳分布. 若以 $\pi$ 初始化, 每个时刻的分布都为 $\pi$. 若存在整数 $r\geq1$ 使 $P^r$ 的所有元素严格为正, 称 $P$ 为本原转移矩阵. 源讲义把这个性质称作遍历性.
\end{definition}
平稳的意思是整个分布不变, 不是每条路径停在某状态. 它是特征值1的左特征向量, 并须同时满足非负与归一化条件. 若用列向量解方程, 正确方程是 $(P^{\mathsf T}-I)\pi^{\mathsf T}=0$.
\begin{example}
设备模型的平稳方程为
$\pi_0=(3/4)\pi_0+(1/2)\pi_1$, $\pi_0+\pi_1=1$.
第一式给 $\pi_0=2\pi_1$, 所以 $\pi=(2/3,1/3)$. 令 $u_n=\Prob(X_n=0)$, 全概率公式给出
\[u_{n+1}=\tfrac34u_n+\tfrac12(1-u_n)=\tfrac12+\tfrac14u_n.\]
减去固定点 $2/3$, 得 $u_{n+1}-2/3=(u_n-2/3)/4$, 因而
\[u_n=\frac23+4^{-n}\left(u_0-\frac23\right).\]
这个计算同时给出平稳分布和任意初始分布的收敛速率, 比只解线性方程多了收敛结论.
\end{example}

\begin{theorem}[有限链的平稳分布存在]
每个有限状态转移矩阵至少有一个平稳分布.
\end{theorem}
\begin{proof}
取任意概率行向量 $\alpha$, 令 $\nu_n=n^{-1}\sum_{k=0}^{n-1}\alpha P^k$. 它们的每个坐标都在 $[0,1]$, 由有限维Bolzano--Weierstrass定理可取收敛子列 $\nu_{n_l}\to\pi$. 极限仍非负且坐标和为1. 又有
$\nu_nP-\nu_n=(\alpha P^n-\alpha)/n\to0$.
矩阵乘法是连续的, 沿子列取极限得 $\pi P=\pi$.
\end{proof}
\begin{lemma}[正矩阵幂带来的收缩]
设 $Q=P^r$ 所有元素严格为正, $\eta=\min_{i,j}Q_{ij}>0$, $\theta=d\eta\leq1$. 对任意两个概率行向量 $\alpha,\beta$,
\[\|\alpha Q-\beta Q\|_1\leq(1-\theta)\|\alpha-\beta\|_1,\qquad \|v\|_1=\sum_i|v_i|.\]
\end{lemma}
\begin{proof}
若 $\theta=1$, 每行和为1迫使 $Q_{ij}=1/d$, 两个分布乘 $Q$ 后都成为均匀分布. 若 $\theta<1$, 令 $U$ 每个元素为 $1/d$, $R_{ij}=(Q_{ij}-\eta)/(1-\theta)$. $R$ 非负, 行和1, 且 $Q=\theta U+(1-\theta)R$. 因概率向量之差的坐标和为零, $(\alpha-\beta)U=0$. 对任何行向量 $v$,
\[\|vR\|_1\leq\sum_j\sum_i|v_i|R_{ij}=\sum_i|v_i|.\]
代入 $v=\alpha-\beta$ 即得结论. 同一计算也说明乘任意转移矩阵不会增大此距离.
\end{proof}
\begin{theorem}[本原有限链的收敛]
若 $P$ 本原, 则平稳分布唯一, 且对任意初始分布 $\alpha$,
\[\alpha P^n\to\pi,\qquad (P^n)_{ij}\to\pi_j.\]
若 $P^r$ 的最小元素为 $\eta$, 则
$\|\alpha P^n-\pi\|_1\leq2(1-d\eta)^{\lfloor n/r\rfloor}$.
\end{theorem}
\begin{proof}
先由存在定理取平稳分布 $\pi$. 若 $\rho$ 也是平稳分布, 对 $Q^k=P^{rk}$ 反复收缩得到
$\|\rho-\pi\|_1\leq(1-d\eta)^k\|\rho-\pi\|_1$, 令 $k\to\infty$ 得唯一性. 写 $n=kr+s$, $0\leq s<r$. 先乘 $Q^k$ 收缩 $k$ 次, 再乘 $P^s$ 不增加距离, 得所列界, 初始距离至多2. 取 $\alpha$ 为状态 $i$ 的点质量就得到逐行收敛.
\end{proof}

\originalsection{不可约、周期与收敛条件}
\begin{definition}
若对每对 $i,j$ 存在 $n\geq0$ 使 $(P^n)_{ij}>0$, 称链不可约. 状态 $i$ 的周期是
$\gcd\{n\geq1:(P^n)_{ii}>0\}$. 周期为1称非周期. 在不可约链中这些集合非空, 且所有状态周期相同.
\end{definition}
最后一条也需要理由. 若有 $i$ 到 $j$ 的正概率路径长 $a$, $j$ 到 $i$ 的路径长 $b$, 则 $a+b$ 是 $i$ 的返回长度; 对任一 $j$ 的返回长度 $n$, $a+n+b$ 也是 $i$ 的返回长度. 所以 $i$ 的周期同时整除 $a+b$ 与 $a+n+b$, 从而整除所有 $n$, 即整除 $j$ 的周期. 交换 $i,j$ 得二者相等.

\begin{lemma}[返回长度的整数工具]
若一组正整数的最大公因数为1, 则其中可取有限个 $a_1,\ldots,a_l$ 最大公因数为1, 且每个足够大的整数是这些数的非负整数线性组合.
\end{lemma}
\begin{proof}
逐步加入整数所得最大公因数是递减的正整数序列, 只会有限次严格下降, 最终为1, 所以可取有限组. 模 $a_1$ 看 $a_2,\ldots,a_l$ 生成的余数. 非负组合的余数集合对加法封闭, 且包含每个元素的加法逆元: 在有限循环群中, $-x=(a_1-1)x$. 因而它是子群. 最大公因数为1和B\'ezout等式说明该子群含余数1, 即含全部余数. 为每个余数 $s$ 选一个非负组合 $b_s$ 代表它. 若 $n\geq\max_s b_s$, 取其余数 $s$, 则 $n=b_s+ka_1$, $k\geq0$, 完成证明.
\end{proof}
\begin{theorem}
对有限状态链, 本原等价于不可约且非周期. 因此不可约非周期的有限链满足上一节的唯一平稳分布与分布收敛结论.
\end{theorem}
\begin{proof}
若 $P^r>0$, 则所有状态互通. 而 $P^{r+1}>0$: 对任意列 $j$, 不可约性保证存在某 $k$ 使 $p_{kj}>0$, 故 $(P^{r+1})_{ij}\geq(P^r)_{ik}p_{kj}>0$. 每个状态可在 $r$ 与 $r+1$ 步返回, 周期为1.

反过来固定状态 $o$. 由非周期性, $o$ 的返回长度最大公因数为1. 引理选出的有限返回长度都有正概率路径; 拼接这些路径说明 $o$ 可以在每个充分大的步数返回. 对每个 $i,j$, 选一条 $i$ 到 $o$ 的路径长 $a_i$ 和一条 $o$ 到 $j$ 的路径长 $b_j$, 允许 $a_o=b_o=0$. 当 $n-a_i-b_j$ 足够大时, 中间插入一次该长度的返回路径, 就得 $(P^n)_{ij}>0$. 状态对有限, 可选一个共同的充分大 $n$, 得 $P^n>0$.
\end{proof}
\begin{example}
$P=\begin{pmatrix}0&1\\1&0\end{pmatrix}$ 不可约, 周期2. 平稳分布为 $(1/2,1/2)$, 但从0出发, 偶数时刻一定在0, 奇数时刻一定在1, 分布不收敛. 相反 $P=I$ 的每个概率向量都平稳, 起点决定最终状态. 前例显示平稳存在甚至唯一不足以保证分布收敛, 后例显示不可约条件在唯一性中发挥作用.
\end{example}

\originalsection{长期时间比例与奖励平均}
平稳分布的坐标 $\pi_j$ 还可以表示单条长期路径中访问 $j$ 的比例. 这是另一个定理, 不能从 $\Prob(X_n=j)\to\pi_j$ 直接推出, 因为各时刻的指示变量通常依赖.
\begin{theorem}[有限不可约链的时间平均]
若有限链不可约, 对任意初始分布和任意函数 $f:S\to\R$, 几乎处处有
\[\frac1n\sum_{k=0}^{n-1}f(X_k)\longrightarrow\sum_{j\in S}\pi_jf(j),\]
其中 $\pi$ 是唯一平稳分布. 此结论不需要非周期条件. 特别地, 取 $f=\ind_{\{j\}}$, 得访问状态 $j$ 的长期比例为 $\pi_j$.
\end{theorem}
\begin{proof}
固定状态 $o$. 不可约性和状态有限意味着存在 $L\geq1$, $c>0$, 从任意状态在随后 $L$ 步内访问 $o$ 的概率至少 $c$: 为每个状态选正概率路径, 取共同长度上界与路径概率的正下界. 在 $o$ 出发时要考虑正时间返回, 可以先走一步, 再使用此界, 把 $L$ 增大1即可. 每个长度 $L$ 的区块, 给定此前历史后仍有至少 $c$ 的成功概率, 所以首次访问或返回时间 $H$ 满足
$\Prob(H>mL)\leq(1-c)^m$. 它几乎处处有限, 并且所有正整数阶矩有限; 例如按区块分层, $\E H^q\leq\sum_{m\geq0}((m+1)L)^q(1-c)^m<\infty$.

先从 $o$ 出发. 令 $\tau_0=0$, $\tau_{l+1}=\min\{n>\tau_l:X_n=o\}$, 周期长度 $D_l=\tau_{l+1}-\tau_l$, 周期奖励 $V_l=\sum_{k=\tau_l}^{\tau_{l+1}-1}f(X_k)$. 这些对 $(D_l,V_l)$ 独立同分布. 这里用到的停时马尔可夫性质可直接由路径概率证明: 将 $\{\tau_l=t\}$ 按所有在 $t$ 结束的历史分解, 对每条历史, 之后任意有限路径的条件概率均为从 $o$ 开始的同一乘积, 不依赖历史; 再对 $t$ 求和. 因访问时间有限, 这些历史的概率和为1. 递归应用即得各周期独立同分布.

有限状态上 $M=\max|f|<\infty$, $|V_l|\leq MD_l$, 因而二者有有限四阶矩. 上章已证明的强大数定律给出
\[\frac{\tau_m}{m}=\frac1m\sum_{l=0}^{m-1}D_l\to d_0=\E D_0>0,\qquad
\frac1m\sum_{l=0}^{m-1}V_l\to v_0=\E V_0.\]
令 $m(n)$ 是时刻 $n$ 前已完成的周期数, 即 $\tau_{m(n)}\leq n<\tau_{m(n)+1}$. 由于 $\tau_m/m\to d_0$, 夹逼给 $m(n)/n\to1/d_0$. 未完成周期不能破坏平均: 对每个 $\eps>0$, 指数尾使 $\sum_l\Prob(D_l>\eps l)<\infty$, Borel--Cantelli给 $D_l/l\to0$ 几乎处处. 故余下奖励绝对值至多 $MD_{m(n)}=o(n)$. 时间平均于是趋于 $v_0/d_0$.

为识别这个比值, 令 $U_j$ 是一个周期中处于 $j$ 的次数, $u_j=\E U_j$. 有 $\sum_j u_j=d_0$. 对周期内从 $i$ 到 $j$ 的转移次数取期望, 得 $u_i p_{ij}$: 写成对 $k\geq0$ 的指示和, 事件 $\{k<\tau_1,X_k=i\}$ 由时刻 $k$ 的历史决定, 条件期望中的下一步概率为 $p_{ij}$. 非负求和允许逐项取期望. 一个周期的终点与起点同为 $o$, 所以每个状态的进入次数等于离开次数, 因而 $u_j=\sum_i u_i p_{ij}$. 所以 $\pi_j=u_j/d_0$ 是平稳分布, 且 $v_0/d_0=\sum_j\pi_jf(j)$.

若初始状态不是 $o$, 到首次访问 $o$ 之前的段落长度几乎处处有限, 奖励也有限, 除以 $n$ 后趋于零. 从该访问时刻开始重复以上周期证明, 得同样结论. 最后, 若 $\rho$ 是任一平稳分布, 从 $\rho$ 初始化后, 每个时间平均的期望是 $\sum_j\rho_jf(j)$. 时间平均有界于 $M$, 其几乎处处极限的期望可交换: 有界变量的收敛可由偏差事件概率趋零与 $2M$ 的上界直接证明. 因而 $\sum_j\rho_jf(j)=\sum_j\pi_jf(j)$. 逐个取指示函数得 $\rho=\pi$, 也证明了不可约链平稳分布唯一.
\end{proof}
\begin{example}
设备模型若正常日奖励为3, 维护日奖励为 $-1$, 则长期每天平均奖励几乎处处趋于
$(2/3)3+(1/3)(-1)=5/3$. 正常日与维护日并不独立, 所以不能直接把独立同分布大数定律用于每日奖励. 上述周期分解正是恢复独立结构的方法. 对周期2的交替链, 时间比例仍为 $1/2,1/2$, 虽然单个时刻的分布始终振荡; 因而时间平均与分布收敛必须分开讨论.
\end{example}

\originalsection{习题与解答}
\begin{exercise}
两状态链的矩阵为 $P=\begin{pmatrix}1-a&a\\b&1-b\end{pmatrix}$, $0<a,b<1$. 求平稳分布与任意初始分布的 $n$ 步概率. 求每次处于0后首次到1的平均等待时间.
\end{exercise}
\begin{solution}
平稳方程给 $a\pi_0=b\pi_1$, 故 $\pi=(b/(a+b),a/(a+b))$. 设 $u_n=\Prob(X_n=0)$, 则 $u_{n+1}=b+(1-a-b)u_n$. 减去固定点得
\[u_n=\frac b{a+b}+(1-a-b)^n\left(u_0-\frac b{a+b}\right).\]
因为 $|1-a-b|<1$, 分布趋稳, 若该因子负则误差交替变号. 从0离开的等待时间为成功概率 $a$ 的几何变量, 均值 $1/a$.
\end{solution}
\begin{exercise}
三状态链有 $P=\begin{pmatrix}0&1&0\\0&0&1\\1/2&0&1/2\end{pmatrix}$. 判断不可约与周期, 求平稳分布, 求状态2的长期访问比例.
\end{exercise}
\begin{solution}
路径 $0\to1\to2\to0$ 使所有状态互通, 所以不可约. 状态2有自环, 返回长度集合含1, 周期为1, 其他状态周期也为1. 平稳方程为 $\pi_0=\pi_2/2$, $\pi_1=\pi_0$, $\pi_2=\pi_1+\pi_2/2$. 归一化给 $\pi=(1/4,1/4,1/2)$. 不可约非周期有限链分布收敛到它; 独立于该分布收敛陈述, 时间平均定理保证单条路径访问2的比例几乎处处趋于 $1/2$.
\end{solution}
\begin{exercise}
设两状态转移矩阵为单位矩阵, 初始状态各以概率 $1/2$ 选取. 每个时刻的分布都是 $(1/2,1/2)$. 判断路径中状态0的访问比例是否趋于 $1/2$.
\end{exercise}
\begin{solution}
路径始终等于初始状态. 状态0的访问比例对所有 $n$ 都是 $\ind_{\{X_0=0\}}$, 因而以相等概率为1或0, 永不等于 $1/2$. 此链可约. 平稳初始化可以保证边缘分布不变, 却不能替代时间平均定理所需的不可约性.
\end{solution}

\originalchapter{熵、条件信息与无噪声编码}
\sourcelecture{33}

连续投掷 $k$ 枚公平硬币有 $2^k$ 个等可能结果, 用 $k$ 个二进制位就能把结果写清楚. 如果结果不等可能, 经常出现的结果应使用短码, 少见结果可以使用长码. 本章把这一直觉变成两个可以证明的结论: 平均码长不能低于熵, 而适当编码可以使它与熵相差不到一个比特. 条件熵则说明已有信息怎样改变剩余的不确定性.

\originalsection{信息量与熵}
本章 $\log_2$ 始终表示以 $2$ 为底的对数, $\ln$ 表示自然对数. 概率为 $p>0$ 的结果的信息量定义为 $-\log_2p$. 两个独立结果的联合概率为乘积, 对数把乘积变为和, 因而信息量相加. 概率 $1/2$ 的结果携带一个比特, 概率 $1/8$ 的结果携带三个比特. 这不是说看到一个罕见结果就一定实际写了三个二进制位, 而是提供了平均编码长度的标尺.

\begin{definition}[离散熵]
设 $X$ 的取值集合有限或可数, 质量函数为 $p(x)$. 定义
\[
 H(X)=-\sum_xp(x)\log_2p(x)=\E[-\log_2p(X)].
\]
约定 $0\log_20=0$, 因为 $\lim_{t\downarrow0}t\log_2t=0$. 可数情形允许 $H(X)=+\infty$; 正概率结果的信息量均非负, 所以此级数没有正负抵消的歧义.
\end{definition}
熵依赖各结果的概率, 不依赖给结果起什么名称. 对概率为零的结果, $-\log_20$ 本身无穷, 但该结果不会发生, 在期望和中其贡献按上述连续延拓取零.

\begin{example}
计算常量、$m$ 点均匀分布、伯努利分布和正整数上的几何分布的熵.
\end{example}
\begin{solution}
常量的熵为 $-1\log_21=0$. $m$ 点均匀分布的熵为
$-m(1/m)\log_2(1/m)=\log_2m$. 若 $X\sim\Bin(1,p)$, 则
\[
 H(X)=h(p):=-p\log_2p-(1-p)\log_2(1-p).
\]
在 $0<p<1$ 上,
\[
 h'(p)=\log_2\frac{1-p}{p},\qquad
 h''(p)=-\frac1{\ln2}\left(\frac1p+\frac1{1-p}\right)<0.
\]
所以 $h$ 在 $p=1/2$ 取得最大值 $1$, 在两端为零.
若 $G$ 的质量函数为 $p(1-p)^{j-1}$, $j\geq1$, $0<p<1$, 则
\begin{align*}
 H(G)&=-\sum_{j\geq1}p(1-p)^{j-1}\bigl[\log_2p+(j-1)\log_2(1-p)\bigr]\\
 &=-\log_2p-\frac{1-p}{p}\log_2(1-p)=\frac{h(p)}p.
\end{align*}
这里使用 $\E(G-1)=(1-p)/p$. 特别地 $p=1/2$ 时 $H(G)=2$, 虽然可能取值有无穷多个, 平均信息量仍有限.
\end{solution}

\begin{proposition}[熵的凹性]
在有限概率单纯形 $\{(p_1,\ldots,p_m):p_i\geq0,\sum p_i=1\}$ 上,
$H(p)=-\sum_i p_i\log_2p_i$ 为严格凹函数. 即对两个不同概率向量 $p,q$ 及 $0<t<1$,
\[
 H(tp+(1-t)q)>tH(p)+(1-t)H(q).
\]
\end{proposition}
\begin{proof}
令 $g(s)=-s\log_2s$, $g(0)=0$. 在 $(0,1]$ 上 $g''(s)=-1/(s\ln2)<0$, 因而任意两个不同正数之间严格凹. 若一个端点为零, 另一个为 $a>0$, 则
\[
 g(ta)-tg(a)=-ta\log_2t>0\quad(0<t<1),
\]
仍有严格凹性. 对每个坐标应用凹性再相加得到不等式. $p\ne q$ 时至少一个坐标不同, 该坐标的不等式严格, 因而总和严格.
\end{proof}
凹性的意思是: 先随机选择一种分布, 又不告诉观察者选了哪一种, 混合分布的熵至少等于选择已知时的平均熵. 后面将以此证明条件化平均减少不确定性.

\originalsection{相对熵与最大熵}
比较两个概率分布需要一个非负量. 我们先证明它, 再用它给出最大熵和编码下界, 这样各结论有同一个明确的依据.
\begin{definition}[相对熵]
对同一有限集合上的分布 $p,q$, 定义
\[
 D(p\Vert q)=\sum_{i:p_i>0}p_i\log_2\frac{p_i}{q_i}.
\]
若某个 $p_i>0$ 而 $q_i=0$, 则定义 $D(p\Vert q)=+\infty$.
\end{definition}
相对熵一般不对称, 因而不是距离. 它衡量的是使用 $q$ 作为概率描述, 与实际分布 $p$ 不符所带来的平均对数损失.

\begin{lemma}[相对熵非负]
$D(p\Vert q)\geq0$, 等号成立当且仅当 $p=q$.
\end{lemma}
\begin{proof}
无穷的情况已经成立. 否则对所有 $p_i>0$ 使用 $\ln z\leq z-1$, 得到
\[
 (\ln2)D(p\Vert q)
 =-\sum_{p_i>0}p_i\ln\frac{q_i}{p_i}
 \geq\sum_{p_i>0}(p_i-q_i)
 =1-\sum_{p_i>0}q_i\geq0.
\]
不等式 $\ln z\leq z-1$ 可由 $z-1-\ln z$ 在 $z=1$ 取得最小值零证明; 等号只在 $z=1$ 成立. 所以总等号要求每个正概率坐标满足 $q_i=p_i$, 并且 $q$ 在其余坐标上总质量为零, 即 $p=q$. 反方向直接代入.
\end{proof}

\begin{theorem}[有限集合上的最大熵]
若 $X$ 至多有 $m$ 个可能取值, 则
\[
 0\leq H(X)\leq\log_2m.
\]
左等号当且仅当 $X$ 几乎必然为常量; 右等号当且仅当在指定的 $m$ 个取值上均匀分布.
\end{theorem}
\begin{proof}
各项 $-p_i\log_2p_i\geq0$. 若其和为零, 每个正 $p_i$ 必须为 $1$, 因而只有一个正概率取值. 取 $q_i=1/m$, 有
\[
 D(p\Vert q)=\sum_i p_i\log_2p_i+\log_2m=\log_2m-H(p).
\]
相对熵非负和等号条件分别给出右侧不等式及其等号条件.
\end{proof}
“有 $m$ 个可能结果”与“有 $\log_2m$ 比特的不确定性”只有在均匀分布时相同. 一个结果概率接近 $1$ 时, 即使其余结果很多, 熵也可能很小.

\originalsection{联合熵、条件熵与链式法则}
\begin{definition}
设 $X,Y$ 取有限多个值, 联合质量函数为 $p(x,y)$, 边缘质量函数为 $p_X,p_Y$. 定义
\[
 H(X,Y)=-\sum_{x,y}p(x,y)\log_2p(x,y).
\]
对 $p_Y(y)>0$, 定义
\[
 H(X\mid Y=y)=-\sum_xp(x\mid y)\log_2p(x\mid y),
 \qquad
 H(X\mid Y)=\sum_yp_Y(y)H(X\mid Y=y).
\]
零概率条件值在平均中不作贡献, 无需给它指定条件分布.
\end{definition}
$H(X\mid Y=y)$ 随观察值 $y$ 改变, 而 $H(X\mid Y)$ 是观察 $Y$ 前计算的平均剩余信息量, 是一个数, 不是条件期望随机变量.

\begin{theorem}[熵的链式法则]
\[
 H(X,Y)=H(Y)+H(X\mid Y)=H(X)+H(Y\mid X).
\]
更一般地,
\[
 H(X_1,\ldots,X_n)=H(X_1)+\sum_{j=2}^nH(X_j\mid X_1,\ldots,X_{j-1}).
\]
\end{theorem}
\begin{proof}
对 $p(x,y)>0$, 因 $p(x,y)=p_Y(y)p(x\mid y)$,
\begin{align*}
 H(X,Y)&=-\sum_{x,y}p(x,y)\log_2p_Y(y)
         -\sum_{x,y}p(x,y)\log_2p(x\mid y)\\
 &=-\sum_yp_Y(y)\log_2p_Y(y)
   +\sum_yp_Y(y)\left[-\sum_xp(x\mid y)\log_2p(x\mid y)\right].
\end{align*}
有限和允许改变求和次序; 零概率项按定义为零. 这就是第一式. 交换 $X,Y$ 得到第二式. 对向量 $(X_1,\ldots,X_{n-1})$ 和 $X_n$ 使用两变量公式并归纳, 得到多变量公式.
\end{proof}

\begin{theorem}[条件化减少平均熵]
\[
 0\leq H(X\mid Y)\leq H(X),\qquad H(X,Y)\leq H(X)+H(Y).
\]
$H(X\mid Y)=H(X)$ 当且仅当 $X,Y$ 独立. $H(X\mid Y)=0$ 当且仅当存在函数 $f$ 使 $X=f(Y)$ 几乎必然成立.
\end{theorem}
\begin{proof}
令 $v_y=(p(x_1\mid y),\ldots,p(x_m\mid y))$. 边缘概率向量为
$v=\sum_y p_Y(y)v_y$. 熵的凹性给出
\[
 H(X)=H\left(\sum_yp_Y(y)v_y\right)
 \geq\sum_yp_Y(y)H(v_y)=H(X\mid Y).
\]
严格凹性表明等号要求所有具有正权重的 $v_y$ 相同, 从而它们都等于 $v$. 这正是 $p(x\mid y)=p_X(x)$, 等价于独立. 也可直接计算
\[
 H(X)-H(X\mid Y)=D(p_{X,Y}\Vert p_Xp_Y)\geq0,
\]
其中任何 $p(x,y)>0$ 的点两边缘概率都正, 不会出现分母为零.
非负性来自定义; 联合熵上界来自链式法则. 平均条件熵为零时, 每个正概率 $y$ 上的条件熵都为零, 该条件分布集中在唯一的 $f(y)$ 上. 反之若 $X=f(Y)$, 这些条件分布都集中在单点, 条件熵为零.
\end{proof}
这个定理是平均意义的陈述. 某个特定观察值可能使熵增加. 例如 $Y$ 以 $9/10$ 的概率等于 $0$, 此时 $X=0$; $Y=1$ 时 $X$ 在 $1,2$ 上各占一半. 则 $H(X\mid Y=1)=1$, 而 $H(X)=h(1/10)+1/10<1$. 罕见观察揭示“目前进入了更不确定的情形”, 与平均条件熵 $H(X\mid Y)=1/10$ 较小并不矛盾.

\begin{definition}[互信息]
$ I(X;Y)=H(X)+H(Y)-H(X,Y)=H(X)-H(X\mid Y)$.
\end{definition}
链式法则说明它对称; 上一定理说明它非负, 且为零恰好是独立. 当 $Y=f(X)$ 时, $I(X;Y)=H(Y)$, 因为获知 $X$ 已经确定了 $Y$.

\begin{example}
$X$ 在 $\{0,1,2,3\}$ 上均匀分布, $Y$ 是 $X$ 的奇偶性. 计算联合熵、条件熵和互信息.
\end{example}
\begin{solution}
$H(X)=2$, $Y$ 为公平二元变量, $H(Y)=1$. 给定奇偶性后仍有两个等可能 $X$ 值, 所以 $H(X\mid Y)=1$. 给定 $X$ 后 $Y$ 确定, 故 $H(Y\mid X)=0$. 链式法则给出 $H(X,Y)=2$, 互信息为 $1$. 观察奇偶性正好消除一个比特的不确定性.
\end{solution}

\originalsection{二叉决策树与前缀码}
一个是非问题给出一个二进制位. 接下来问什么可以依赖已有回答, 所以完整策略是一棵二叉树: 根是开始, 内部结点是问题, 两条边是两个回答, 叶子是已经确定的结果. 某个结果所需问题数就是对应叶子的深度.

\begin{definition}[前缀码]
为每个可能结果 $x_i$ 指定有限二进制串 $c_i$. 若任一个码字都不是另一个码字的前缀, 则称为前缀码. 其码长为 $\ell_i$, 平均码长为 $L=\sum_ip_i\ell_i$.
\end{definition}
前缀条件允许即时解码: 从左到右走二叉树, 一到叶子就读出一个结果并回到根, 无需额外分隔符. 只有一个可能结果时允许空码字, 长度为零. 一般的唯一可解码码不必是前缀码; 本章证明针对与是非问题策略直接对应的前缀码.

\begin{lemma}[Kraft 不等式及其逆命题]
有限前缀码的长度满足 $\sum_i2^{-\ell_i}\leq1$. 反之, 非负整数 $\ell_i$ 若满足此不等式, 就存在具有这些长度的前缀码.
\end{lemma}
\begin{proof}
令 $L_*=\max_i\ell_i$, 考虑所有长 $L_*$ 的二进制串, 共 $2^{L_*}$ 个. 以 $c_i$ 为开头的串有 $2^{L_*-\ell_i}$ 个. 前缀条件保证这些集合不相交, 因而相加不超过 $2^{L_*}$; 除以它便得不等式.

逆命题将长度从小到大排列. 在深度 $\ell_i$ 处, 每个已选的深度 $\ell_j\leq\ell_i$ 的叶子禁用 $2^{\ell_i-\ell_j}$ 个位置. 因而选择第 $i$ 个码字时已有禁用位置数为
$\sum_{j<i}2^{\ell_i-\ell_j}$. 由 $\sum_{j\leq i}2^{-\ell_j}\leq1$ 得
\[
 \sum_{j<i}2^{\ell_i-\ell_j}\leq2^{\ell_i}-1.
\]
所以至少有一个位置可选. 选定它, 并在之后禁用其后代. 因选择按深度递增, 新叶子不会成为旧叶子的祖先, 也不是旧叶子的后代. 依次完成, 得到所需前缀码.
\end{proof}

\begin{theorem}[无噪声编码界]
对有限分布 $p$, 任意前缀码的平均长度满足 $L\geq H(p)$. 存在前缀码使
\[
 H(p)\leq L<H(p)+1.
\]
当所有正概率均为 $2$ 的负整数次幂时, 可以取得 $L=H(p)$.
\end{theorem}
\begin{proof}
只编码正概率结果. 对给定码令 $a=\sum_i2^{-\ell_i}\leq1$, $q_i=2^{-\ell_i}/a$. 相对熵非负给出
\[
 D(p\Vert q)=\sum_ip_i\log_2p_i+\sum_ip_i\ell_i+\log_2a
 =L-H(p)+\log_2a\geq0.
\]
所以 $L\geq H(p)-\log_2a\geq H(p)$.

为构造上界, 取 $\ell_i=\lceil-\log_2p_i\rceil$. 因为
$2^{-\ell_i}\leq p_i$, Kraft 逆命题保证这些长度可实现. 又因
$-\log_2p_i\leq\ell_i<-\log_2p_i+1$, 乘以 $p_i$ 后求和得所需界.
若 $p_i=2^{-k_i}$, 取 $\ell_i=k_i$, 则 Kraft 和为 $1$, 且 $L=\sum_ip_ik_i=H(p)$.
\end{proof}
下界证明还给出等号条件: $a=1$ 且 $p_i=q_i$, 即 $p_i=2^{-\ell_i}$. 在决策树中, 这意味着每个被访问的分叉都把余下概率分成两半. 一般概率没有整数比特数, 因而不能期待单个符号编码恰好达到熵.

\begin{example}
一个传感器输出四种状态, 概率为 $(1/2,1/4,1/8,1/8)$. 构造达到熵下界的编码, 并解释如何解码.
\end{example}
\begin{solution}
四种状态记为 $a,b,c,d$, 选码字 $0,10,110,111$. 码长分别为 $1,2,3,3$, 没有码字是另一个的前缀. 平均长度
\[
 L=\frac12+\frac24+\frac38+\frac38=\frac74=H(X).
\]
例如收到串 $101110110$, 从根逐位读得 $10|111|0|110$, 即 $b,d,a,c$. 竖线是解码后标出的边界, 传输时不需要发送它们.
\end{solution}

\originalsection{Huffman 编码与分组压缩}
上述向上取整法很容易实现, 但不总是最优. Huffman 算法反复合并最小的两个概率, 最后从合并树读出二进制码. 为什么可以安全地把最小概率放在最深处, 需要一个交换论证.

\begin{proposition}[Huffman 算法的最优性]
对有限个正概率结果, 反复将两个最小概率 $p_a,p_b$ 合并为一个概率 $p_a+p_b$, 对缩小后的集合编码, 再把合并结果的码字后添 $0,1$, 得到平均码长最小的前缀码.
\end{proposition}
\begin{proof}
只有一个结果时结论成立. 对至少两个结果, 最优树可取为每个内部结点都有两个孩子的树: 若只有一个孩子, 删除该无效分叉会缩短其下所有叶子的码长. 这种树最多有 $m-1$ 个内部结点, 所以深度不超过 $m-1$, 可能树及叶子标号只有有限多个, 最小值确实存在.

最深的一个叶子的兄弟也是叶子, 否则兄弟还有更深的后代. 若 $p_i>p_j$ 而 $\ell_i>\ell_j$, 交换这两个叶子的标签会使平均码长改变
\[
 p_i\ell_j+p_j\ell_i-p_i\ell_i-p_j\ell_j
 =-(p_i-p_j)(\ell_i-\ell_j)<0.
\]
因此可以让两个最小概率占据最深的一对兄弟叶子, 而不增加平均码长. 将这对兄弟合并, 新树平均长度比原树减少 $p_a+p_b$. 合并后的树必须最优, 否则用更好的小树替换再拆开会改善原树. 归纳假设给出缩小后 Huffman 树最优, 拆开时增加的 $p_a+p_b$ 与小树选择无关, 故所得原树最优.
\end{proof}

\begin{example}
对分布 $(0.4,0.3,0.2,0.1)$ 求一个最优前缀码.
\end{example}
\begin{solution}
先合并 $0.1,0.2$ 为 $0.3$, 再将两个 $0.3$ 合并为 $0.6$, 最后合并 $0.4,0.6$. 可以取码字 $0,10,110,111$, 分别对应 $0.4,0.3,0.2,0.1$. 平均长度为
$0.4+2(0.3)+3(0.2)+3(0.1)=1.9$.
熵为 $-\sum p_i\log_2p_i\approx1.8464$, 小于 $1.9$. 两个 $0.3$ 的处理顺序及每次左右分支可以互换, 最优码未必唯一.
\end{solution}

\begin{corollary}[分组编码接近熵率]
若 $X_1,\ldots,X_n$ 独立同分布, 每个变量取有限多个值, 则存在对整个向量的一次前缀编码, 平均总长度 $L_n$ 满足
\[
 nH(X_1)\leq L_n<nH(X_1)+1,
 \qquad \frac{L_n}{n}\longrightarrow H(X_1).
\]
\end{corollary}
\begin{proof}
独立性使每个条件熵等于无条件熵, 链式法则给出
$H(X_1,\ldots,X_n)=nH(X_1)$. 将整个向量当作一个具有有限个可能取值的随机变量应用编码定理即可. 最后的一个比特误差摊到 $n$ 个符号上成为 $1/n$.
\end{proof}
这解释了分组压缩的收益: 单个符号的码长必须是整数, 但一组符号可以共同承担取整误差. 这一定理是平均长度结果, 并不保证每一条输入序列都得到短码; 罕见序列仍可能很长.

\originalsection{练习与解答}
\begin{exercise}
$X$ 为公平比特, $B$ 为独立的伯努利 $(q)$ 比特, $Y=X+B\pmod2$. 求 $H(Y)$、$H(Y\mid X)$、$I(X;Y)$, 并解释 $q=0,1/2,1$.
\end{exercise}
\begin{solution}
$\Prob(Y=1)=\tfrac12q+\tfrac12(1-q)=1/2$, 故 $H(Y)=1$. 给定 $X$, $Y$ 的条件分布是 $(q,1-q)$ 或 $(1-q,q)$, 两者熵均为 $h(q)$. 因而 $H(Y\mid X)=h(q)$, $I(X;Y)=1-h(q)$. $q=0$ 时 $Y=X$, $q=1$ 时 $Y=1-X$, 二者均完全确定 $X$, 互信息为 $1$. $q=1/2$ 时输出不含关于输入的信息, 两者独立, 互信息为零.
\end{solution}

\begin{exercise}
三种输出的概率为 $(1/2,1/3,1/6)$. 求最优前缀码的平均长度, 并比较单符号编码与 $n$ 符号分组编码的每符号误差上界.
\end{exercise}
\begin{solution}
Huffman 算法合并 $1/3,1/6$, 得到两叶概率 $1/2,1/2$. 取码字 $0,10,11$, 平均长度为 $3/2$. 熵为
\[
 H=\frac12+\frac13\log_23+\frac16\log_26
 =\frac23+\frac12\log_23\approx1.4591.
\]
单符号最优码多花约 $0.0409$ 比特; 一般分组上界则保证 $L_n/n-H<1/n$. 分组编码要处理更多码字, 所以接近熵率同时带来存储和延迟方面的代价.
\end{solution}

\begin{exercise}
证明 $H(X\mid Y,Z)\leq H(X\mid Y)$, 并给出等号条件.
\end{exercise}
\begin{solution}
在每个正概率 $Y=y$ 的条件空间上应用条件化减少熵定理,
\[
 H(X\mid Y=y)\geq\sum_z\Prob(Z=z\mid Y=y)H(X\mid Y=y,Z=z).
\]
乘以 $\Prob(Y=y)$ 求和得到所需不等式. 等号成立当且仅当对每个正概率 $y$, 在该条件分布下 $X,Z$ 独立, 即
$p(x,z\mid y)=p(x\mid y)p(z\mid y)$.
\end{solution}

\originalchapter{鞅、可选停止与风险中性定价}
\sourcelecture{34--36}
本章同时使用官方补充讲义.

条件期望描述给定信息后的平均判断. 鞅把这种判断沿时间排列: 今天的数值等于考虑今天全部信息后对明天数值的条件期望. 这一性质比“各时刻均值相同”强得多. 本章先研究公平赌博与停止规则, 再在可直接计算的市场模型中建立风险中性概率, 最后推导欧洲期权的 Black--Scholes 公式. 金融部分始终讨论给定假设下的数学定价, 不把市场价格当作客观事件频率.

\originalsection{信息、过滤与鞅}
在掷硬币实验中, 时刻 $n$ 已知前 $n$ 次结果, 但未知下一次结果. 为准确表达“此刻知道什么”, 把所有此刻可以判定真假的事件收集起来.

\begin{definition}[过滤与适应过程]
在概率空间 $(\Omega,\mathcal F,\Prob)$ 上, 一个 $\sigma$ 代数是含 $\Omega$、对补集和可数并集封闭的事件集合. 过滤是一列递增的 $\sigma$ 代数
\[
 \mathcal F_0\subseteq\mathcal F_1\subseteq\cdots\subseteq\mathcal F.
\]
$\mathcal F_n$ 表示时刻 $n$ 的信息. 若每个 $X_n$ 都是 $\mathcal F_n$ 可测的, 即对每个实数 $a$, $\{X_n\leq a\}\in\mathcal F_n$, 则过程 $(X_n)$ 称为适应该过滤. 由 $X_0,\ldots,X_n$ 产生的最小 $\sigma$ 代数记为 $\sigma(X_0,\ldots,X_n)$, 给出自然过滤.
\end{definition}
可测性在这里的含义是 $X_n$ 的值可以由当前信息确定. 若另有一组观察 $Y_0,\ldots,Y_n$, 过滤可以包含它们, 不能默认只有价格本身的历史. 同一个过程在不同信息集合下, 可能具有不同的鞅性质.

在有限信息树上, $\E(Z\mid\mathcal F_n)$ 的定义很具体: 对每个当前已知的历史节点, 用条件概率在其可能后续结果上取平均. 一般可积随机变量的条件期望是一个 $\mathcal F_n$ 可测可积变量 $W$, 满足
\[
 \E[W\ind_A]=\E[Z\ind_A]\quad(A\in\mathcal F_n).
\]
它在几乎处处意义下唯一. 有限树和离散随机游走的实例可按有限或可数历史逐项条件化; 连续取值的增量及布朗运动实例则使用第8章已声明的一般条件期望存在性工具, 不在这里重新构造该定理. 使用的性质是线性、已知有界因子可提出、独立变量的条件期望等于其均值, 以及塔式法则
\[
 \E[\E(Z\mid\mathcal F_{n+1})\mid\mathcal F_n]=\E(Z\mid\mathcal F_n).
\]
塔式法则可从上述定义核对: 左边对任何 $A\in\mathcal F_n$ 的积分等于 $\E[Z\ind_A]$, 且对 $\mathcal F_n$ 可测, 因而由唯一性得到等式. 在有限树上这正是先按后一天分组平均, 再按今天分组平均等于直接按今天平均.

\begin{definition}[鞅]
适应过程 $(M_n)_{n\geq0}$ 若满足每个 $n$ 上 $\E|M_n|<\infty$, 且
\[
 \E(M_{n+1}\mid\mathcal F_n)=M_n\quad\text{几乎必然},
\]
则称为相对于 $(\mathcal F_n)$ 和 $\Prob$ 的鞅. 若等号换为 $\geq$ 或 $\leq$, 则分别称为次鞅或上鞅.
\end{definition}
可积性保证条件期望有意义. 适应性保证右侧是当前已知的量. 条件均值条件则要求对每一种具有正概率的当前信息都公平, 不能只要求整体平均公平. 对鞅反复使用塔式法则得
\[
 \E(M_m\mid\mathcal F_n)=M_n\quad(m\geq n),\qquad
 \E M_n=\E M_0.
\]
第一个公式用 $m-n$ 归纳, 第二个公式再取无条件期望.

\begin{example}
设 $A_1,A_2,\ldots$ 独立可积, $\mathcal F_n=\sigma(A_1,\ldots,A_n)$. 证明下列鞅构造: 零均值增量之和, 单位均值因子之积, 单位方差增量的平方补偿, 以及逐次更新的条件期望.
\end{example}
\begin{solution}
若 $\E A_i=0$, 则 $S_n=\sum_{i=1}^nA_i$ 适应且可积, 并且
\[
 \E(S_{n+1}\mid\mathcal F_n)=S_n+\E A_{n+1}=S_n.
\]
若改为 $\E A_i=1$, 则 $P_n=\prod_{i=1}^nA_i$, $P_0=1$, 也是鞅. 独立性给出 $\E|P_n|=\prod_{i=1}^n\E|A_i|<\infty$, 而
$\E(P_{n+1}\mid\mathcal F_n)=P_n\E A_{n+1}=P_n$.
这包括每步独立乘以 $1+a$ 或 $1-a$、两者各概率 $1/2$ 的价格模型, $0<a<1$.

若 $\E A_i=0$, $\E A_i^2=1$, 则 $S_n^2-n$ 可积, 且
\begin{align*}
 \E[S_{n+1}^2-(n+1)\mid\mathcal F_n]
 &=S_n^2+2S_n\E A_{n+1}+\E A_{n+1}^2-(n+1)\\
 &=S_n^2-n.
\end{align*}
其中 $S_nA_{n+1}$ 可积由独立性和两者可积保证. 最后若 $Z$ 可积, $M_n=\E(Z\mid\mathcal F_n)$ 适应且
$\E|M_n|\leq\E|Z|$, 塔式法则给出
$\E(M_{n+1}\mid\mathcal F_n)=M_n$.
\end{solution}

\begin{example}
$A$ 以相同概率取 $1,-1$, 令 $X_0=0$, $X_n=(-1)^nA$ ($n\geq1$), 使用自然过滤. 判定它是否为鞅.
\end{example}
\begin{solution}
每个 $X_n$ 均值为零, 但观察 $X_1$ 后便确定 $A$. 所以 $n\geq1$ 时下一步已经完全可预测,
\[
 \E(X_{n+1}\mid\mathcal F_n)=-X_n\ne X_n.
\]
故不是鞅. 即使随机游走本来是鞅, 若把所有未来掷币结果在时刻零就告诉观察者, 它也不再对扩大后的过滤是鞅.
\end{solution}

若 $Z=\ind_A$, 条件期望鞅就是 $\Prob(A\mid\mathcal F_n)$. 例如观察一场比赛逐分钟更新获胜概率, 只要所有更新来自同一个事先选定的概率模型, 就形成鞅. 一条具体轨迹可以不断上升或下降; 鞅的等式是在每个历史下对尚未看到的信息取平均. 若事先概率模型估计不准, 该模型下仍可有鞅, 所以“概率更新是鞅”本身不证明估计符合现实.

\originalsection{停止时间与停止过程}
\begin{definition}[停止时间]
随机变量 $T:\Omega\to\{0,1,2,\ldots\}\cup\{\infty\}$ 称为停止时间, 若对每个 $n\geq0$,
\[
 \{T\leq n\}\in\mathcal F_n.
\]
即到时刻 $n$ 时能判断是否已经停止.
\end{definition}
这等价于 $\{T=n\}\in\mathcal F_n$ 对每个 $n$ 成立. 一方向使用
$\{T=n\}=\{T\leq n\}\setminus\{T\leq n-1\}$ 和过滤递增性; 在 $n=0$ 时前一个事件视为空集. 另一方向使用
$\{T\leq n\}=\bigcup_{k=0}^n\{T=k\}$.

首次到达集合 $B$ 的时间 $T=\inf\{n:X_n\in B\}$ 是停止时间, 因为
$\{T\leq n\}=\bigcup_{k=0}^n\{X_k\in B\}$ 属于 $\mathcal F_n$. 第 $j$ 次访问 $B$ 的时间也一样, 因为截至 $n$ 的访问次数 $\sum_{k=0}^n\ind_{\{X_k\in B\}}$ 可由当前历史计算. 这里要求 $B$ 为可测集合, 空集合的首次到达时间约定为 $\infty$.

“未来十步内的最后一次到达零”通常不是停止时间. 对简单对称随机游走, 定义 $L=\max\{0\leq n\leq2:S_n=0\}$. 同样的时刻零信息下, 两条后续路径 $++$ 与 $+-$ 分别使 $L=0$ 与 $L=2$, 所以 $\{L=0\}$ 不属于平凡的 $\mathcal F_0$. “截至固定期限取得最大值的时刻”也需要看未来, 一般不是停止时间. 在某个特定退化过程中它们可能是, 因而判断应看事件判据, 不能只靠规则名称.

\begin{proposition}[停止后的过程仍是鞅]
若 $M$ 是鞅, $T$ 是停止时间, 则 $M_{n\wedge T}$ 也是鞅. 其中 $n\wedge T=\min(n,T)$.
\end{proposition}
\begin{proof}
$M_{n\wedge T}=\sum_{k=0}^{n-1}M_k\ind_{\{T=k\}}+M_n\ind_{\{T\geq n\}}$, 所以适应且可积. 一步增量为
\[
 M_{(n+1)\wedge T}-M_{n\wedge T}
 =\ind_{\{T>n\}}(M_{n+1}-M_n).
\]
$\{T>n\}\in\mathcal F_n$, 因而取条件期望并将已知指示因子提出, 得到零. 这正是停止过程的鞅条件.
\end{proof}
停止之后把数值冻结, 于是每一步都只在“尚未停止”时参与下一次公平变化. 是否参与由变化发生之前的信息决定, 所以一步平均仍为零.

\originalsection{可选停止定理及其边界}
\begin{theorem}[有界停止时间]
若 $M$ 为可积鞅, $T$ 为停止时间, 且 $T\leq N$ 几乎必然, 其中 $N$ 是确定的有限整数, 则 $M_T$ 可积且
\[
 \E M_T=\E M_0.
\]
\end{theorem}
\begin{proof}
可积性由 $|M_T|\leq\sum_{k=0}^N|M_k|$ 保证. 路径上有恒等式
\[
 M_T=M_0+\sum_{k=0}^{N-1}\ind_{\{T>k\}}(M_{k+1}-M_k).
\]
对每项, 指示变量由时刻 $k$ 的信息决定, 所以
\[
 \E[\ind_{\{T>k\}}(M_{k+1}-M_k)]
 =\E\left[\ind_{\{T>k\}}\E(M_{k+1}-M_k\mid\mathcal F_k)\right]=0.
\]
有限求和后得到结论. 这也揭示了为何必须用停止时间: 若是否参与下一步取决于下一步输赢, 指示变量不能从条件期望提出.
\end{proof}

\begin{theorem}[有界过程与几乎必然有限的停止时间]
若 $M$ 为鞅, $T<\infty$ 几乎必然, 且存在确定常数 $C$ 使
$|M_{n\wedge T}|\leq C$ 对所有 $n$ 几乎必然成立, 则
$\E M_T=\E M_0$. 特别地, 一致有界鞅满足此条件.
\end{theorem}
\begin{proof}
$T\wedge N$ 是有界停止时间, 所以前一定理给出
$\E M_{T\wedge N}=\E M_0$. 在 $\{T\leq N\}$ 上,
$M_{T\wedge N}=M_T$; 在其余事件上两者的差绝对值至多 $2C$. 因而
\[
 |\E M_{T\wedge N}-\E M_T|\leq2C\Prob(T>N)\longrightarrow0.
\]
最后一个极限使用 $\{T>N\}$ 递减到 $\{T=\infty\}$, 后者概率为零. 这给出期待等式, 无需把极限交换当作无条件操作.
\end{proof}
这里明确要求 $T$ 几乎必然有限, 否则 $M_T$ 未定义, 除非先证明并定义无穷时刻的极限. 所谓鞅有界是所有时刻共用一个确定界, 不是“对每个固定 $n$, $M_n$ 都有界”. 简单随机游走在每个固定 $n$ 都满足 $|S_n|\leq n$, 仍不满足一致有界条件.

\begin{proposition}[可预测下注]
设 $H_k$ 是 $\mathcal F_k$ 可测有界变量. 若 $M$ 是鞅, 则截至固定期限 $N$ 的净收益
\[
 G_N=\sum_{k=0}^{N-1}H_k(M_{k+1}-M_k)
\]
具有期望零; $G_n$ 本身是鞅.
\end{proposition}
\begin{proof}
有限和可积. 每项取条件期望为
$H_k\E(M_{k+1}-M_k\mid\mathcal F_k)=0$. 一步增量同样满足条件均值为零, 因而收益过程是鞅.
\end{proof}
下注大小可以依赖已发生的输赢, 但不能偷看本次结果. 这一定理允许在固定期限内采用相当复杂的策略, 不把无穷期、无资金上限的策略自动包括进来.

\originalsection{赌场破产、往返访问与无界停止反例}
设整数 $i,b$ 满足 $0<i<b$, $S_0=i$, 每步独立以概率 $1/2$ 加 $1$ 或减 $1$. 定义
$T=\inf\{n:S_n\in\{0,b\}\}$. 这是首次到达边界的停止时间.

首先证明 $T<\infty$ 几乎必然. 任意尚未退出的状态, 接下来连续 $b$ 次向上都必然触及上边界, 该串的概率为 $2^{-b}$. 将时间分成长度 $b$ 的块, 每块给定前史后仍有至少 $2^{-b}$ 的退出概率, 因而
\[
 \Prob(T>kb)\leq(1-2^{-b})^k\longrightarrow0.
\]
这个界虽不精确, 但足以支持下面所有停止论证.

\begin{theorem}[公平赌场破产概率与平均时长]
上述随机游走满足
\[
 \Prob(S_T=b)=\frac ib,\qquad \Prob(S_T=0)=1-\frac ib,
 \qquad \E T=i(b-i).
\]
\end{theorem}
\begin{proof}
停止过程 $S_{n\wedge T}$ 值在 $[0,b]$ 内, 是有界鞅, 所以
$\E S_T=i$. 若成功概率为 $p$, 则 $bp=i$, 得 $p=i/b$.

由平方补偿例子, $S_n^2-n$ 是鞅. 对有界停止时间 $T\wedge N$ 使用可选停止,
\[
 \E[S_{T\wedge N}^2]-\E(T\wedge N)=i^2.
\]
左边第一项因 $0\leq S_{T\wedge N}\leq b$ 及 $T<\infty$ 而趋于
$\E S_T^2=b^2(i/b)=bi$. 具体误差至多 $b^2\Prob(T>N)$.
又 $\E(T\wedge N)=\sum_{k=0}^{N-1}\Prob(T>k)$, 非负级数的部分和趋于 $\E T$ (可由整数变量的尾和公式直接定义). 所以
$\E T=bi-i^2=i(b-i)$, 同时证明了平均时长有限.
\end{proof}
在一般区间 $[a,b]$ 内, 若从 $c$ 出发的鞅几乎必然在有限时刻无跳越地触及某一端点, 且停止过程一致有界, 则先触及 $b$ 的概率为 $(c-a)/(b-a)$. 若过程可能越过端点, 终值未必恰为 $a,b$, 不能照搬这条线性公式.

\begin{example}
公平随机游走从 $6$ 出发, 遇到 $0$ 或 $12$ 便结束. 求它在结束前依次达到 $4,8,4,8$ 的概率.
\end{example}
\begin{solution}
先从 $6$ 在 $[4,12]$ 内求先达 $4$ 的概率, 为 $(12-6)/(12-4)=3/4$. 达到 $4$ 后, 在 $[0,8]$ 内先达 $8$ 的概率为 $4/8=1/2$. 从 $8$ 再在 $[4,12]$ 内先达 $4$ 的概率为 $1/2$. 最后由 $4$ 再先达 $8$ 的概率为 $1/2$. 总概率为 $3/32$.
这里每次到达后未来掷币仍独立公平. 为避免把这个事实仅作为定理名引用, 可对每个可能到达时间 $n$ 条件化: 该到达事件只由前 $n$ 次掷币决定, 而任一固定长的后续掷币串与这前史独立. 后续串的概率不随 $n$ 或具体前史改变, 再对 $n$ 求和, 就得到每次重启时相同的退出公式和概率乘积.
\end{solution}

\begin{example}
从 $S_0=0$ 出发的简单对称随机游走, 令 $T=\inf\{n:S_n=1\}$. 证明 $T<\infty$ 几乎必然, 但 $\E S_T\ne\E S_0$.
\end{example}
\begin{solution}
对每个正整数 $a$, 在 $[-a,1]$ 内先到 $1$ 的概率为 $a/(a+1)$. 这一事件包含于 $\{T<\infty\}$, 所以
$\Prob(T<\infty)\geq a/(a+1)$ 对每个 $a$ 成立, 令 $a\to\infty$ 得概率为 $1$. 然而 $S_T=1$, 所以 $\E S_T=1\ne0$.
每个 $T\wedge N$ 仍满足 $\E S_{T\wedge N}=0$. 虽然
$S_{T\wedge N}\to1$ 几乎必然, 期望却不收敛到 $1$: 未停止的少数路径可能已经走到很负的值, 它们维持整体均值为零. 这里不存在共同的确定界控制所有停止值.
\end{solution}

另一反例是输后加倍. 第一次下注 $1$, 每次输后下注额翻倍, 首次赢后停止. 前 $n$ 次全输的概率为 $2^{-n}$, 损失为 $2^n-1$; 只要前 $n$ 次内赢过一次, 收益为 $1$. 因而截至 $n$ 的期望为
\[
 (1-2^{-n})\cdot1+2^{-n}\bigl(-(2^n-1)\bigr)=0.
\]
无限策略几乎必然最终赢且收益 $1$, 但最大资金需求没有界. 这与有界期限策略的零期望不矛盾, 也表明不能仅凭几乎必然有限的停止时间就交换期望和极限.

\originalsection{无套利与一期二叉树}
现在以股票和无风险账户建立一个完全可以求解的市场. 股票现价 $S>0$, 一期后只能为 $uS$ 或 $dS$, 其中 $0<d<u$. 无风险账户一期增长因子为 $R=1+r>0$. 可以买卖任意实数股, 允许借贷和卖空, 同一资产买卖价格相同, 无交易费用、税费或违约, 股票不分红. 两种状态都具有正的实际概率. 这些是假设, 不由概率公理推导.

\begin{definition}[套利与复制]
一期套利是初始净投入为零、两种终值均非负、至少一种终值严格为正的组合. 若一个组合在每一种状态的终值均与某合约相同, 称为该合约的复制组合.
\end{definition}
初始负成本且终值非负的策略也能变为零成本套利: 将收到的钱存入无风险账户. 反之若两个相同终值的组合价格不同, 买便宜的、卖贵的, 再存入差价, 即构成套利. 所以无套利迫使复制组合与合约价格相同.

\begin{theorem}[一期无套利条件]
上述市场无套利当且仅当
\[
 d<R<u.
\]
此时唯一的风险中性向上概率为
\[
 q=\frac{R-d}{u-d}\in(0,1),\qquad qu+(1-q)d=R.
\]
\end{theorem}
\begin{proof}
若 $R\leq d$, 借入 $S$ 买一股, 初始成本零, 终值为
$S(u-R)$ 或 $S(d-R)$, 均非负且上状态严格正, 是套利.
若 $R\geq u$, 卖空一股把所得 $S$ 存入账户, 终值为
$S(R-u)$ 或 $S(R-d)$, 同样是套利.

反之若 $d<R<u$, 上述 $q$ 在 $(0,1)$. 任一组合由 $\Delta$ 股股票及金额 $B$ 的账户组成, 初价 $V_0=\Delta S+B$, 两终值为
$V_u=\Delta uS+RB$, $V_d=\Delta dS+RB$. 由 $qu+(1-q)d=R$,
\[
 qV_u+(1-q)V_d=R(\Delta S+B)=RV_0.
\]
若 $V_0=0$, 且终值非负并至少一个严格正, 左侧由于权重均正便严格正, 与右侧零矛盾. 所以无套利. 解均值方程给出的 $q$ 唯一.
\end{proof}
这里实际向上概率 $p$ 没有出现在公式里. $q$ 的作用是把价格写成贴现均值, 不要求 $q=p$. 无套利禁止无风险的确定盈利, 允许带风险的正期望投资; 因而“实际概率下有正期望收益”不等同于套利.

\begin{theorem}[复制与一期定价]
合约在上涨和下跌时分别支付 $H_u,H_d$. 其唯一复制组合为
\[
 \Delta=\frac{H_u-H_d}{S(u-d)},\qquad
 B=\frac{uH_d-dH_u}{R(u-d)}.
\]
无套利价格为
\[
 V_0=\Delta S+B=\frac1R\bigl(qH_u+(1-q)H_d\bigr).
\]
\end{theorem}
\begin{proof}
要求 $\Delta uS+RB=H_u$, $\Delta dS+RB=H_d$. 两式相减得到 $\Delta$, 代回得 $B$. 因 $S(u-d)\ne0$, 线性方程唯一可解. 对两终值用 $q,1-q$ 加权, 得 $RV_0=qH_u+(1-q)H_d$. 若合约报价偏离 $V_0$, 同终值组合的买卖差价就给出套利, 故价格唯一.
\end{proof}

\begin{example}
股票现价 $80$, 一期后为 $100$ 或 $60$, 账户增长因子为 $1.05$. 执行价为 $85$ 的看涨合约终值为 $(S_1-85)^+$. 求复制组合和价格.
\end{example}
\begin{solution}
$u=1.25$, $d=0.75$, $q=(1.05-0.75)/0.5=0.6$. 两种支付为 $15,0$. 因而
\[
 \Delta=\frac{15}{40}=\frac38,\qquad
 B=\frac{-0.75(15)}{1.05(0.5)}=-\frac{150}7\approx-21.4286.
\]
购买 $3/8$ 股并借入 $150/7$, 初始净投入 $30-150/7=60/7$. 上涨终值 $37.5-22.5=15$, 下跌终值 $22.5-22.5=0$. 同样,
$V_0=(0.6\cdot15)/1.05=60/7$.
\end{solution}

\originalsection{多期复制与贴现鞅}
把一期模型重复 $N$ 次. 每个历史节点的下一步仍为当前价乘以 $u$ 或 $d$, 两个分支的实际条件概率均为正, 无风险增长因子仍为 $R$, 且 $d<R<u$. 在完整的 $2^N$ 条涨跌路径上, 定义新概率 $Q$: 每次上涨概率为 $q$, 下跌为 $1-q$, 独立. 实际概率与 $Q$ 因而具有相同的零概率事件. 对一条有 $k$ 次上涨的具体路径, $Q$ 概率为 $q^k(1-q)^{N-k}$; 同一终价对应多条路径时应加上路径数, 不能把它们当作单一路径.

\begin{proposition}
在 $Q$ 下, 贴现股价 $S_n/R^n$ 是鞅. 任意终期支付 $H$ (可依赖完整路径) 的价格为
\[
 V_n=R^{-(N-n)}\E_Q(H\mid\mathcal F_n),\qquad
 V_0=R^{-N}\E_QH.
\]
此价格可以由只使用当前信息的股票和账户交易逐期复制, 且 $V_n/R^n$ 为 $Q$ 鞅.
\end{proposition}
\begin{proof}
每个节点上
$\E_Q(S_{n+1}\mid\mathcal F_n)=S_n[qu+(1-q)d]=RS_n$, 所以贴现过程为鞅. 从终期令 $V_N=H$, 逐层向后定义
\[
 V_n=\frac1R(qV_{n+1}^{u}+(1-q)V_{n+1}^{d}).
\]
在该节点应用一期复制公式, 选择
\[
 \Delta_n=\frac{V_{n+1}^{u}-V_{n+1}^{d}}{S_n(u-d)},\qquad
 B_n=V_n-\Delta_n S_n.
\]
下一步组合恰有价值 $V_{n+1}$. 到达该下一节点时, 将组合重新调整成该节点所需的股票和账户持仓; 调整前后总价值同为 $V_{n+1}$, 不需要外部注资, 这就是自融资. 归纳直到 $N$, 便复制 $H$.

逐层条件化或塔式法则把递归展开为
$V_n=R^{-(N-n)}\E_Q(H\mid\mathcal F_n)$. 因而
$V_n/R^n=R^{-N}\E_Q(H\mid\mathcal F_n)$ 是条件期望鞅. 同终值复制组合的无套利论证逐期成立, 故得到价格唯一.
\end{proof}
这段论证实际构造了概率和策略, 没有调用一般金融基本定理. 完整二叉树中任何终期支付都能复制, 称为市场完备. 如果每个节点有三个可能股票价而只有股票和账户两个交易工具, 一般有三个终值方程但仅两个未知持仓, 不能保证所有合约都可复制, 风险中性概率也可能不唯一.

\begin{example}
股票 $S_0=100$, 两期内每期乘 $1.2$ 或 $0.8$, 账户 $R=1$. 求执行价 $100$ 的两期看涨价格及第一期持仓.
\end{example}
\begin{solution}
$q=1/2$. 终价为 $144,96,96,64$, 终值为 $44,0,0,0$. 因而 $V_0=44/4=11$. 一期上涨节点价格为 $22$, 下跌节点为 $0$. 初始持股
$\Delta_0=(22-0)/(120-80)=0.55$, 账户金额 $B_0=11-55=-44$.
若上涨到 $120$, 组合价值为 $22$, 下一步持股
$\Delta_1^u=44/(144-96)=11/12$, 账户金额 $22-(11/12)120=-88$.
若下跌到 $80$, 合约已无可能取得正终值, 该节点两种子节点支付均为零, 所以持股和账户金额均为零. 调整没有新增现金, 所有路径末端终值都正确.
\end{solution}

\originalsection{状态价格、风险中性概率与计价单位}
二叉树展示了风险中性概率如何由复制产生. 还有一个一般的有限状态说明, 可以把市场里的二元事件合约与概率公理联系起来. 固定终期 $T$, 有限个状态为 $\omega_1,\ldots,\omega_m$, 假设每个支付 $\ind_{\{\omega_i\}}$ 的合约都可买卖且被履约, 其当前价格为 $\pi_i$. 无风险支付 $1$ 的现价为 $D=\ee^{-rT}$.

同支付同价格迫使 $\sum_i\pi_i=D$, 因为各状态合约合起来总支付 $1$. 无套利迫使 $\pi_i\geq0$, 否则买该合约收到钱, 存入账户后即得非负且正的终值. 若每个状态实际有正概率且零价合约可交易, 还要求 $\pi_i>0$, 否则免费买入就形成套利. 定义
\[
 Q(\{\omega_i\})=\frac{\pi_i}{D},\qquad
 Q(A)=\sum_{\omega_i\in A}\frac{\pi_i}{D}.
\]
它确为有限状态空间上的概率. 任意支付 $H$ 为各状态合约的线性组合, 所以
\[
 \operatorname{Price}(H)=\sum_i\pi_iH(\omega_i)
 =D\E_QH.
\]
对于无限状态, 单靠无套利的有限买卖只能推出有限可加性; 要从无限事件合约推出可数可加性还需价格对单调极限的连续性等条件. 本章有限树证明没有这项缺口; 后面的连续模型直接假设有一个可数可加的风险中性概率, 不把其一般存在性误说成已证.

风险中性概率与实际频率不同, 因为不同状态的一元货币对人的价值可能不同. 经济不景气时多收到一元可能更有用, 保险支付可能因此更贵. 这只是解释价格与频率可能分离的经济直觉, 不是从无套利单独推出的严格大小关系. 买卖价差、交易限制及履约风险若存在, 同终值组合的差价也未必能按上述理想方式套利.

计价单位也影响概率. 在上述有限状态模型中, 设一个正资产终值为 $A_i>0$, 当前价 $A_0=\sum_i\pi_iA_i$. 以该资产为计价单位, 支付“一单位资产且仅在状态 $i$ 发生时”的现价相对于 $A_0$ 为
\[
 Q^A(\{\omega_i\})=\frac{\pi_iA_i}{A_0}
 =\frac{Q(\{\omega_i\})A_i}{\E_QA_T}.
\]
这些权重相加为 $1$, 且
\[
 \operatorname{Price}(H)=A_0\E_{Q^A}\left(\frac{H}{A_T}\right).
\]
所以更换计价单位确实改变概率权重. 这个有限求和推导给出精确公式, 无需调用抽象测度变换定理. 在某状态终值较低的计价资产, 该状态在新概率中的相对权重也较低.

\originalsection{看涨函数、看跌函数与分布恢复}
先暂时不贴现金融价格. 对非负且可积随机变量 $X$, 定义
\[
 \mathcal C(K)=\E(X-K)^+,\qquad
 \mathcal P(K)=\E(K-X)^+,\qquad K\geq0,
\]
其中 $a^+=\max(a,0)$, $F(x)=\Prob(X\leq x)$. 这两个函数只由概率分布决定.

\begin{proposition}[尾积分与看涨函数]
\[
 \mathcal C(K)=\int_K^\infty(1-F(x))\dd x,\quad
 \mathcal P(K)=\int_0^K F(x)\dd x,\quad
 \mathcal C(K)-\mathcal P(K)=\E X-K.
\]
$\mathcal C$ 非负、递减且凸, $\mathcal C(0)=\E X$, $\mathcal C(K)\to0$ 当 $K\to\infty$. 它的右、左导数分别为
\[
 \mathcal C'_+(K)=-\Prob(X>K),\qquad
 \mathcal C'_-(K)=-\Prob(X\geq K)\quad(K>0).
\]
若 $F$ 连续, 则 $\mathcal C'(K)=F(K)-1$. 若 $X$ 具有在 $K$ 连续的密度 $f$, 则 $\mathcal C''(K)=f(K)$.
\end{proposition}
\begin{proof}
对每个确定 $x\geq0$,
$(x-K)^+=\int_K^\infty\ind_{\{x>t\}}\dd t$,
$(K-x)^+=\int_0^K\ind_{\{x<t\}}\dd t$.
非负函数可以交换积分次序; 可由矩形上的非负简单函数逐次逼近证明. 所以第一式成立. 第二式的被积函数是 $\Prob(X<t)$, 它与 $F(t)$ 仅在原子点不同; 正概率原子至多可数, 不改变积分, 得所写公式. 也可用
$(x-K)^+-(K-x)^+=x-K$ 直接得到平价恒等式.

每个固定 $x$ 上 $(x-K)^+$ 都递减且凸, 取期望保留这两性质. 又
\[
 0\leq\mathcal C(K)\leq\E[X\ind_{\{X>K\}}]\longrightarrow0,
\]
尾部期望趋零由 $X$ 可积保证. 从右侧差商
\[
 \frac{(X-K-h)^+-(X-K)^+}{h},\qquad h>0,
\]
绝对值不超过 $1$, 点态极限为 $-\ind_{\{X>K\}}$. 对这些绝对值不超过 $1$ 的差商可交换期望和极限, 得右导数. 左侧相同论证给出 $-\ind_{\{X\geq K\}}$. 在无原子点两者相同, 得到一阶导数. $F'=f$ 在密度连续处成立, 再求导得到二阶公式.
\end{proof}
原子处两侧斜率的差为 $\Prob(X=K)$, 所以折角记录离散概率质量. 不能对任意分布写一个普通函数 $f=\mathcal C''$; 有原子时应使用斜率跳跃或分布意义的导数.

若 $X=S_T$ 是风险中性终价, 欧洲看涨和看跌期权的当前价格分别为
\[
 C_0(K)=\ee^{-rT}\mathcal C_Q(K),\qquad
 P_0(K)=\ee^{-rT}\mathcal P_Q(K).
\]
因为无分红股票满足 $\ee^{-rT}\E_QS_T=S_0$, 得到看涨看跌平价
\[
 C_0(K)-P_0(K)=S_0-K\ee^{-rT}.
\]
连续分布下
\[
 Q(S_T>K)=-\ee^{rT}C_0'(K),\qquad
 f_Q(K)=\ee^{rT}C_0''(K).
\]
第一阶导数来自很窄执行价区间的看涨价差, 第二阶导数来自相邻执行价的二阶差分. 真实报价只在有限个执行价上存在, 所以只能估计导数, 还受到噪声、买卖价差与插值影响. 这些式子恢复风险中性分布, 不直接恢复实际频率.

\begin{example}
$X$ 在 $1,3$ 两点上概率分别为 $2/3,1/3$. 写出 $\mathcal C$, 并说明如何由斜率恢复质量.
\end{example}
\begin{solution}
\[
 \mathcal C(K)=
 \begin{cases}
 5/3-K,&0\leq K\leq1,\\
 1-K/3,&1\leq K\leq3,\\
 0,&K\geq3.
 \end{cases}
\]
斜率从 $-1$ 在 $K=1$ 跳到 $-1/3$, 跳幅 $2/3$; 在 $K=3$ 再跳到 $0$, 跳幅 $1/3$. 这恰是两点概率. 普通二阶导数在其余点为零, 所以忽略折角会错失整个分布.
\end{solution}

\originalsection{二叉树的对数正态极限}
Black--Scholes 模型的核心假设是风险中性终价对数正态. 为解释其来源, 将总时长 $T$ 划为 $n$ 段, $h=T/n$, 取
\[
 u_h=\ee^{\sigma\sqrt h},\qquad d_h=\ee^{-\sigma\sqrt h},
 \qquad R_h=\ee^{rh},\qquad \sigma>0.
\]
$h$ 充分小时 $d_h<R_h<u_h$, 因为对数上的严格不等式为
$-\sigma\sqrt h<rh<\sigma\sqrt h$. 风险中性向上概率为
\[
 q_h=\frac{\ee^{rh}-\ee^{-\sigma\sqrt h}}
 {\ee^{\sigma\sqrt h}-\ee^{-\sigma\sqrt h}}.
\]
设第 $j$ 步的对数增量为 $Y_{h,j}$, 以概率 $q_h$ 取 $\sigma\sqrt h$, 否则取 $-\sigma\sqrt h$, 各步独立. 则
$\ln S_T^{(n)}=\ln S_0+\sum_{j=1}^nY_{h,j}$.

\begin{proposition}[固定终期的正态极限]
上述模型中
\[
 \ln S_T^{(n)}\ \Longrightarrow\
 \Normal\left(\ln S_0+(r-\sigma^2/2)T,\,\sigma^2T\right).
\]
因此 $S_T^{(n)}$ 分布收敛到相应对数正态变量.
\end{proposition}
\begin{proof}
令 $a=\sigma\sqrt h$. 泰勒展开的余项在 $h\downarrow0$ 时统一控制, 得
\begin{align*}
 \ee^{rh}-\ee^{-a}&=a+(r-\sigma^2/2)h+O(h^{3/2}),\\
 \ee^a-\ee^{-a}&=2a+O(h^{3/2}),
\end{align*}
所以
\[
 q_h=\frac12+\frac{r-\sigma^2/2}{2\sigma}\sqrt h+O(h),
 \quad m_h:=\E_QY_{h,1}=(r-\sigma^2/2)h+O(h^{3/2}),
\]
而 $v_h:=\Var_QY_{h,1}=\sigma^2h-m_h^2=\sigma^2h+O(h^2)$.
于是 $nm_h\to(r-\sigma^2/2)T$, $nv_h\to\sigma^2T$.
这里每行的分布随 $n$ 改变, 不能直接当作固定分布的独立同分布中心极限定理. 我们给出特征函数计算来核实所需版本.

设 $Z_h=Y_{h,1}-m_h$. 对固定实数 $t$, 使用
$\ee^{itZ_h}=1+itZ_h-t^2Z_h^2/2+O(|t|^3|Z_h|^3)$.
因为 $\E Z_h=0$ 且 $|Z_h|\leq2\sigma\sqrt h$, 有
\[
 \E\ee^{itZ_h}=1-\frac{t^2v_h}2+O(h^{3/2}).
\]
独立性使中心化和的特征函数为此式的 $n$ 次幂. 利用
$\ln(1+w)=w+O(w^2)$ (复数 $w$ 近零), 其对数为
\[
 n\ln\E\ee^{itZ_h}
 =-\frac{nt^2v_h}2+O(nh^{3/2})+O(nh^2)
 \longrightarrow-\frac{t^2\sigma^2T}2.
\]
加上均值因子便趋于所陈述正态分布的特征函数. 这里使用特征函数连续性定理: 若特征函数点态趋于某个在零连续的概率分布的特征函数, 则相应分布弱收敛. 此定理是特征函数极限判据, 本章不重新证明一般傅里叶反演理论; 上面的行列变化、误差累积与极限参数均已直接计算. 指数函数连续, 所以指数变换的分布也收敛; 等价地在每个 $x>0$ 上对 $\ln x$ 比较分布函数即可.
\end{proof}
注意 $-\sigma^2/2$ 修正不能省略. 股价的风险中性均值增长率为 $r$, 但对数的均值增长率为 $r-\sigma^2/2$, 原因是对数凹性与指数的正态均值修正.

仅有分布收敛不足以保证看涨期权的期望收敛, 因为支付 $(S-K)^+$ 无界. 在这个特定树上还可建立足够的尾部控制. 一步二阶矩因子为
\[
 q_h\ee^{2a}+(1-q_h)\ee^{-2a}
 =1+(2r+\sigma^2)h+O(h^{3/2}),
\]
于是独立性给出
\[
 \E_Q[(S_T^{(n)})^2]
 =S_0^2\bigl[1+(2r+\sigma^2)h+O(h^{3/2})\bigr]^n,
\]
对所有充分大的 $n$ 一致有界. 对非负变量 $S$, 尾部估计
$\E[S\ind_{\{S>M\}}]\leq\E S^2/M$ 说明这些股价的尾部期望一致趋零. 将支付截断为
$\min\{(S-K)^+,M\}$. 用非负尾积分公式,
\[
 \E\min\{(S-K)^+,M\}=\int_0^M\Prob(S>K+t)\dd t.
\]
极限对数正态分布的分布函数处处连续, 所以每个尾概率都收敛; 被积函数又在有限区间 $[0,M]$ 上被 $1$ 控制, 可交换极限与积分. 这直接证明截断支付的期望收敛, 无需另引一般的弱收敛等价刻画. 截断误差至多 $\E[S\ind_{\{S>K+M\}}]$, 由上述二阶估计一致控制; 极限对数正态变量也有有限二阶矩. 先令 $n\to\infty$, 再令 $M\to\infty$, 得到看涨期望也收敛. 每个树上的贴现因子 $R_h^{-n}=\ee^{-rT}$ 相同, 所以树价格趋于下面的公式.

这是一种固定终期分布及价格的极限证明, 不是整条随机路径收敛或连续时间复制策略收敛的证明. 若要严格构造连续时间的布朗运动及随机积分, 需要更进一步的概率理论. 对本章的欧洲终期支付, 固定终期极限已经够用.

\originalsection{Black--Scholes 模型与公式}\label{sec:bs}
连续时间的理想模型假定无分红股票、恒定无风险利率 $r$、恒定波动率 $\sigma>0$, 无交易摩擦且可连续交易. 在风险中性概率 $Q$ 下, 布朗运动 $W_t$ 有连续路径, $W_0=0$, 不重叠时间段的增量独立, $W_t-W_s\sim\Normal(0,t-s)$. 这里取自然信息 $\mathcal F_s=\sigma(W_u:0\le u\le s)$; 不额外预先透露未来增量. 定义
\[
 S_t=S_0\exp\left((r-\sigma^2/2)t+\sigma W_t\right).
\]
仅为了终期定价, 可以不先构造完整过程, 而直接假定
\[
 N:=\ln S_T\sim\Normal(\mu,\sigma^2T),\qquad
 S_0=\ee^{-rT}\E_QS_T.
\]
后一个等式是模型的贴现价格条件. 正态指数矩 $\E\ee^N=\ee^{\mu+\sigma^2T/2}$ 给出
\[
 \mu=\ln S_0+(r-\sigma^2/2)T.
\]
因此无需另行猜测真实世界下股票上涨趋势来填写 $\mu$. 它在风险中性模型里已由 $S_0,r,\sigma,T$ 决定.

完整过程也确实满足贴现鞅条件: 令 $A_t=\ee^{rt}$, 对 $s<t$,
\[
 \frac{S_t}{A_t}=\frac{S_s}{A_s}
 \exp\left(-\frac{\sigma^2}2(t-s)+\sigma(W_t-W_s)\right).
\]
未来增量与当前信息独立, 正态指数矩使后一个因子的均值为 $1$, 因而
$\E_Q(S_t/A_t\mid\mathcal F_s)=S_s/A_s$. 本章是在 $Q$ 下直接定义过程, 并未从任意实际概率下的股票模型证明风险中性测度存在, 也未调用 Girsanov 定理.

\begin{theorem}[欧洲看涨期权公式]
设 $S_0>0$, $K>0$, $T>0$, $\sigma>0$. 在上述模型中, 到期支付 $(S_T-K)^+$ 的当前价为
\[
 C_0=S_0\Phi(d_1)-K\ee^{-rT}\Phi(d_2),
\]
其中 $\Phi(z)=\int_{-\infty}^z(2\pi)^{-1/2}\ee^{-x^2/2}\dd x$,
\[
 d_1=\frac{\ln(S_0/K)+(r+\sigma^2/2)T}{\sigma\sqrt T},
 \qquad d_2=d_1-\sigma\sqrt T.
\]
\end{theorem}
\begin{proof}
设 $v=\sigma^2T$, $\mu=\ln S_0+(r-\sigma^2/2)T$, $k=\ln K$. 将支付拆成两个截断期望,
\[
 C_0=\ee^{-rT}\left(\E_Q[\ee^N\ind_{\{N>k\}}]
                 -KQ(N>k)\right).
\]
先算概率项. 标准化给出
\[
 Q(N>k)=1-\Phi\left(\frac{k-\mu}{\sqrt v}\right)
 =\Phi\left(\frac{\mu-k}{\sqrt v}\right)=\Phi(d_2),
\]
最后使用标准正态关于零的对称性.

指数项的困难是密度之外多出 $\ee^x$. 在指数中配方,
\[
 x-\frac{(x-\mu)^2}{2v}
 =\mu+\frac v2-\frac{(x-\mu-v)^2}{2v}.
\]
因而
\begin{align*}
 \E_Q[\ee^N\ind_{\{N>k\}}]
 &=\frac1{\sqrt{2\pi v}}\int_k^\infty
   \exp\left(x-\frac{(x-\mu)^2}{2v}\right)\dd x\\
 &=\ee^{\mu+v/2}\frac1{\sqrt{2\pi v}}\int_k^\infty
   \exp\left(-\frac{(x-\mu-v)^2}{2v}\right)\dd x\\
 &=\ee^{\mu+v/2}\Phi\left(\frac{\mu+v-k}{\sqrt v}\right)
 =S_0\ee^{rT}\Phi(d_1).
\end{align*}
将两项代回得到公式. 这也直接验证支付期望有限, 因为 $0\leq(S_T-K)^+\leq S_T$ 且 $\E_QS_T=S_0\ee^{rT}$.
\end{proof}
$d_2$ 给出的 $\Phi(d_2)$ 是风险中性行权概率 $Q(S_T>K)$.
$\Phi(d_1)$ 通常不是这个概率: 它来自额外乘以股价后的截断期望. 若按前述计价单位公式改用股票作为正计价资产, 则它确可解释为新概率下的行权概率. 这两个正态参数相差 $\sigma\sqrt T$, 正是配方造成的均值平移.

由看涨看跌平价, 欧洲看跌价格为
\[
 P_0=K\ee^{-rT}\Phi(-d_2)-S_0\Phi(-d_1).
\]
边界参数不必强行代入除以 $\sigma\sqrt T$ 的表达式: $T=0$ 时价格是 $(S_0-K)^+$; $K=0$ 时看涨支付就是股票, 价格为 $S_0$; $\sigma=0$ 时终价确定为 $S_0\ee^{rT}$, 看涨价格为 $(S_0-K\ee^{-rT})^+$.

\begin{example}
设 $S_0=K=100$, $r=0$, $T=1$, $\sigma=0.2$. 求看涨和看跌价格, 以及风险中性行权概率.
\end{example}
\begin{solution}
$d_1=0.1$, $d_2=-0.1$. 因而
\[
 C_0=100[\Phi(0.1)-\Phi(-0.1)]
 =100[2\Phi(0.1)-1]\approx7.9656.
\]
平价给出 $P_0=C_0$, 因为现价等于执行价且无利息. 风险中性行权概率为
$\Phi(-0.1)\approx0.46017$, 小于 $1/2$. 尽管 $\E_QS_T=100$, 对数均值为 $\ln100-0.02$, 中位数低于 $100$, 少量较大上涨补足均值.
\end{solution}

\originalsection{隐含波动率与模型限制}
波动率 $\sigma$ 表示单位时间对数增量的标准差, $\sigma^2$ 是方差率, 两者不能混用. 给定 $S_0,K,r,T$ 和看涨报价, 把使公式匹配报价的 $\sigma$ 称为隐含波动率. 先看公式对 $\sigma$ 是否严格递增, 才能讨论反解的唯一性.

\begin{proposition}[波动率敏感度]
令 $\phi(z)=(2\pi)^{-1/2}\ee^{-z^2/2}$. 对 $T,K,\sigma>0$,
\[
 \frac{\partial C_0}{\partial\sigma}
 =S_0\phi(d_1)\sqrt T>0.
\]
因此报价位于 $((S_0-K\ee^{-rT})^+,S_0)$ 内时, 存在唯一的正隐含波动率.
\end{proposition}
\begin{proof}
由 $d_1-d_2=\sigma\sqrt T$ 及
$d_1+d_2=2[\ln(S_0/K)+rT]/(\sigma\sqrt T)$,
\[
 \frac{\phi(d_1)}{\phi(d_2)}
 =\exp\left(-\frac{d_1^2-d_2^2}{2}\right)
 =\exp[-\ln(S_0/K)-rT],
\]
所以 $S_0\phi(d_1)=K\ee^{-rT}\phi(d_2)$.
对价格公式求导, 两项相消后剩下
\[
 \frac{\partial C_0}{\partial\sigma}
 =S_0\phi(d_1)\left(\frac{\partial d_1}{\partial\sigma}
             -\frac{\partial d_2}{\partial\sigma}\right)
 =S_0\phi(d_1)\sqrt T.
\]
当 $\sigma\downarrow0$, 公式极限为 $(S_0-K\ee^{-rT})^+$, 包括 $S_0=K\ee^{-rT}$ 时极限零. 当 $\sigma\to\infty$, $d_1\to\infty$, $d_2\to-\infty$, 所以价格趋于 $S_0$. 连续性和严格单调性给出开区间内反解存在且唯一.
\end{proof}
在同一股票、同一期限的理想恒波动率模型中, 所有执行价应给出同一个隐含波动率. 若不同执行价反推出不同数值, 它们描述的是报价相对于模型的差异, 而不是证明存在多个实际恒波动率. 波动率微笑或偏斜可以反映更厚的尾部、跳跃或波动率随状态变化等现象.

二叉树极限依赖小而独立、尺度为 $\sqrt h$ 的对数增量及恒定参数. 日内或跨日涨跌幅可能变化, 大跳跃不能按小增量泰勒余项控制, 真实利率和未来波动率也可能随机. 这些变化会改变终期分布. 本章没有证明任意真实股票都对数正态, 没有证明实际概率等于风险中性概率, 也没有用历史标准差自动替代未来模型参数. 连续时间精确复制涉及随机积分与自融资策略的可行性; 本章提供的是每棵有限二叉树的精确复制和其欧洲价格极限, 并完整计算给定对数正态风险中性模型下的公式.

\originalsection{练习与解答}
\begin{exercise}
公平随机游走从 $3$ 出发, 首次到达 $0$ 或 $10$ 时停止. 求到达 $10$ 的概率、平均停止时间, 并说明为什么不能直接对无界随机游走调用有界过程版本.
\end{exercise}
\begin{solution}
成功概率为 $3/10$, 平均停止时间为 $3(10-3)=21$. 使用的有界鞅是停止后的过程 $S_{n\wedge T}$, 它的值在 $[0,10]$ 内; 原过程 $S_n$ 本身没有统一界. 平均时长用 $S_n^2-n$ 在有界停止时间 $T\wedge N$ 上计算, 随后分别控制平方项和非负时间项的极限, 不能说 $S_n^2-n$ 是有界鞅.
\end{solution}

\begin{exercise}
股票现价 $50$, 一期后为 $65$ 或 $40$, 账户一期增长因子 $1.02$. 合约上涨支付 $8$, 下跌支付 $3$. 求风险中性概率、复制持仓和当前价格.
\end{exercise}
\begin{solution}
$u=1.3$, $d=0.8$, $q=(1.02-0.8)/(1.3-0.8)=0.44$. 持股
$\Delta=(8-3)/(65-40)=1/5$. 账户终值由下状态为
$3-(1/5)40=-5$, 所以当前账户金额 $B=-5/1.02=-250/51$.
价格 $V_0=10-250/51=260/51\approx5.09804$.
风险中性计算也为
$[0.44(8)+0.56(3)]/1.02=5.2/1.02=260/51$.
实际上涨概率可以为 $0.7$, 但不会因此改变这组复制终值或无套利价格.
\end{solution}

\begin{exercise}
非负终价 $X$ 的风险中性分布连续, 无风险贴现因子为 $D$. 证明对 $0<K_1<K_2$, 看涨价差满足
\[
 D\Prob_Q(X>K_2)\leq
 \frac{C_0(K_1)-C_0(K_2)}{K_2-K_1}
 \leq D\Prob_Q(X>K_1).
\]
说明怎样由很窄的价差估计一个尾概率.
\end{exercise}
\begin{solution}
由尾积分公式, 中间量等于
$D(K_2-K_1)^{-1}\int_{K_1}^{K_2}\Prob_Q(X>x)\dd x$.
尾概率随 $x$ 递减, 所以区间平均介于两端之间. 当两端同时趋于连续点 $K$, 两端尾概率趋于同一个值, 夹逼给出 $D\Prob_Q(X>K)$. 因此窄价差除以执行价差, 再除以 $D$, 可估计尾概率. 若有原子, 左右逼近分别给出包含和不包含该点的尾概率.
\end{solution}

\begin{exercise}
在 Black--Scholes 模型中求只在 $S_T>K$ 时支付 $1$ 的二元合约价格, 并与看涨价格的执行价导数比较.
\end{exercise}
\begin{solution}
终价对数标准化给出 $Q(S_T>K)=\Phi(d_2)$, 因而二元合约价格为 $\ee^{-rT}\Phi(d_2)$.
从看涨函数的一阶导数公式知
$-\partial C_0/\partial K=\ee^{-rT}\Phi(d_2)$.
也可以直接对 Black--Scholes 公式求导: 两个密度项因
$S_0\phi(d_1)=K\ee^{-rT}\phi(d_2)$ 及
$\partial d_1/\partial K=\partial d_2/\partial K=-1/(K\sigma\sqrt T)$ 相消, 留下 $-\ee^{-rT}\Phi(d_2)$. 这将概率计算、复制定价的含义和从价格恢复分布联系起来.
\end{solution}

\originalchapter{综合练习与逐讲对应}
\sourcelecture{15,28,37--39}
前两次复习讲义分别把第1--14讲与第17--27讲的内容重新串联. 其中的计数、公理、分布与矩计算已在前文推导, 不把同一证明按课件渐显页重复排印. 最后三次复习强调多个工具的配合: 先用指标拆计数, 再检查条件分布, 或把停止问题、极限定理和定价计算联系起来. 本章按这些训练目标设置完整练习, 参数与情境另行组织.

\originalsection{指标计数与条件分布}
\begin{exercise}
两次独立排名均在 $n!$ 个排列中均匀选择, $n\ge2$. 令 $F$ 为名次不变的对象数, $U_s$ 为恰好前进 $s$ 位的对象数, $1\le s<n$. 求二者期望与方差. 后退 $s$ 位的对象数又如何?
\end{exercise}
\begin{solution}
固定第一次排名, 第二次仍为均匀排列. 令 $I_i$ 表示原第 $i$ 名最后仍第 $i$ 名. 有 $\E I_i=1/n$, 对 $i\ne j$, $\E I_iI_j=1/[n(n-1)]$. 因而
\[
\E F=1,\qquad \E F^2=n\frac1n+n(n-1)\frac1{n(n-1)}=2,
\qquad\Var F=1.
\]
对前进数, 只有原名次 $s+1,\ldots,n$ 可以前进 $s$ 位. 记 $m=n-s$, 相应指标为 $J_i=\ind_{\{\text{新名次}=i-s\}}$. 任意两个要求的目的名次不同, 所以单项与双项的概率仍是 $1/n$ 与 $1/[n(n-1)]$. 得
\[
\E U_s=\frac mn,\qquad
\Var U_s=\frac mn+\frac{m(m-1)}{n(n-1)}-\frac{m^2}{n^2}.
\]
后退情形同理, 可行原名次为 $1,\ldots,n-s$, 因而期望与方差相同. 例如 $n=10,s=3$ 时, $\E U_3=7/10$, $\Var U_3=203/300$. 指标之间并不独立; 直接代入二项方差会丢掉排列的约束.
\end{solution}

\begin{exercise}
十二次独立抽签, 每次等概率取得五种符号之一. 已知第三种符号恰好出现三次, 求第一种符号恰好出现五次的条件概率.
\end{exercise}
\begin{solution}
先固定第三种符号所在的三个位置. 条件下其余九个位置独立且均匀地取其余四种符号, 所以第一种符号的计数为 $\Bin(9,1/4)$, 答案是
\[
\binom95\left(\frac14\right)^5\left(\frac34\right)^4.
\]
也可直接用比值核验. 交集概率是 $\binom{12}{3}\binom95 3^4/5^{12}$, 条件事件概率是 $\binom{12}{3}4^9/5^{12}$, 相除得到同一值. 不指定三个位置时仍可求这个二项分布, 因为每一种位置安排给出同样的条件计数分布.
\end{solution}

\originalsection{连续分布与独立和}
\begin{exercise}
两类报警分别是速率为2与5的独立泊松过程, 时间单位为一个月. 令 $A,B$ 是各自首次报警时间. 求 $\E A^2$, $\Cov(A,B)$, $\min(A,B)$ 的密度, 以及一个月分成六个等长区间后, 没有任何报警的区间数的期望和方差.
\end{exercise}
\begin{solution}
等待时间分别为 $\Exp(2)$ 与 $\Exp(5)$, 独立性给 $\Cov(A,B)=0$. 指数矩公式给 $\E A^2=2!/2^2=1/2$. 对 $t\ge0$,
\[
\Prob(\min(A,B)>t)=\Prob(A>t)\Prob(B>t)=\ee^{-7t},
\]
所以最小值密度为 $7\ee^{-7t}$, $t>0$, 其他处为0. 每个区间无报警的概率为 $p=\ee^{-7/6}$. 不同区间的两类增量共同独立, 故无报警区间数 $K\sim\Bin(6,p)$, 从而 $\E K=6\ee^{-7/6}$, $\Var K=6p(1-p)$. 只有期望时, 无须先证明六个指标独立; 求这个方差时独立增量才真正派上用场.
\end{solution}

\begin{exercise}
$X$ 在 $[-2,2]$ 均匀分布. 求 $\Var(X^2)$. 再令 $X_1,\ldots,X_n$ 为独立同分布样本, 求最小值 $M$ 的密度和期望.
\end{exercise}
\begin{solution}
对正整数 $k$, 对称积分给 $\E X^{2k}=2^{2k}/(2k+1)$. 因而
\[
\Var(X^2)=\frac{16}{5}-\left(\frac43\right)^2=\frac{64}{45}.
\]
当 $-2<x<2$ 时, $\Prob(M>x)=((2-x)/4)^n$, 求导得
\[
f_M(x)=\frac n4\left(\frac{2-x}{4}\right)^{n-1},
\]
支持集之外为0. 为算期望, 令 $V=(M+2)/4$, 则 $V$ 是 $n$ 个 $(0,1)$ 均匀样本的最小值. 尾积分给 $\E V=\int_0^1(1-t)^n\dd t=1/(n+1)$. 所以 $\E M=-2+4/(n+1)$. 大样本的最小值向左端点集中, 和单个样本均值0是不同问题.
\end{solution}

\begin{exercise}
设 $X_1,\ldots,X_n$ 独立同分布, $X$ 的矩母函数 $M_X(t)$ 在零的某邻域有限. 求样本平均 $\overline X_n$ 的矩母函数. 解释标准柯西样本平均为什么不会依概率收敛到一个有限常数.
\end{exercise}
\begin{solution}
指数把和变成积, 独立性再使期望分解:
\[
M_{\overline X_n}(t)=\E\prod_{i=1}^n\ee^{tX_i/n}
=\bigl(M_X(t/n)\bigr)^n.
\]
标准柯西分布的特征函数为 $\ee^{-|t|}$, 因而柯西平均的特征函数仍为 $(\ee^{-|t|/n})^n=\ee^{-|t|}$, 即它仍具有同一柯西分布. 若收敛到常数 $c$, 则对任意 $\eps>0$ 都应有 $\Prob(|\overline X_n-c|>\eps)\to0$. 但这个概率不随 $n$ 改变且为正, 因为柯西密度在整条实轴为正. 不存在这样的收敛. 大数定律的可积条件不能由独立同分布替代.
\end{solution}

\originalsection{熵、状态与时间平均}
\begin{exercise}
$X,Y$ 独立, 均以概率 $p\in(0,1)$ 取1, 以概率 $1-p$ 取0. 比较 $H(X+Y)$ 与 $H(X,Y)$, 并求差值. 证明对任意有限离散变量, $H(X+Y)\le H(X,Y)$.
\end{exercise}
\begin{solution}
记 $q=1-p$, 与第11章相同, 用以2为底的对数计熵. $X+Y$ 的概率分别为 $q^2,2pq,p^2$, 而联合分布的四个概率为 $q^2,pq,pq,p^2$. 因而
\[
H(X,Y)=2[-p\log_2 p-q\log_2 q],\qquad
H(X,Y)-H(X+Y)=2pq>0.
\]
相加只在和为1时丢失``哪个变量取1''的信息, 两种可能此时等概率, 所以差值正是这个条件下的一比特信息乘发生概率. 一般设 $Z=X+Y$. 因 $Z$ 是 $(X,Y)$ 的函数, 链式公式给
\[
H(X,Y)=H(Z)+H(X,Y\mid Z)\ge H(Z).
\]
有限状态保证所有熵有限; 若扩展到可数状态, 仍可按非负求和建立不等式, 但不应把两个无穷熵形式相减.
\end{solution}

\begin{exercise}
仓库有两件可重复使用的工具. 每天有两个操作各执行一次, 次序由独立公平硬币决定. 操作 $D$ 在架上有工具时移走一件, 无工具时不改变; 操作 $R$ 在架上至少有一件时不改变, 在架上没有时把两件都放回. 以每天开始时架上数量为状态. 求转移矩阵、稳态分布, 以及长期 $D$ 遇到空架的日比例.
\end{exercise}
\begin{solution}
状态依次为0,1,2. 逐个复合两种操作, $RD$ 把状态映为 $(1,0,1)$, $DR$ 把状态映为 $(2,2,1)$. 所以行向量约定 下
\[
P=\begin{pmatrix}0&1/2&1/2\\1/2&0&1/2\\0&1&0\end{pmatrix}.
\]
令 $\pi P=\pi$, 解得 $\pi_0=\pi_1/2$, $\pi_2=(\pi_0+\pi_1)/2$, 再用和为1得 $\pi=(2/9,4/9,1/3)$. 所有状态互达, 状态0有长度2和3的回路, 其公约数为1, 故有限状态不可约非周期收敛定理适用.

$D$ 遇空架只会在日初状态为0且 $D$ 先执行时发生; 若 $R$ 先执行就已有两件. 这个日事件的稳态概率为 $\pi_0/2=1/9$. 为得到实际时间比例, 将日初数量 $X_n$ 与当天操作次序 $C_n$ 一起作为六状态链. 从 $(x,c)$ 转到 $(g_c(x),d)$ 的概率为 $1/2$, 其中 $d$ 是次日的新次序. 原来的数量链不可约, 未来各次序都可正概率指定, 所以扩大后的链也不可约. 分布 $\widetilde\pi(x,c)=\pi_x/2$ 平稳: 对任意 $(y,d)$, 流入质量为 $\tfrac12\sum_x\pi_xP_{xy}=\pi_y/2$. 对这个六状态链的奖励函数 $f(x,c)=\ind_{\{x=0,\ c=D\text{先}\}}$ 使用第10章的时间平均定理, 得实际长期日比例几乎必然为 $\sum_{x,c}\widetilde\pi(x,c)f(x,c)=1/9$. 因而这里得到的是路径上的比例, 还不只是某个遥远日的事件概率.
\end{solution}

\originalsection{停止、极限与指数计算}
\begin{exercise}
独立公平步长 $\xi_i\in\{-1,1\}$, $S_n=\sum_{i=1}^n\xi_i$, $S_0=0$. 求首次到达 $-a$ 或 $b$ 时到达下边界的概率及平均到达时间, 其中 $a,b$ 为正整数. 另用中心极限定理近似 $\Prob(S_{4\cdot10^6}>4000)$.
\end{exercise}
\begin{solution}
停止时间 $T=\inf\{n:S_n\in\{-a,b\}\}$ 几乎必然有限且有有限期望: 在每段 $a+b$ 步内, 全向右走的概率为 $2^{-(a+b)}$, 从任一内部位置出发都会先到边界. 所以 $\Prob(T>k(a+b))\le(1-2^{-(a+b)})^k$, 给出几何尾界.

$S_n$ 与 $S_n^2-n$ 都是鞅. 先对有界时间 $T\wedge n$ 使用可选停止, 得 $\E S_{T\wedge n}=0$ 和 $\E S_{T\wedge n}^2=\E(T\wedge n)$. 停止后位置始终在 $[-a,b]$, 可用有界控制令 $n\to\infty$; 时间部分则由单调收敛和上面的有限期望成立. 设 $p=\Prob(S_T=-a)$, 则 $-ap+b(1-p)=0$, 所以 $p=b/(a+b)$. 再代入第二式,
\[
\E T=a^2\frac b{a+b}+b^2\frac a{a+b}=ab.
\]
这里不能直接说``停止时间有限, 所以可选停止''. 真正支持极限的是停止位置有界及停时的尾估计.

固定大时间时, $S_n$ 均值0、方差 $n$. 本题标准差为2000, 阈值是两倍标准差, 所以正态近似为 $1-\Phi(2)\approx0.02275$. 中心极限定理给近似极限, 不给此有限 $n$ 的自动误差保证; 若要误差上界, 还需额外的定量定理.
\end{solution}

\begin{exercise}
$X_i$ 独立且可积, $\E X_i=0$. 在自然过滤下, 判断以下过程是否必为鞅:
$A_n=\sum_{i=1}^n iX_i$,
$B_n=\sum_{i=1}^nX_i^2-n$,
$C_n=\prod_{i=1}^n(1+X_i)$,
$D_n=\prod_{i=1}^n(X_i-1)$.
\end{exercise}
\begin{solution}
$A_n$ 的新增项 $(n+1)X_{n+1}$ 条件均值为0, 故是鞅. $B_n$ 不必可积, 即使可积, 新增项条件均值也为 $\E X_{n+1}^2-1$, 不一定为0. 取所有 $X_i=0$ 就有 $B_n=-n$, 给出反例. 若额外假定各二阶矩为1, 它才是鞅.

对 $C_n$, 有限个独立可积因子的绝对乘积期望等于绝对期望乘积, 所以可积; 条件期望为 $C_n\E(1+X_{n+1})=C_n$, 故是鞅. $D_n$ 的条件期望是 $-D_n$, 一般不同于自身; 同样取 $X_i=0$ 得 $D_n=(-1)^n$. 把它改成 $(-1)^nD_n=\prod_{i=1}^n(1-X_i)$ 则成为鞅.
\end{solution}

\begin{exercise}
$Z\sim\Normal(0,1)$, $\Phi$ 为标准正态分布函数. 对实数 $c,d$ 与 $a<b$, 求 $\E\ee^{cZ+d}$ 及 $\E[\ee^{cZ+d}\ind_{\{a<Z<b\}}]$.
\end{exercise}
\begin{solution}
配方是计算的核心: $cz-z^2/2=c^2/2-(z-c)^2/2$. 因而
\[
\begin{aligned}
\E[\ee^{cZ+d}\ind_{\{a<Z<b\}}]
&=\ee^{d+c^2/2}\int_a^b\frac1{\sqrt{2\pi}}\ee^{-(z-c)^2/2}\dd z\\
&=\ee^{d+c^2/2}\bigl[\Phi(b-c)-\Phi(a-c)\bigr].
\end{aligned}
\]
令区间扩大为整条实轴, 得 $\E\ee^{cZ+d}=\ee^{d+c^2/2}$. 用在对数正态股票价格时, 指数因子与阈值一起平移; 忘掉阈值平移就会把看涨期权的第一个正态概率算错.
\end{solution}

\originalsection{逐讲来源对应}
全部来源是官方18.600 Fall 2019讲义目录中的38份文件. 下表中的页数是原始PDF页数, 包括渐显重复、标题与来源页, 不是本中文稿页数. 第16、29讲无讲义. 前两次复习重复数学内容归入相应章节, 期末复习的训练目标由本章练习完整体现. 第39讲末的课程推荐与告别页没有新增数学结论, 不作为教材正文复制.

\begin{longtable}{r p{8.0cm} r}
\toprule
讲次 & 主要内容与中文对应 & 原PDF页数\\
\midrule
\endhead
1 & 排列组合; 第1章 & 81\\
2 & 多项式系数与计数; 第1章 & 53\\
3 & 样本空间与集合; 第1章 & 47\\
4 & 概率公理; 第1章 & 54\\
5 & 等可能模型; 第1章 & 42\\
6 & 条件概率; 第2章 & 39\\
7 & Bayes公式与独立; 第2章 & 62\\
8 & 离散随机变量; 第3章 & 53\\
9 & 离散期望; 第3章 & 62\\
10 & 方差; 第3章 & 73\\
11 & 二项分布与投资组合; 第3章 & 62\\
12 & 泊松分布; 第4章 & 65\\
13 & 泊松过程; 第4章 & 59\\
14 & 更多离散分布; 第4章 & 63\\
15 & 第一次期中复习; 第13章 & 139\\
17 & 连续随机变量; 第5章 & 103\\
18 & 正态分布; 第5章 & 56\\
19 & 指数分布; 第5章 & 70\\
20 & 更多连续分布; 第5章 & 59\\
21 & 联合分布; 第6章 & 69\\
22 & 独立变量之和; 第6章 & 44\\
23 & 和的期望与次序统计量; 第7章 & 61\\
24 & 协方差与条件练习; 第7章 & 70\\
25 & 条件期望; 第8章 & 59\\
26 & 矩母函数与特征函数; 第8章 & 75\\
27 & 弱大数定律; 第9章 & 52\\
28 & 第二次期中复习; 第13章 & 184\\
30 & 中心极限定理; 第9章 & 78\\
31 & 强大数定律与Jensen; 第9章 & 64\\
32 & 马尔可夫链; 第10章 & 54\\
33 & 熵; 第11章 & 58\\
34 & 鞅与可选停止; 第12章 & 72\\
35 & 鞅与风险中性概率; 第12章 & 80\\
36 & 风险中性概率与Black--Scholes; 第12章 & 62\\
37 & 期末复习I; 第13章 & 20\\
38 & 期末复习II; 第13章 & 15\\
39 & 期末复习III; 第13章 & 40\\
补充 & 鞅、风险中性概率与Black--Scholes; 第12章 & 14\\
\midrule
合计 & 37份主讲义与1份补充 & 2413\\
\bottomrule
\end{longtable}
逐文件原始直链、资源页、下载日期、版权核查与SHA-256见源码包的source-manifest.json. 同一主源只算一个具体学科, 不把旧课号18.440另行计数. 原课与本稿的关系是数学内容重构与补充, 不把本稿新增证明归为原作者已经写出的证明.

\begin{thebibliography}{9}
\bibitem{mit} Scott Sheffield. \textit{18.600 Probability and Random Variables, Fall 2019}. MIT OpenCourseWare. \href{https://ocw.mit.edu/courses/18-600-probability-and-random-variables-fall-2019/pages/lecture-notes/}{官方讲义目录}: 37份主讲义及1份补充讲义.
\bibitem{license} MIT OpenCourseWare. \href{https://ocw.mit.edu/pages/privacy-and-terms-of-use/}{使用条款}; \href{https://creativecommons.org/licenses/by-nc-sa/4.0/}{CC BY-NC-SA 4.0}. 单独权利标记见 source-manifest.json.
\end{thebibliography}
\end{document}
